\documentclass[a4paper,11pt]{article}
\usepackage{fontspec}
\usepackage{polyglossia}
\setmainlanguage{french}
\setotherlanguage{english}
\setmainfont{Latin Modern Roman}
\setsansfont{Latin Modern Sans}
\setmonofont{Latin Modern Mono}
\usepackage{amsmath,amssymb}
\usepackage{unicode-math}
\setmathfont{Latin Modern Math}
\usepackage[margin=2.2cm]{geometry}
\usepackage{longtable,booktabs,array,calc}
\usepackage{enumitem}
\usepackage[dvipsnames]{xcolor}
\usepackage[most]{tcolorbox}
\usepackage{xurl}
\usepackage[hidelinks]{hyperref}
\usepackage{newunicodechar}
\usepackage{graphicx}
\usepackage{fancyhdr}
\setlength{\parindent}{0pt}
\setlength{\parskip}{0.45em}
\setlength{\emergencystretch}{3em}
\providecommand{\tightlist}{\setlength{\itemsep}{0pt}\setlength{\parskip}{0pt}}
\providecommand{\pandocbounded}[1]{#1}
\newunicodechar{≠}{\ensuremath{\neq}}
\pagestyle{fancy}
\fancyhf{}
\fancyhead[L]{\small Hypothèses du directeur et leurs appuis — v1}
\fancyhead[R]{\small Note de travail, classe C}
\fancyfoot[C]{\small\thepage}
\renewcommand{\headrulewidth}{0.3pt}

% Encadrés
\newtcolorbox{exemple}[1][]{colback=gray!6,colframe=gray!55,boxrule=0.5pt,arc=1pt,
  left=6pt,right=6pt,top=4pt,bottom=4pt,fonttitle=\bfseries\small,title={Exemple concret#1},breakable}
\newtcolorbox{enonce}{colback=blue!3,colframe=blue!40,boxrule=0.5pt,arc=1pt,
  left=6pt,right=6pt,top=4pt,bottom=4pt,fonttitle=\bfseries\small,title={Énoncé (mots du directeur, resserrés, orthographe rétablie)},breakable}
\newtcolorbox{avertissement}{colback=orange!6,colframe=orange!70!black,boxrule=0.6pt,arc=1pt,
  left=6pt,right=6pt,top=4pt,bottom=4pt,fonttitle=\bfseries,title={Avertissements de lecture},breakable}

% Tableau des avis des trois modèles
\newenvironment{avis}{\small\renewcommand{\arraystretch}{1.15}%
\begin{longtable}{@{}>{\raggedright\arraybackslash}p{1.35cm}>{\raggedright\arraybackslash}p{2.45cm}>{\raggedright\arraybackslash}p{7.15cm}>{\raggedright\arraybackslash}p{4.4cm}@{}}
\toprule \textbf{Modèle} & \textbf{Statut} & \textbf{Raison (résumée)} & \textbf{Références citées}\\ \midrule \endhead
\bottomrule \endfoot}{\end{longtable}}
\newcommand{\st}[1]{\textbf{#1}}
\newcommand{\cours}[1]{\par\smallskip\textbf{\textsf{#1}}\enspace}

\title{Les hypothèses du directeur sur ce qui fait une idée intéressante,\\ ce qui les soutient, et un cours accéléré\\[0.4em]\large Note de travail, version 1 — fusionnée avec les réponses d'Argon, Carbon et Neon}
\author{Session fille (Opus 5.5, effort élevé), mandat GO-2026-11-23-25}
\date{23 novembre 2026 (heures UTC relevées par \texttt{date -u})}

\begin{document}
\maketitle
\thispagestyle{fancy}
\begin{avertissement}
\begin{enumerate}[leftmargin=1.3em,itemsep=2pt]
\item \textbf{Question non neutre.} Argon, Carbon et Neon (trois modèles de Google) ont reçu la version~1 de la demande reformulée (\texttt{demande-reformulee-modele-idees-interessantes-v1.md}), jugée \emph{non neutre} (0,55) par le lecteur de neutralité du collisionneur (registre des décisions, R-5176). La demande exposait les hypothèses H1 à H10 comme celles d'un chercheur : leurs verdicts peuvent pencher vers l'accord. Leur convergence n'est donc qu'un indice faible (voir la partie~3).
\item \textbf{Ce qu'ils n'ont pas vu.} Les trois modèles n'ont lu ni les messages 4 à 6 du directeur ni la question de suivi sur l'homomorphisme de graphes typés. Pour les idées de ces messages (section~1.11 et suivantes), les tableaux donnent ce que leurs réponses contiennent \emph{qui s'y rapporte} ; le statut y est une lecture de la présente note, signalée comme telle.
\item \textbf{Références.} Chaque référence de la note a été cherchée en ligne le 23 novembre 2026 (Crossref par DOI ou par recherche bibliographique, interface d'arXiv, OpenLibrary pour les livres sans DOI, Wikisource pour Poincaré). La bibliographie finale donne, pour chacune, l'identifiant trouvé et le mode de vérification ; « à vérifier » quand rien n'a été trouvé. Vérifier qu'une référence existe ne vérifie pas ce qu'on lui fait dire : les résumés des travaux empiriques ont été lus (Crossref, Europe PMC, arXiv), les textes intégraux non, sauf Poincaré (Wikisource).
\item \textbf{Chiffres.} Les parties 1 à 3 ne donnent aucun chiffre tiré d'un travail sans sa source et sa page ; quand un modèle en cite un, il est rapporté comme sien, non vérifié. Les annexes reproduisent les réponses sans en changer le fond : leurs chiffres n'y sont \emph{pas} vérifiés.
\item \textbf{Statut.} Note de travail (classe C), non contre-lue, ni papier ni préenregistrement. Une version~2 intégrera la revue bibliographique en cours et les réponses de Fable, d'Opus et du collisionneur.
\end{enumerate}
\end{avertissement}

\paragraph{Entrées (sur \texttt{main}, sous \texttt{docs/notes/idee-interessante-v1/}), empreintes vérifiées le 23 novembre 2026 à 13:49 UTC.}
{\small
\begin{longtable}{@{}p{7.2cm}p{9.2cm}@{}}
\toprule Fichier & Empreinte SHA-256\\ \midrule
\nolinkurl{intuitions-directeur-2026-11-23-v1.md} (messages 1 à 6) & \texttt{e99c99115218c1cc928ea1cd005730db\-359e3863bb4345772df734bba3914380}\\
\nolinkurl{demande-reformulee-modele-idees-interessantes-v1.md} (H1 à H10) & \texttt{230585b902f912cd1c492225e18661f9\-ce7cfd20ddf7450abb18a8eef71fd3fa}\\
\nolinkurl{reponses-google-v1/reponse-argon-v1.md} & \texttt{8604a3e67d4b1e1acadaf55db71f74da\-6d4b14c3d81fb3a9ab1901efd0363521}\\
\nolinkurl{reponses-google-v1/reponse-carbon-v1.md} & \texttt{a9532fbb60c34d46647b81de1a4c52cd\-b6cec3521cf1728e7b16fc6e80db5ca6}\\
\nolinkurl{reponses-google-v1/reponse-neon-v1.md} & \texttt{b86ac5d1e4f18d79e8a36c16dd172f79\-ba18d084e683df18062cdcb4f5c43799}\\
\nolinkurl{question-suivi-homomorphisme-v1.md} (lue, non servie ici) & \texttt{9092bd3744928a298c5e7b427f4164ae\-b3bc319230bd8bbd889454318f882ce8}\\
\bottomrule
\end{longtable}}

\paragraph{Plan.} Partie~1 : les hypothèses, une par section, avec l'avis des trois modèles. Partie~2 : cours accéléré sur les travaux qui occupent déjà ces idées. Partie~3 : accords, désaccords, erreurs relevées et propositions retenues. Partie~4 (annexes) : les trois réponses en \LaTeX{}. Bibliographie vérifiée à la fin.

\tableofcontents
\newpage

%=====================================================================
\section{Les hypothèses du directeur}
%=====================================================================

\subsection*{Trois questions à ne pas confondre}
Les trois modèles commencent par la même mise en garde, et elle est juste : les hypothèses mêlent trois questions qui n'ont ni les mêmes données ni les mêmes tests.
\begin{enumerate}[itemsep=1pt]
\item \textbf{L'intérêt} : pourquoi une idée capte l'attention d'un public \emph{donné}, relativement à ce qu'il croit au moment~$t$ (Davis, 1971 ; Itti et Baldi, 2009). Une idée fausse peut être passionnante.
\item \textbf{La vérité et la fécondité} : pourquoi elle survit à l'épreuve du réel et engendre d'autres résultats (Kitcher, 1981 ; Lakatos, 1970).
\item \textbf{La production} : par quel mécanisme cognitif, social ou algorithmique elle apparaît (Koestler, 1964 ; Swanson, 1986 ; Simonton, 2004).
\end{enumerate}
Chaque hypothèse ci-dessous est rattachée à celle de ces questions qu'elle vise d'abord.

\subsection*{Notations communes}
Les trois réponses, indépendamment les unes des autres (mais sur la même consigne), construisent le même squelette ; on l'adopte ici pour exprimer les hypothèses (annexes A à C, § 2).
\begin{itemize}[itemsep=1pt]
\item $\mathcal{G}_t$ : un graphe de concepts \emph{daté} (construit seulement sur ce qui est publié avant $t$), à trois couches : l'\emph{usage} (cooccurrences, plongements de mots $x_u(t)$), la \emph{structure} (relations typées : cause, inhibe, conserve, est proportionnel à… ; cône de dépendances $G_u$ de rayon $k$ autour de $u$), les \emph{communautés} humaines (profil disciplinaire $p_u(t)$ des auteurs et revues).
\item $d_{\text{surf}}(u,v;t)$ : distance de surface, par exemple $\tfrac{\alpha}{2}(1-\cos(x_u,x_v)) + (1-\alpha)\,D_{\mathrm{JS}}(p_u\Vert p_v)$.
\item $\phi : G_u \rightharpoonup G_v$ : une application partielle injective avec $\phi(u)=v$ ; sa \emph{couverture} $\kappa(\phi)$ (part du réseau amont appariée), sa \emph{fidélité} (part des relations $(a,r,b)$ telles que $(\phi(a),r,\phi(b))$ existe) et ses \emph{conflits} $\chi(\phi)$ (relations envoyées sur une relation incompatible : excite contre inhibe).
\item $K(u,v;t)=\max_{\phi(u)=v}\ \kappa(\phi)\,\sigma(\phi)\,e^{-\gamma\chi(\phi)}$ : le score « serrure et clé » ; $d_{\text{struct}} = 1-K$, ou une distance de Gromov-Wasserstein entre $G_u$ et $G_v$ privés de leurs étiquettes.
\item $F(u,v;t)=\sum_m \alpha_m(t)\,h_m(u,t)\,h_m(v,t)$ : le « champ », co-exposition de $u$ et $v$ à un même instrument ou formalisme $m$ dont l'adoption croît au rythme $\alpha_m(t)$.
\item $\lambda_{uv}(t)$ : l'intensité (taux de hasard) de naissance du lien $u$–$v$ dans $[t,t+dt]$.
\end{itemize}

\subsection*{Comment lire les tableaux}
\emph{Statut} reprend le vocabulaire du mandat (soutenue, contredite, nuancée, métaphore à élaguer), suivi du mot exact du modèle quand il diffère. \emph{Raison} résume la réponse, sans les chiffres. \emph{Références} reprend celles que le modèle cite à cet endroit (vérifiées en bibliographie).
%---------------------------------------------------------------------
\subsection{H1 — Le mariage éloigné}
\begin{enonce}
« Une idée intéressante […] associe deux concepts, deux sous-idées, deux éléments qui n'ont pas de connexion logique directe, de lien évident entre eux, mais qui en réalité en ont un caché, parfois contre-intuitif. » « Sur une dimension cachée et abstraite, grande proximité, grande similitude ; sur une autre dimension plus immédiate […] grand éloignement, grande différence, et justement le signal surprise est lié à cette distance sur cette dimension-là. »
\end{enonce}

\textbf{Explication.} L'hypothèse pose \emph{deux} distances entre les parents $u$ et $v$ et leur donne deux rôles distincts. La distance de surface $d_{\text{surf}}$ (mots, disciplines, communautés) fixe la \emph{surprise} : plus elle est grande, moins le lien était attendu, au sens où sa probabilité a priori $P_{\text{surf}}(u\!-\!v)$ était faible et sa « surprise » $-\log P_{\text{surf}}$ grande. La distance de structure $d_{\text{struct}}$ (forme des relations, indépendamment des noms) fixe la \emph{viabilité} : le lien tient parce que les deux réseaux amont ont la même forme. Le cœur testable n'est donc pas « loin, c'est bien », mais une \emph{interaction} : l'effet de $d_{\text{surf}}$ sur la fécondité doit être positif \emph{seulement} quand $K$ est élevé. Carbon et Neon placent précisément ce terme $d_{\text{surf}}\times K$ au centre de leur expérience (annexes B et C, § 7). L'hypothèse vise surtout les questions 1 (intérêt) et 2 (fécondité) ; elle ne dit pas encore comment trouver de tels couples (question 3).

\begin{avis}
Argon & \st{Soutenue} & Correspond à la bisociation de Koestler et à l'appariement de structures de Gentner (les relations d'ordre supérieur sont projetées, les attributs de surface ignorés) ; en scientométrie, les articles à fort impact combinent un socle conventionnel et une queue de combinaisons atypiques. & Koestler 1964 ; Gentner 1983 ; Hofstadter et Sander 2013 ; Uzzi et al. 2013\\ \midrule
Carbon & \st{Soutenue} (questions 1, 2 et 3) & Même socle cognitif ; ajoute que les stratégies qui franchissent les frontières sont rares et risquées mais concentrent les percées. & Koestler 1964 ; Gentner 1983 ; Hofstadter et Sander 2013 ; Foster et al. 2015\\ \midrule
Neon & \st{Nuancée} (« soutenue avec nuance forte ») & Soutenue, mais \emph{incomplète sans l'ancrage conventionnel} : les sauts lointains ont une forte variance, un fort risque d'échec et une reconnaissance retardée ; la recombinaison locale produit aussi des percées. & Gentner 1983 ; Uzzi et al. 2013 ; Shi et Evans 2023 ; Foster et al. 2015 ; Wang et al. 2017 ; Kaplan et Vakili 2015\\
\end{avis}

\begin{exemple}[ : l'horizon acoustique d'Unruh]
Unruh (1981) remarque que le son, dans un fluide qui s'écoule plus vite que le son lui-même, obéit à la même équation d'onde qu'un champ près de l'horizon d'un trou noir : hydrodynamique et gravitation quantique sont très loin en surface (communautés, revues, vocabulaire), très proches en structure (même équation sur une « métrique effective »). C'est l'exemple de « mécanique des fluides pour étudier les trous noirs » du message~2 ; Carbon le cite (annexe B, H5). Il illustre aussi la limite : l'analogie reproduit la cinématique de l'horizon, non la dynamique d'Einstein ; ce que l'expérience de fluide peut apprendre de la gravitation est précisément ce que l'invariant conserve.
\end{exemple}

%---------------------------------------------------------------------
\subsection{H2 — Cascades et transferts}
\begin{enonce}
« C'est une cascade de mariages d'idées : a priori douleur et conscience n'ont pas grand-chose à voir, sauf que conscience et anesthésie sont liées et que douleur et anesthésie aussi. Et ce qui est connu dans le cas du cerveau animal est ensuite investigué comme lien dans les modèles de langage. » « C'est soit l'application d'un résultat connu dans un domaine à un nouveau domaine, […] tant que le domaine est jeune et en expansion. »
\end{enonce}

\textbf{Explication.} Deux mécanismes distincts sont réunis. (a) \emph{La fermeture transitive} : si $A\!-\!B$ et $B\!-\!C$ sont attestés dans deux littératures qui ne se citent pas, $A\!-\!C$ est un candidat ; son poids se mesure par les chemins à deux pas, par exemple $C_{uv}=\sum_w \tilde W_{uw}\tilde W_{wv}$ (annexe A, § 2.5). C'est le modèle « ABC » de Swanson. (b) \emph{Le transfert orienté} d'un domaine mûr (le cerveau, où douleur et anesthésie sont mesurées) vers un domaine jeune (les états internes d'un modèle de langage). Le premier mécanisme propose ; il ne valide pas : le terme pivot $B$ peut être polysémique, et les nœuds très connectés créent des chemins partout. Le second n'est valide que si les \emph{types} de relations sont conservés (c'est H3) et, plus exigeant, si le mécanisme causal l'est (Neon, annexe C, § 6.7 : un même graphe de corrélations n'implique pas un même modèle causal).

\begin{avis}
Argon & \st{Soutenue} & Inférence transitive entre littératures disjointes (Swanson, huile de poisson et syndrome de Raynaud) ; transfert d'un domaine mûr vers un domaine jeune documenté en sciences cognitives. & Swanson 1986 (\emph{Library Quarterly}) ; Holyoak et Thagard 1995\\ \midrule
Carbon & \st{Nuancée} (« soutenue pour la production, insuffisante seule pour la validité ») & Sans filtre structurel, le chaînage ABC explose en faux positifs par polysémie ; l'exemple anesthésie–modèle de langage est exposé à ce piège tant qu'on ne dit pas quel invariant de calcul remplace les boucles thalamo-corticales. & Swanson 1986 (\emph{Perspect. Biol. Med.}) ; Holyoak et Thagard 1995\\ \midrule
Neon & \st{Nuancée} (« soutenue pour la production, contredite comme critère suffisant de validité ») & Même argument : explosion combinatoire, biais des nœuds très connectés ; sans contraintes mécanistes, le transfert risque la fausse homologie. & Swanson 1986 ; Swanson et Smalheiser 1997\\
\end{avis}

\begin{exemple}[ : huile de poisson et syndrome de Raynaud]
Swanson (1986) relie deux littératures qui ne se citaient pas : l'une montre que les huiles de poisson diminuent la viscosité du sang et l'agrégation des plaquettes ; l'autre que le syndrome de Raynaud s'accompagne d'une viscosité élevée. D'où l'hypothèse « l'huile de poisson pourrait soulager le syndrome de Raynaud », ensuite éprouvée par un essai contrôlé en double aveugle (DiGiacomo et al., 1989 ; résumé lu, résultat non détaillé ici). Le pivot $B$ (viscosité) est une \emph{grandeur physiologique} commune, non un mot : c'est ce qui rend ce chaînage solide.
\end{exemple}

%---------------------------------------------------------------------
\subsection{H3 — Serrure et clé : invariance et réseaux amont}
\begin{enonce}
« Pour que deux théories ou idées s'emboîtent, il faut que les réseaux respectifs de concepts ayant généré ces théories ou idées, jusqu'à un certain point à définir, puissent s'emboîter aussi. Il ne suffit pas que deux points s'emboîtent […] ; l'invariance doit se préserver quelle que soit la rotation ou projection appliquée. » « S'il y a invariance, c'est parce que les parents peuvent s'accoupler […] donc théorème de Noether, symétrie quelque part ; il faudrait peut-être sortir l'artillerie lourde, lagrangien, action. »
\end{enonce}

\textbf{Explication.} L'hypothèse a deux étages qu'il faut séparer.
\begin{enumerate}[itemsep=1pt]
\item \emph{L'alignement des réseaux amont.} Une analogie viable n'est pas une correspondance $u\mapsto v$ isolée mais une application $\phi$ du cône $G_u$ vers $G_v$ qui préserve les relations : $(a,r,b)\in G_u \Rightarrow (\phi(a),r,\phi(b))\in G_v$, sans conflit. C'est exactement un \emph{homomorphisme de graphes typés} (partiel), et c'est le « principe de systématicité » de Gentner (1983) : on préfère les correspondances qui emportent un système connexe de relations, gouverné par des relations d'ordre supérieur (cause, implique). Mesurable par $\kappa$, $\sigma$, $\chi$ et $K$.
\item \emph{Symétrie, Noether, lagrangien.} Le théorème de Noether relie une symétrie continue d'une \emph{action} $S=\int L(q,\dot q,t)\,dt$ à une quantité conservée. Un graphe de concepts n'a ni action ni dynamique canonique ; les trois modèles écartent donc Noether. La note nuance leur argument : l'obstacle n'est pas le caractère \emph{discret} (il existe des théorèmes de Noether discrets, pour la mécanique variationnelle discrète, Marsden et West 2001, et pour les processus de Markov, Baez et Fong 2013), mais l'absence d'une dynamique (action ou générateur) dont la symétrie serait une propriété. Tant qu'aucune dynamique des idées n'est posée et mesurée, « Noether » reste une métaphore ; ce qui en survit est l'invariance par $\phi$, et, au sens des catégories, la \emph{naturalité} (le diagramme commute quelle que soit l'étape où l'on transporte).
\end{enumerate}

\begin{avis}
Argon & \st{Soutenue} (alignement) ; \st{métaphore à élaguer} (Noether, lagrangien) & Systématicité de Gentner, moteur SME, ologs de Spivak et Kent ; Noether exige des symétries continues d'une action différentiable, absentes d'un réseau discret orienté. Équivalent proposé : homomorphisme de graphes typés ou équivariance d'opérateurs discrets. & Gentner 1983 ; Falkenhainer et al. 1989 ; Spivak et Kent 2012\\ \midrule
Carbon & \st{Soutenue} (alignement) ; \st{métaphore à élaguer} (Noether) & Systématicité formalisable par des foncteurs qui préservent la structure ; Noether demande un groupe de Lie agissant sur une intégrale d'action. & Gentner 1983 ; Mac Lane 1971 ; Spivak 2014\\ \midrule
Neon & \st{Soutenue} (alignement) ; \st{métaphore à élaguer} (Noether) & Systématicité, SME, cohérence explicative de Thagard ; sans principe variationnel mesurable dans un graphe discret et dissipatif, le lagrangien est métaphorique. & Gentner 1983 ; Falkenhainer et al. 1989 ; Thagard 1989 ; Noether 1918\\
\end{avis}

\begin{exemple}[ : la Lune et la pomme, le « test de la Lune »]
L'application $\phi$ : pomme $\mapsto$ Lune, Terre $\mapsto$ Terre, « tombe vers » $\mapsto$ « tombe vers » (l'orbite est une chute continue). L'invariant conservé est la relation « l'accélération vers le centre décroît comme l'inverse du carré de la distance ». Newton ne se contente pas de la correspondance des deux points : il vérifie que l'accélération centripète de la Lune, calculée à partir de son orbite, est celle de la pesanteur à la surface réduite dans le rapport du carré des distances. C'est un contrôle de serrure et clé sur tout le réseau amont (distance, période, accélération), non sur les seuls mots « chute » et « orbite ».
\end{exemple}

%---------------------------------------------------------------------
\subsection{H4 — Ni trop près ni trop loin}
\begin{enonce}
« Qu'est-ce qui prédit la naissance d'une fille viable ? Deux parents ni trop proches, car trop proches = bassin large, ni trop lointains, car singularité, trou trop étroit, autrement dit rien. »
\end{enonce}

\textbf{Explication.} Formellement, l'intensité de naissance est non monotone en $d_{\text{surf}}$ : par exemple $\lambda\propto d^{\beta_1}e^{-\beta_2 d}$, maximale en $d^*=\beta_1/\beta_2$ (annexe A), ou une cloche $\exp(-(d-d^*)^2/2\sigma^2)$ (annexe C). Deux précisions. D'abord, la courbe n'est pas la même pour l'intérêt et pour la fécondité : la curiosité suit une courbe en U renversé (Berlyne), mais un saut local peut être très fécond (Kaplan et Vakili, 2015). Ensuite, le résultat d'Uzzi et al. (2013) n'est pas « une distance moyenne optimale » : c'est un \emph{mélange} — un socle très conventionnel plus une petite queue de combinaisons très atypiques. L'hypothèse devient donc multidimensionnelle (distribution des distances dans un même travail) plutôt que scalaire. Test proposé par les trois : une spline sur $d_{\text{surf}}$ dans un modèle additif ; réfutation si la relation est monotone.

\begin{avis}
Argon & \st{Soutenue} & Courbe en U renversé de Berlyne ; « distance cognitive optimale » entre partenaires (trop faible, rien de neuf ; trop grande, au-delà de la capacité d'absorption) ; arbitrage entre consolidation locale et sauts distants. & Berlyne 1960 ; Nooteboom et al. 2007 ; Foster et al. 2015\\ \midrule
Carbon & \st{Nuancée} (« soutenue avec nuance asymétrique ») & Le U renversé vaut pour la curiosité ; pour la fécondité, Uzzi et al. trouvent un socle conventionnel plus une queue atypique, non une distance intermédiaire uniforme. & Berlyne 1960 ; Uzzi et al. 2013\\ \midrule
Neon & \st{Nuancée} (« soutenue pour l'intérêt ; partiellement pour la fécondité ») & Le « trop près » est moins pénalisé en vérité qu'en intérêt perçu : la recombinaison locale reste productive. & Berlyne 1970 ; Loewenstein 1994 ; Davis 1971 ; Uzzi et al. 2013 ; Kaplan et Vakili 2015\\
\end{avis}

\begin{exemple}[ : l'article « conventionnel avec une pointe »]
Un article de biologie cellulaire dont presque toutes les références viennent de revues habituellement citées ensemble, plus deux références qui associent la biologie à une revue de physique statistique rarement co-citée avec elle. Selon Uzzi et al., c'est ce profil (forte conventionnalité médiane, forte nouveauté en queue), plutôt qu'un article uniformément « moyennement original », qui est associé aux plus fortes chances d'impact.
\end{exemple}

%---------------------------------------------------------------------
\subsection{H5 — Champ d'interaction et air du temps}
\begin{enonce}
« S'ils sont très éloignés, ils interagissent potentiellement beaucoup moins, SAUF SI ils répondent tous les deux au champ d'interaction, comme un matériau ferromagnétique. » « Voilà le champ d'interaction : il vient déformer la structure d'un réseau complexe en le projetant dans un nouvel espace de pensée qui, par invariance, préserve les connexions passées, jusqu'à ce que, dans cette nouvelle perspective, deux concepts éloignés finissent par se rejoindre. » « Les grandes découvertes sont faites à plusieurs endroits quasi en même temps […] ensuite c'est mécanique et aléatoire, comme avec les protéines. »
\end{enonce}

\textbf{Explication.} Trois affirmations. (a) \emph{Un facteur commun} : un cadre nouveau (instrument, technologie, formalisme) touche en même temps des concepts lointains ; mesurable comme co-exposition $F(u,v;t)=\sum_m\alpha_m(t)h_m(u,t)h_m(v,t)$, où $h_m(u,t)$ est la part des travaux sur $u$ qui emploient $m$ et $\alpha_m(t)$ le taux de croissance de $m$. (b) \emph{Les découvertes multiples} : si plusieurs équipes indépendantes forment le même lien avant de se lire, le nombre $N_{uv}$ d'équipes suit une loi de Poisson de moyenne $\Lambda_{uv}=\int\lambda_{uv}$, et la probabilité d'une découverte multiple vaut $1-(1+\Lambda)e^{-\Lambda}$ (annexes B et C) : un champ qui monte doit précéder les multiples. (c) \emph{L'emboîtement aléatoire} : une fois les réseaux rapprochés, qui conclut le premier relève du hasard — c'est la thèse « chance » de Simonton contre la thèse « air du temps » au sens fort. Le ferromagnétisme (aimantation spontanée, symétrie de spin) n'a pas d'équivalent mesurable ; seul survit le terme bilinéaire de co-exposition, qui ressemble à un champ extérieur dans un modèle d'Ising sans en avoir la physique.

\begin{avis}
Argon & \st{Soutenue} (découvertes multiples) ; « métaphore partielle » (ferromagnétisme, protéines) & Multiples recensés et théorisés (Ogburn et Thomas, Merton) ; Galison pour le substrat technique commun d'Einstein et de Poincaré ; appariement après agitation rapproché de la variation aveugle et sélection de Simonton. \emph{Mais} sa table du § 2.8 \emph{retient} le ferromagnétisme (champ d'Ising) et l'appariement des protéines. & Ogburn et Thomas 1922 ; Merton 1961 ; Galison 2003 ; Simonton 2004\\ \midrule
Carbon & \st{Soutenue} (histoire, sociologie) ; \st{métaphore à élaguer} sauf réduction & Ajoute l'analogie hydrodynamique d'Unruh et la variation aveugle de Campbell ; ferromagnétisme et protéines n'ont de valeur qu'une fois réduits à un processus ponctuel à facteur latent partagé (ce qu'il fait au § 2.5 : champ et polarisation « conservés sous forme statistique »). & Ogburn et Thomas 1922 ; Merton 1961 ; Galison 2003 ; Unruh 1981 ; Campbell 1960 ; Simonton 2004\\ \midrule
Neon & \st{Soutenue} (contexte matériel) ; \st{métaphore à élaguer} (ferromagnétisme, protéines) & La distribution des multiples s'ajuste à un processus de Poisson (Simonton 1979) ; ferromagnétisme et protéines n'ajoutent aucune prédiction propre, écartés. & Ogburn et Thomas 1922 ; Merton 1961 ; Galison 2003 ; Simonton 1979\\
\end{avis}

\begin{exemple}[ : les horloges de Berne et de Paris]
Selon Galison (2003), au tournant du siècle, la synchronisation des horloges par signaux électriques est un problème technique pressant (chemins de fer, télégraphe, détermination des longitudes). Einstein examine des brevets à Berne ; Poincaré travaille au Bureau des longitudes. Tous deux arrivent à une définition \emph{opératoire} de la simultanéité par échange de signaux. Dans les notations : l'instrument $m$ = « synchronisation électrique des horloges » ; $h_m$ élevé à la fois pour « simultanéité » (électrodynamique des corps en mouvement) et pour « mesure du temps à distance » (géodésie). C'est une thèse d'historien sur un cas, non une démonstration ; le modèle en fait une prédiction testable seulement sur un catalogue de multiples (section~\ref{sec:galison}).
\end{exemple}
%---------------------------------------------------------------------
\subsection{H6 — Connexion et distinction}
\begin{enonce}
« De nouvelles idées peuvent émerger parce que des liens qui semblent évidents entre deux choses, une invariance, une identité, ne le sont en réalité pas […] ; en tirant les ficelles par projection dans d'autres dimensions, on découvre soit des structures cachées, soit une distinction entre ce qui semblait être deux choses identiques. » « Il y a dans l'émergence de nouvelles idées fructueuses autant de travail de connexion que de distinction. » « Ce qui résonne ensemble est gardé, ce qui se cogne […] sans possibilité d'emboîtement est éliminé comme fausse invariance ou fausse identité. »
\end{enonce}

\textbf{Explication.} C'est le symétrique de H1 : deux concepts \emph{proches} en surface ($d_{\text{surf}}\approx 0$, souvent le même mot) se révèlent \emph{incompatibles} en structure quand on les confronte à un troisième concept « sonde » $v$. Deux mesures proposées : l'écart $|K(u,v)-K(u',v)|$ entre les deux candidats à l'identité (Carbon) ; le score de scission $S_{\text{split}}=(1-d_{\text{surf}})\,\mathrm{Conf}(\phi^*)$ (Neon). Le mécanisme est celui des « différences alignables » de Gentner et Markman : aligner deux structures fait apparaître exactement les points où elles divergent. Neon remarque à juste titre que les modèles quantitatifs de la science mesurent presque toujours la \emph{recombinaison} et presque jamais la \emph{scission} : H6 est la plus neuve des dix au plan de la mesure.

\begin{avis}
Argon & \st{Soutenue} & L'alignement fait émerger des différences alignables ; chez Lakatos, la collision d'une conjecture et d'un contre-exemple scinde des concepts confondus. & Gentner et Markman 1997 ; Lakatos 1976\\ \midrule
Carbon & \st{Soutenue} & Davis : une thèse intéressante unifie ce qu'on croyait hétérogène \emph{ou} dissocie ce qu'on croyait identique ; différences alignables. & Davis 1971 ; Gentner et Markman 1997\\ \midrule
Neon & \st{Nuancée} (« soutenue historiquement et cognitivement ; inconnue à grande échelle ») & Analogies négatives de Hesse ; chaleur et température séparées (Wiser et Carey) ; mais la scission n'est presque jamais testée en bibliométrie. & Hesse 1966 ; Wiser et Carey 1983\\
\end{avis}

\begin{exemple}[ : chaleur et température ; masse inerte et masse grave]
Pendant longtemps, « chaleur » et « température » désignent presque la même chose (Wiser et Carey, 1983, à partir des expérimentateurs florentins du \textsc{xvii}\ieme{} siècle ; détail à vérifier dans le chapitre, non lu). La collision avec les mélanges et les changements d'état (de la glace qui fond absorbe de la chaleur sans changer de température) oblige à les distinguer : quantité contre intensité. L'exemple du directeur, les deux masses, est plus subtil et montre les deux mouvements à la fois : Newton \emph{distingue} conceptuellement masse inerte (résistance à l'accélération) et masse grave (source et cible de la pesanteur), l'expérience les trouve égales, et le principe d'équivalence d'Einstein \emph{explique} cette égalité en en faisant une identité. Distinction d'abord, connexion ensuite.
\end{exemple}

%---------------------------------------------------------------------
\subsection{H7 — Profondeur de l'idée fille}
\begin{enonce}
« L'expérience menée en mariant deux parents doit donner une graine qui, elle, est immédiatement proche de concepts importants, des concepts “nœuds” dans notre graphe de connaissances. » « Une nouvelle idée doit expliquer a posteriori pourquoi son apparition était si surprenante… elle lève le mystère sur ce qui semblait étranger et inapplicable. »
\end{enonce}

\textbf{Explication.} Deux critères à ne pas confondre. (a) \emph{L'unification} : l'idée fille permet de dériver beaucoup de choses avec peu de schémas (Kitcher, 1981). Version mesurable \emph{avant} l'épreuve : le gain de longueur de description, $\Delta\mathcal L = \mathcal L(G_u\cup G_v) - [\mathcal L(G_{\text{abs}}) + \mathcal L(G_u\mid G_{\text{abs}})+\mathcal L(G_v\mid G_{\text{abs}})]$, où $G_{\text{abs}}$ est le schéma commun (annexe B, § 2.4) ; c'est une formalisation proposée par Carbon et Neon, non par Kitcher. (b) \emph{La position} : l'idée tombe près des nœuds centraux. Ici Neon apporte un contre-indice sérieux : les travaux les plus \emph{disruptifs} (indice $CD$ de Funk et Owen-Smith) éclipsent les anciens nœuds au lieu de s'y arrimer. Le critère (b) mesure donc plutôt la consolidation, et (a) la profondeur. Enfin, « expliquer après coup pourquoi c'était surprenant » est exposé au biais rétrospectif (Fischhoff, 1975) : à mesurer sur des jugements recueillis \emph{avant} l'issue.

\begin{avis}
Argon & \st{Soutenue} & Unification de Kitcher ; la capacité d'un nœud neuf à court-circuiter les chemins et à éclipser ses parents est mesurée par l'indice de disruption. & Kitcher 1981 ; Funk et Owen-Smith 2017 ; Wu et al. 2019\\ \midrule
Carbon & \st{Soutenue} & Profondeur = unification (synthèse newtonienne) ; l'idée qui supplante ses antécédents se mesure par la disruption. & Kitcher 1981 ; Funk et Owen-Smith 2017 ; Wu et al. 2019\\ \midrule
Neon & \st{Nuancée} (« soutenue pour l'unification ; partiellement contredite pour la proximité aux nœuds ») & Unification = compression (longueur de description minimale) ; mais les idées disruptives éclipsent les nœuds ; et l'impression de comprendre après coup est contaminée par le biais rétrospectif. & Kitcher 1981 ; Grünwald 2007 ; Funk et Owen-Smith 2017 ; Wu et al. 2019 ; Fischhoff 1975\\
\end{avis}

\begin{exemple}[ : Maxwell et la lumière]
En réunissant électricité et magnétisme dans un même système d'équations, Maxwell obtient des ondes dont la vitesse se calcule à partir de deux constantes mesurées sur des circuits et des aimants ; elle coïncide avec la vitesse de la lumière. Une seule structure explique trois domaines (électricité, magnétisme, optique) : c'est l'unification au sens (a). Et la surprise s'explique après coup : la lumière « semblait étrangère » aux bobines parce que rien, dans la surface des deux domaines, ne les rapprochait.
\end{exemple}

%---------------------------------------------------------------------
\subsection{H8 — Émergence, percolation, surface de contact}
\begin{enonce}
« Peut-être modéliser l'émergence d'une idée comme une percolation ? Il faut suffisamment de connexions entre chaînes de concepts pour relier parent A à parent B, sauf que notre percolation doit se faire entre deux milieux différents. » « C'est clairement un phénomène émergent et une transition de phase. » « Pourquoi collisionner, échanger des idées est-il un moyen fertile ? […] on augmente la surface d'attaque, de contact, donc la surface de chance de rencontre heureuse. »
\end{enonce}

\textbf{Explication.} La percolation a ici un sens exact. Sur le graphe biparti des paires viables entre deux domaines $D_1$ et $D_2$, chaque arête s'active avec une probabilité $p_{uv}$ ; une composante géante hybride (une sous-discipline) apparaît quand le rayon spectral d'un opérateur de non-retour pondéré dépasse 1 (Karrer, Newman et Zdeborová, 2014 ; annexe B, § 2.4), ou, dans une version simplifiée, quand $R=p\,(\langle k^2\rangle-\langle k\rangle)/\langle k\rangle>1$ (critère de Molloy et Reed ; annexe C, § 2.7). La « surface de contact » est mesurable : nombre de chercheurs ou de chemins à deux pas entre les communautés (Burt, 2004, pour l'avantage des courtiers). Ce qui ne survit pas : la transition de phase \emph{thermodynamique} au sens strict (exposants critiques, limite de taille infinie), que rien ne permet d'observer sur des réseaux finis et hétérogènes.

\begin{avis}
Argon & \st{Soutenue} & Expansion dans le possible adjacent (urnes de Pólya avec déclenchement) ; courtiers aux trous structuraux. & Kauffman 1995 ; Tria et al. 2014 ; Burt 2004\\ \midrule
Carbon & \st{Nuancée} (« partiellement soutenue ») & Avantage des courtiers établi ; la naissance d'une discipline présente une densification assimilable à une percolation dans le réseau des coauteurs ; la percolation \emph{conceptuelle} pure reste ouverte. & Burt 2004 ; Bettencourt et al. 2008\\ \midrule
Neon & \st{Soutenue} (contact, seuils) ; \st{métaphore à élaguer} (transition de phase stricte) & Dynamiques épidémiques à seuil pour l'émergence des champs ; pas d'exposants critiques vérifiés sur des réseaux finis. & Burt 2004 ; Bettencourt et al. 2008\\
\end{avis}

\begin{exemple}[ : un champ hybride]
Prenons deux communautés, physiciens des systèmes désordonnés et informaticiens de l'apprentissage. Tant que peu de chercheurs lisent les deux, chaque pont (verre de spin $\leftrightarrow$ réseau de neurones de Hopfield, par exemple) reste isolé. Le modèle prédit qu'au-delà d'une densité de ponts viables et de passeurs, les ponts s'enchaînent (chaque alignement réussi expose d'autres paires candidates) et une composante hybride apparaît brusquement. Exemple illustratif, non mesuré ici.
\end{exemple}

%---------------------------------------------------------------------
\subsection{H9 — La structure cachée du langage}
\begin{enonce}
« Le langage possède une structure mathématique cachée qui est le reflet des lois mathématiques qui gouvernent tout l'univers, le réel et la nature. » « Peut-être lancer une analyse en composantes principales ou un modèle entraîné non supervisé pour identifier les principales dimensions qui séparent les champs disjoints éloignés. »
\end{enonce}

\textbf{Explication.} Deux affirmations de force très différente. (a) \emph{Faible, établie} : les statistiques de cooccurrence d'un corpus, comprimées en vecteurs (plongements), portent des régularités relationnelles (« roi $-$ homme $+$ femme $\approx$ reine », Mikolov et al. 2013) et des connaissances latentes d'un domaine (Tshitoyan et al. 2019). (b) \emph{Forte, non établie} : cette structure \emph{reflète les lois de l'univers}. Ce que montrent les travaux, c'est qu'elle reflète \emph{ce que les scientifiques ont écrit} — observations accumulées, mais aussi habitudes, communautés et angles morts (Sourati et Evans, 2023). Surtout, un plongement de cooccurrence mesure d'abord la proximité de surface, exactement la distance que H1 demande de dépasser : pour la structure relationnelle, il faut des représentations explicites des relations (Lu, Wu et Holyoak 2019) ou des équations. Une analyse en composantes principales des plongements donnera les axes de \emph{surface} ; les axes « cachés » de H1 exigent la couche de structure.

\begin{avis}
Argon & \st{Soutenue} & Relations sémantiques en translations linéaires ; plongements non supervisés qui retrouvent le tableau périodique et « prédisent » des thermoélectriques avant leur découverte (Argon écrit « ont prouvé »). & Mikolov et al. 2013 ; Tshitoyan et al. 2019\\ \midrule
Carbon & \st{Nuancée} (« soutenue avec limite critique ») & Les vecteurs de cooccurrence confondent proximité thématique et isomorphisme relationnel ; il faut représenter les relations explicitement. & Mikolov et al. 2013 ; Tshitoyan et al. 2019 ; Lu et al. 2019\\ \midrule
Neon & \st{Nuancée} (« partiellement soutenue, partiellement contredite ») & Prédiction de composés confirmée, mais les vecteurs échouent sur les analogies relationnelles complexes : ils codent d'abord la surface. & Tshitoyan et al. 2019 ; Rogers et al. 2017 ; Mitchell 2021\\
\end{avis}

\begin{exemple}[ : les thermoélectriques de Tshitoyan et al.]
Des plongements de mots appris sur des résumés de science des matériaux, sans aucune connaissance chimique insérée, retrouvent la structure du tableau périodique et des relations structure–propriété ; une méthode non supervisée y recommande des matériaux pour une fonction (thermoélectricité) plusieurs années avant leur découverte (résumé de l'article). Lecture prudente : les auteurs concluent que la connaissance latente des découvertes futures est « dans une large mesure » déjà contenue dans les publications passées — c'est la « connaissance publique non découverte » de Swanson, à grande échelle, plutôt qu'un miroir direct du réel.
\end{exemple}

%---------------------------------------------------------------------
\subsection{H10 — Un signal de surprise mesurable}
\begin{enonce}
« Peut-être que quand notre organisme explore des idées, on pourrait monitorer avec une sonde le signal dopamine et repérer les moments où l'agent est surpris ? »
\end{enonce}

\textbf{Explication.} Trois grandeurs différentes se cachent sous « surprise ». (1) L'\emph{erreur de prédiction de récompense}, signée, que codent les neurones dopaminergiques (Schultz, Dayan et Montague, 1997) : elle dépend d'une récompense, pas seulement de l'improbable. (2) La \emph{surprise bayésienne}, $D_{\mathrm{KL}}(P(\theta\mid\text{donnée})\Vert P(\theta))$, ce dont l'observateur a dû réviser ses croyances (Itti et Baldi) — distincte de la simple improbabilité $-\log p$. (3) Le \emph{progrès de compression} ou d'apprentissage (Schmidhuber ; Oudeyer et al.) : l'amélioration de la capacité à prédire. La distinction compte : maximiser l'improbabilité attire vers le bruit (une neige d'écran est toujours improbable), alors que le bruit n'apporte ni révision durable ni progrès. Pour un agent de langage, des équivalents mesurables existent : la surprise d'un juge fixe sur la sortie, la divergence entre ses croyances avant et après sur un jeu de questions tenu à l'écart, ou une sonde sur les activations. Dans le programme, une sonde de production ne moyenne pas sur les jetons ; ce signal serait d'abord un instrument d'exploration.

\begin{avis}
Argon & \st{Soutenue} & Dopamine = erreur de prédiction de récompense ; curiosité = dérivée de la compression ; surprise bayésienne = divergence entre a priori et a posteriori. & Schultz et al. 1997 ; Schmidhuber 2010 ; Itti et Baldi 2009\\ \midrule
Carbon & \st{Nuancée} (« soutenue sous réserve d'une distinction ») & La dopamine code une erreur \emph{signée} de récompense, non une surprise informationnelle ; pour l'agent, viser la surprise bayésienne ou le gain de compression. & Schultz et al. 1997 ; Itti et Baldi 2009 ; Schmidhuber 2010\\ \midrule
Neon & \st{Nuancée} (« soutenue pour l'information intrinsèque ; contredite pour la surprise brute ») & Maximiser la surprise brute échoue sur le bruit ; il faut le progrès de compression ou d'apprentissage. & Schultz et al. 1997 ; Itti et Baldi 2009 ; Oudeyer et al. 2007 ; Schmidhuber 2010\\
\end{avis}

\begin{exemple}[ : le jus et la lumière]
Chez le singe, un neurone dopaminergique répond d'abord au jus inattendu ; après apprentissage, il répond à la lumière qui annonce le jus et plus au jus lui-même ; si le jus annoncé ne vient pas, son activité baisse au moment attendu (Schultz et al., 1997). Ce n'est pas « la surprise » en général : c'est l'écart entre récompense obtenue et récompense prédite. Pour l'organisme d'idéation, l'analogue exact suppose donc une récompense définie (par exemple le verdict d'un juge fixe), sans quoi la mesure pertinente est la surprise bayésienne ou le progrès.
\end{exemple}
%---------------------------------------------------------------------
\subsection*{Idées des messages 4 à 6 (non vues par les trois modèles)}
Pour les cinq sections qui suivent, les trois modèles n'ont pas lu les messages du directeur. Le tableau donne ce que leurs réponses contiennent \emph{qui s'y rapporte} ; le statut est une \textbf{lecture de la note}.

\subsection{Le champ et la distance (message 4, précisé au message 6)}\label{sec:champ-distance}
\begin{enonce}
« Au fur et à mesure que deux réseaux, ou deux bouts d'un réseau déformé, se rapprochent, le champ d'interaction doit devenir plus puissant, donc augmenter avec l'inverse d'une fonction croissante de la distance. Toute la difficulté, c'est de définir la bonne distance, dans la bonne dimension. Mais il ne suffit pas que deux réseaux se rapprochent pour qu'ils interagissent, il faut que des collisions soient activement menées, des expériences, des publications. » « Le champ augmenterait au fur et à mesure qu'on se rapproche de la bonne invariance ? » Puis, au message 6, accord du directeur sur la reformulation : le champ est \emph{la direction dans laquelle un nouveau cadre rend deux réseaux comparables} (même espace de description), et non celle où l'invariance apparaît évidente.
\end{enonce}

\textbf{Explication.} Trois idées. (a) \emph{Décroissance avec la distance} : $F\propto 1/g(d)$ avec $g$ croissante. Si $d$ est la distance de \emph{surface}, c'est l'inverse de H1 (les lointains interagiraient moins), ce que le directeur dit lui-même au message 2 ; si $d$ est la distance \emph{de structure dans la représentation courante}, alors c'est une tautologie utile : le champ mesure la vitesse à laquelle une représentation neuve fait baisser $d_{\text{struct}}$. La reformulation acceptée au message 6 tranche pour cette seconde lecture. On peut l'écrire : soit $\mathcal R_t$ l'ensemble des cadres de description disponibles à $t$ (instruments, formalismes) ; le meilleur alignement possible est $K^\star_t(u,v)=\max_{\rho\in\mathcal R_t}K(\rho G_u,\rho G_v)$ ; le champ est son accroissement, $F^\star_t(u,v)=K^\star_t-K^\star_{t-1}$. (b) \emph{Les collisions doivent être menées} : le rapprochement ne suffit pas ; il faut un terme de contact (chercheurs passeurs, publications communes), que Carbon et Neon ont mis dans $\lambda$ (ci-dessous). (c) \emph{Poincaré} « avait trouvé la bonne invariance » : il avait établi que les transformations de Lorentz forment un groupe ; pourquoi il n'a pas tiré les conséquences qu'en tire Einstein est un débat d'historiens que la note ne tranche pas.

\begin{avis}
Argon & \st{Nuancée} (lecture de la note) & Son champ $H_t$ (annexe A, § 2.4) est une co-exposition à un catalyseur, sans terme de distance ; la distance de surface entre dans $\lambda$ par une fenêtre en U renversé, pas par le champ. & Kleinberg 2003 (détection des poussées)\\ \midrule
Carbon & \st{Soutenue} pour (b) ; (a) non traité & $\lambda=(\eta_0 C_{\text{soc}}+\eta_1F)\cdot\sigma(\beta_1K-\beta_2(d_{\text{surf}}-\mu)^2-\beta_0)$ : le contact par passeurs $C_{\text{soc}}$ est un facteur séparé, additif au champ. & Burt 2004 (implicite) ; Galison 2003\\ \midrule
Neon & \st{Soutenue} pour (b) ; (a) non traité & $\lambda=\lambda_0B^{\rho}e^{\eta\Phi}\,\mathrm{cloche}(d_{\text{surf}})\,K^{\theta}$ : la surface de contact $B_{ij}$ est un facteur multiplicatif. & Burt 2004\\
\end{avis}

\begin{exemple}[ : la diffusion de la chaleur et Black et Scholes]
L'équation de la chaleur et l'équation de Black et Scholes (prix d'une option) sont, après changement de variables, la même équation parabolique (Argon, annexe A, § 7.1). Avant que le calcul stochastique ne fournisse le cadre commun (marche aléatoire, mouvement brownien), finance et physique de la chaleur n'avaient pas de « même espace de description » ; après, la distance de structure tombe d'un coup. $F^\star$ mesure ce saut ; il est mesurable avant que les deux communautés ne se citent.
\end{exemple}

%---------------------------------------------------------------------
\subsection{Le champ, les invariants et les probabilités de rencontre (message 5)}\label{sec:rencontres}
\begin{enonce}
« Le champ d'interaction doit tenir compte des invariants autant que des possibilités, potentialités, probabilités de rencontre, donc de la densité du réseau, et de comment une transformation par invariance rapproche deux réseaux jusqu'à les collisionner. »
\end{enonce}

\textbf{Explication.} Le message 5 décrit une intensité \emph{produit} de deux facteurs : la fréquence des rencontres (densité, contact) et la probabilité qu'une rencontre « prenne » (compatibilité des invariants). C'est la forme de la \emph{théorie des collisions} en cinétique chimique (vitesse = fréquence des chocs $\times$ fraction des chocs bien orientés et assez énergétiques) : l'image des protéines du message 2 trouve là son équivalent exact. Écrit dans les notations :
\[
\lambda_{uv}(t)=\underbrace{\big(\eta_0\,C_{\text{contact}}(u,v;t)+\eta_1\,F(u,v;t)\big)}_{\text{rencontres}}\times\underbrace{\Pr\big(\text{prise}\mid K(u,v;t),\,d_{\text{surf}}\big)}_{\text{emboîtement}}.
\]
Carbon et Neon ont écrit, sans avoir lu le message 5, exactement cette forme (annexe B, § 2.3 ; annexe C, § 2.5). C'est un accord sur la \emph{forme} ; il ne dit rien de la valeur des coefficients. Test : les deux facteurs doivent avoir des effets séparables (un contact élevé sans compatibilité ne produit que des liens stériles ; une compatibilité élevée sans contact ne produit rien, ou produit tard — les « belles au bois dormant » de Ke et al. 2015).

\begin{avis}
Argon & \st{Nuancée} (lecture de la note) & Produit de termes ; le contact humain n'y entre qu'à la partie 6 (« porteurs bilingues »), non dans $\lambda$. & Burt 2004\\ \midrule
Carbon & \st{Soutenue} (forme) & Rencontres $(\eta_0C_{\text{soc}}+\eta_1F)$ $\times$ sigmoïde de la compatibilité. & Burt 2004 ; Simonton 2004\\ \midrule
Neon & \st{Soutenue} (forme) & Contact $B^{\rho}$ $\times$ champ $e^{\eta\Phi}$ $\times$ cloche de surface $\times$ $K^{\theta}$ ; les protéines « réduites sans perte » au taux de Poisson conditionné par le contact et la compatibilité. & Burt 2004 ; Simonton 1979\\
\end{avis}

\begin{exemple}[ : deux laboratoires voisins]
Deux équipes travaillent dans le même bâtiment, l'une sur la vision par ordinateur, l'autre sur la rétine (contact élevé). Elles se croisent chaque jour sans rien produire ensemble tant que leurs descriptions n'ont pas d'invariant commun ($K$ faible). Le jour où l'une adopte un formalisme (filtres convolutifs) qui décrit aussi les champs récepteurs de l'autre, $K$ monte : la même fréquence de rencontres devient féconde. Exemple illustratif.
\end{exemple}

%---------------------------------------------------------------------
\subsection{Mesurer directement les invariances (message 6)}\label{sec:mesure-directe}
\begin{enonce}
Sur l'objection de circularité (« il faut une mesure du champ indépendante de l'issue ») : « Et pourquoi ne pas mesurer tout simplement les invariances par collision et rapprochement des structures mathématiques sous-jacentes ou des réseaux sous-tendant les idées ? » Sur l'homomorphisme de graphes typés comme équivalent rigoureux : « intéressant, creusons ça ». Sur « l'invariant conservé sous l'application qui aligne les deux réseaux, qui est déjà la systématicité de Gentner » : « qui les rapproche, en tout cas ».
\end{enonce}

\textbf{Explication.} La proposition est juste et elle est réalisable : c'est le score $K(u,v;t)$, calculé \emph{avant} l'issue. Elle échappe à la circularité à trois conditions. (1) Les réseaux $G_u$, $G_v$ sont construits sur les seuls documents antérieurs à $t$. (2) L'alphabet des relations abstraites $\pi:\mathcal R\to\mathcal R_{\text{abs}}$ et la table des incompatibilités sont fixés \emph{avant} de regarder les paires (sinon on règle l'alphabet pour que les liens connus s'alignent). (3) $K$ est calculé pour \emph{toutes} les paires candidates, réussies ou non. Reste une circularité résiduelle que Carbon nomme (annexe B, § 6.4) : les grandes découvertes créent souvent l'abstraction qui rend l'alignement possible ($\mathcal R_{\text{abs}}(t)\subsetneq\mathcal R_{\text{abs}}(t+1)$) ; un $K$ calculé avec l'alphabet d'avant ne la verra pas.

\emph{Quel objet mathématique ?} Un \emph{graphe typé} est un graphe muni d'un morphisme vers un graphe des types (Ehrig et al., 2006) ; un homomorphisme $f:G\to H$ préserve arêtes et types ; partiel, il peut omettre des nœuds ; injectif, c'est un plongement de sous-graphe ; « gradué », on maximise la part des relations préservées. Au-dessus, les \emph{foncteurs} entre catégories (ologs de Spivak et Kent) préservent aussi la composition des relations, et la \emph{naturalité} exprime qu'une transformation commute avec toutes les flèches — l'énoncé le plus proche de « l'invariance se préserve quelle que soit la projection ». Pour des réseaux de vocabulaires différents et bruités, la version continue est le transport optimal de Gromov-Wasserstein (partie~2). La nuance du directeur (« qui les rapproche ») est importante : la systématicité \emph{classe} les alignements ; elle ne dit pas qu'un alignement de haut score a une valeur causale. La question de suivi sur l'homomorphisme, servie aux modèles, sera lue en version 2.

\begin{avis}
Argon & \st{Soutenue} (lecture de la note) & Propose lui-même l'homomorphisme de graphes typés (H3) et son calcul exact par VF2++ ou programmation par ensembles de réponses ; tout est calculé sur les archives $\le t$. & Jüttner et Madarasi 2018 ; Vayer et al. 2019 ; Spivak et Kent 2012\\ \midrule
Carbon & \st{Soutenue}, avec la réserve de l'alphabet & Calcul exact en programmation en nombres entiers ou par Gromov-Wasserstein fusionné « sans accès aux étiquettes des nœuds » ; contrôle par permutation des types de relations ; réserve : genèse de nouvelles abstractions. & Vayer et al. 2020 ; Mac Lane 1971 ; Carey 2009\\ \midrule
Neon & \st{Soutenue} (lecture de la note) & Criblage par Gromov-Wasserstein fusionné, puis appariement exact (SME, McSplit) ; met en garde : isomorphisme de graphe $\neq$ invariance causale. & Vayer et al. 2019 ; Falkenhainer et al. 1989 ; Pearl 2009\\
\end{avis}

\begin{exemple}[ : le transfert eau–chaleur de SME]
Le moteur SME (Falkenhainer et al., 1989) aligne « l'eau s'écoule d'un bécher vers une fiole par un tuyau parce que la pression est plus haute dans le bécher » et « la chaleur passe du café au glaçon par une barre parce que la température du café est plus haute ». L'invariant conservé : $\textsc{cause}(\textsc{plus-grand}(\text{pression}_1,\text{pression}_2),\ \textsc{flux}(\ldots))$ $\mapsto$ $\textsc{cause}(\textsc{plus-grand}(\text{temp}_1,\text{temp}_2),\ \textsc{flux}(\ldots))$. L'attribut « liquide » de l'eau n'est pas transporté. Mesurer cette invariance « directement », c'est calculer cet alignement sur des réseaux extraits automatiquement — ce qu'aucun des trois modèles n'a encore éprouvé à grande échelle.
\end{exemple}

%---------------------------------------------------------------------
\subsection{Galison : comment modéliser la synchronisation des horloges ? (message 6)}\label{sec:galison}
\begin{enonce}
Sur Galison (2003), la synchronisation des horloges comme substrat commun à Einstein et à Poincaré : « oui, mais comment modéliser ça mathématiquement ? »
\end{enonce}

\textbf{Explication (proposition de la note, à éprouver).} Trois objets suffisent.
\begin{enumerate}[itemsep=1pt]
\item Un graphe biparti daté \emph{concepts–pratiques} : les concepts (« simultanéité », « éther », « temps local », « longitude ») et les pratiques ou instruments (« signal télégraphique horaire », « horloge électrique asservie », « réseau ferroviaire »), reliés quand un document de l'année $t$ les associe. $h_m(u,t)$ = part des documents sur $u$ qui mentionnent la pratique $m$.
\item Le champ $F(u,v;t)=\sum_m\alpha_m(t)h_m(u,t)h_m(v,t)$ entre « simultanéité » (physique) et « mesure du temps à distance » (géodésie, chemins de fer). Prédiction : $F$ monte avant 1905 et le rapprochement \emph{structurel} ($K^\star$, section~\ref{sec:champ-distance}) aussi : « simultanéité » acquiert la même relation \textsc{définie-par}(échange de signaux, correction du temps de transit) que « synchronisation ».
\item Le test : un cas n'est pas un test. Il faut un \emph{catalogue} de découvertes multiples et de singletons appariés (même époque, même discipline), et vérifier que $\Delta F$ et $\Delta K^\star$ précèdent plus souvent les multiples que les singletons (annexes A à C, § 5).
\end{enumerate}
Limites : les archives de l'époque ne sont pas un corpus numérisé homogène ; la thèse de Galison est une interprétation (il ne documente pas un lien direct entre tel brevet examiné et l'article de 1905) ; la modélisation ne la prouve pas, elle la rend réfutable sur d'autres cas.

\begin{avis}
Argon & \st{Soutenue} (Galison « a démontré » : trop fort) & Galison cité pour H5 ; son champ $H_t$ est la forme proposée ci-dessus. & Galison 2003 ; Kleinberg 2003\\ \midrule
Carbon & \st{Soutenue} & « Rôle exact » de la synchronisation ferroviaire « démontré » (même réserve) ; champ $F$ et test des jumeaux. & Galison 2003 ; Bikard et al. 2015\\ \midrule
Neon & \st{Soutenue} & « Confirme le rôle matériel » ; test cas–témoins des multiples sur arXiv et bioRxiv. & Galison 2003 ; Simonton 1979\\
\end{avis}

\begin{exemple}[ : ce que le modèle prédirait pour un multiple récent]
Deux prépublications indépendantes, à quelques semaines d'intervalle, sans auteurs communs ni citations mutuelles, introduisent le même couple de concepts. Le modèle prédit que, dans les années précédentes, les deux concepts avaient commencé à être décrits avec un même outil en forte croissance (un logiciel, une base de données, une méthode). Un singleton apparié n'aurait pas cette montée.
\end{exemple}

%---------------------------------------------------------------------
\subsection{Tshitoyan et la structure du langage (message 6)}\label{sec:tshitoyan}
\begin{enonce}
« Incroyable, qu'est-ce que cela implique pour le langage ? Il posséderait une structure mathématique qui miroiterait celle de l'univers ? Donc en étudiant les liens cachés dans le langage, des vrais liens mathématiques, pas des métaphores, on pourrait faire des découvertes ? Comment s'y prendre et combiner cela avec le reste ? » (Protocole X lancé sur cette question, GO-2026-11-23-24.)
\end{enonce}

\textbf{Explication.} Ce que l'article établit, ce qu'il n'établit pas, et comment s'en servir.
\begin{itemize}[itemsep=1pt]
\item \emph{Établi} : la littérature d'un domaine contient une connaissance latente que des représentations non supervisées extraient (structure du tableau périodique, relations structure–propriété) et qui suffit à classer des candidats pour une fonction mieux que le hasard, plusieurs années avant leur découverte (résumé de l'article).
\item \emph{Non établi} : que le langage « miroite l'univers ». La structure vient de ce que des chimistes ont mesuré et écrit ; elle hérite aussi de leurs habitudes. Sourati et Evans (2023) montrent que la prédiction des découvertes s'améliore nettement quand on modélise \emph{qui} écrit (la distribution des experts) : ces modèles réussissent en partie « en prédisant les prédictions humaines et les scientifiques qui les feront ». Une part du « signal » est donc sociale, non physique.
\item \emph{Comment s'en servir} : (1) passer de la cooccurrence à des couches de \emph{relations} (triplets typés, arbres d'opérateurs des équations), où la structure est explicite ; (2) aligner ces couches entre domaines (section~\ref{sec:mesure-directe}) ; (3) éprouver de façon prospective, contre des bases fortes — dans le concours Science4Cast, les meilleures méthodes de prédiction des liens futurs entre concepts utilisent des traits de réseau construits à la main plutôt qu'une approche de bout en bout (Krenn et al., 2023, résumé) ; (4) chercher aussi, comme Sourati et Evans, les hypothèses « étrangères », loin de la densité humaine, que la science ne visiterait pas d'elle-même.
\end{itemize}

\begin{avis}
Argon & \st{Soutenue} (« ont prouvé » : trop fort) & Cité pour H9 ; plongements qui prédisent des thermoélectriques. & Tshitoyan et al. 2019 ; Mikolov et al. 2013\\ \midrule
Carbon & \st{Nuancée} & Limite critique : la cooccurrence confond thème et isomorphisme relationnel. & Tshitoyan et al. 2019 ; Lu et al. 2019\\ \midrule
Neon & \st{Nuancée} & Les vecteurs codent d'abord la surface ; échecs sur les analogies relationnelles. & Tshitoyan et al. 2019 ; Rogers et al. 2017 ; Mitchell 2021\\
\end{avis}
Aucun des trois ne cite Sourati et Evans (2023), Krenn et Zeilinger (2020), Krenn et al. (2023) ni Rzhetsky et al. (2015), qui occupent directement la prédiction des liens futurs entre concepts (partie~3).

\begin{exemple}[ : une expérience à la Tshitoyan, mais structurelle]
Couper un corpus de physique en 2014. Construire pour chaque concept (a) un plongement de surface et (b) la forêt de ses équations anonymisées. Prédire les paires de concepts qui seront réunies en 2015–2018 avec (a) seul, (b) seul, et (a)+(b), contre les traits de réseau à la Science4Cast. Si (b) n'ajoute rien, l'idée de « vrais liens mathématiques » cachés dans le corpus est réfutée sous cette forme ; si (b) ajoute, on tient un instrument que la cooccurrence seule n'avait pas. (C'est, à peu de chose près, l'expérience d'Argon, annexe A, § 7.)
\end{exemple}
%---------------------------------------------------------------------
\subsection{Message 7 (versé après le mandat de cette note)}\label{sec:message7}
Le fichier des intuitions a reçu, après le mandat, un message 7 (vers 14:10Z à 15:00Z, après lecture des réponses de Google ; empreinte du fichier sur \texttt{main} le 24 novembre : \texttt{6f6397ee…}). La note ne le traite pas en entier ; elle indique où chaque réaction du directeur se rattache.
{\small
\begin{longtable}{@{}>{\raggedright\arraybackslash}p{6.6cm}>{\raggedright\arraybackslash}p{9.6cm}@{}}
\toprule \textbf{Réaction du directeur (résumée)} & \textbf{Où la note en traite ; ce qui reste ouvert}\\ \midrule\endhead
\bottomrule\endfoot
Maximiser le progrès de compression ; comment l'exploiter dans l'organisme ? & H10 et la fiche Schmidhuber : mesurer le progrès d'un apprenant fixe sur un jeu tenu à l'écart, non l'improbabilité. Conception à confier (pièce délicate).\\ \midrule
« Le papier dopamine chez les modèles de langage du MIT » & Référence non identifiée : \textbf{à vérifier} (titre, auteurs, identifiant), revue bibliographique.\\ \midrule
Symétries cachées du réel détectées par celles des cooccurrences (« oui ») ; représentations explicites des relations (« oui ») & H9 et section~\ref{sec:tshitoyan} : passer des cooccurrences aux couches de relations et d'équations.\\ \midrule
Percolation conceptuelle pure « à tester » ; transitions de phase « il faut les chercher » & H8 : percolation bipartie sur réseaux finis, paramètre d'ordre = taille de la composante hybride ; chercher une variation brusque contre un modèle nul de croissance lisse.\\ \midrule
Traction par les anomalies (« brillant ») & Partie 3, ajout 5 : une couche d'anomalies (résidus empiriques) dans le graphe.\\ \midrule
Structure mathématique préexistante contre alphabet relationnel humain ; deux phases (découvrir des objets formels, puis les appliquer) ; une troisième couche « méta » du graphe & Section~\ref{sec:mesure-directe} (circularité de l'alphabet) et Carbon, annexe B, § 6.4. La question neutre « méta-espace mathématique » existe sur \texttt{main} (\texttt{question-meta-espace-mathematique-v1.md}) ; réponses en version 2.\\ \midrule
« On ne peut pas trouver un principe variationnel ? » & Fiche « théorèmes de Noether discrets » : il en faut une dynamique (action discrète ou générateur de Markov) ; la question devient « existe-t-il un générateur mesurable de l'évolution du graphe ? ».\\ \midrule
La scission : « il faut le faire dans l'organisme » & H6 et partie 3, proposition 6 : pièce la moins occupée selon Neon, à vérifier par la revue avant emploi.\\
\end{longtable}}
%=====================================================================
\section{Cours accéléré : les travaux qui occupent déjà ces idées}
%=====================================================================
Chaque fiche suit le même plan : \textbf{ce que dit le travail}, avec un exemple ; \textbf{ce qu'il démontre et ce qu'il ne démontre pas} ; \textbf{son lien avec les hypothèses}. Les fiches s'appuient sur les résumés lus en ligne et sur la connaissance générale de ces travaux ; aucun chiffre n'y est donné. Elles sont groupées en quatre familles : l'analogie et la créativité ; la découverte assistée par la littérature et la science de la science ; l'intérêt et la surprise ; la sociologie, l'histoire et la philosophie ; puis les outils mathématiques.

\subsection*{A. Analogie et créativité}

\subsubsection{Koestler (1964) : la bisociation}
\cours{Ce que dit le travail.} \emph{The Act of Creation} propose qu'un même acte mental est au cœur de l'humour, de la découverte et de l'art : la \emph{bisociation}, percevoir une situation à la fois dans deux cadres de référence (« matrices ») chacun cohérent mais habituellement incompatibles. Le mot d'esprit fait se heurter les deux cadres (rire) ; la découverte les fusionne (« eurêka ») ; l'art les juxtapose. Exemple emblématique : Archimède, dans son bain, voit d'un même regard le corps qui déplace l'eau et le problème de mesurer le volume de la couronne.
\cours{Ce qu'il démontre et ne démontre pas.} C'est une théorie descriptive, appuyée sur des anecdotes choisies parmi des réussites célèbres. Elle ne fournit ni mesure de la distance entre matrices, ni critère qui sépare les bisociations fécondes des stériles, ni test. Elle souffre du biais du survivant que H1 doit éviter : on ne raconte que les bains qui ont donné un théorème.
\cours{Lien.} C'est l'ancêtre direct de H1 (deux cadres éloignés, un point commun caché). Ce que H1 ajoute, et que Koestler n'a pas : deux distances mesurables et la condition de structure (H3).

\subsubsection{Gentner (1983) ; Falkenhainer, Forbus et Gentner (1989) : l'appariement de structures, la systématicité, le moteur SME}
\cours{Ce que dit le travail.} Une analogie projette des \emph{relations} d'un domaine source vers un domaine cible, non des attributs. « L'atome est comme le système solaire » transporte « gravite autour de », « attire », « plus massif que » ; il ne transporte pas « chaud » ni « jaune ». Le \emph{principe de systématicité} : parmi les correspondances possibles, on préfère celle qui emporte un système connexe de relations gouverné par des relations d'ordre supérieur (\textsc{cause}, \textsc{implique}). Le \emph{moteur SME} (1989) le met en algorithme : il forme des hypothèses locales entre prédicats identiques, les assemble en correspondances globales cohérentes (une à une, « connexion parallèle » : si deux relations sont appariées, leurs arguments le sont), et les note en faisant « ruisseler » le score des relations d'ordre supérieur vers leurs arguments. Exemple canonique : écoulement d'eau entre deux récipients et flux de chaleur entre café et glaçon (section~\ref{sec:mesure-directe}).
\cours{Ce qu'il démontre et ne démontre pas.} Des expériences de psychologie montrent que les gens jugent plus « solides » les analogies qui partagent des relations que celles qui partagent des attributs, et SME reproduit ces jugements sur les cas étudiés. Ce n'est pas démontré : que la systématicité prédise la \emph{fécondité scientifique} à grande échelle. Et SME travaille sur des représentations écrites à la main, avec la contrainte d'identité des prédicats : tout dépend de qui a écrit \textsc{flux} des deux côtés. C'est la critique de Chalmers, French et Hofstadter (1992) : le vrai problème est de \emph{construire} la représentation.
\cours{Lien.} H3 en est la reformulation directe ; H6 en découle (les « différences alignables » de Gentner et Markman, 1997). La critique de la représentation est exactement la circularité de la section~\ref{sec:mesure-directe} : l'alphabet des relations doit être fixé avant de chercher.

\subsubsection{Hofstadter et Mitchell : Copycat}
\cours{Ce que dit le travail.} Copycat (Mitchell et Hofstadter, 1990 ; Mitchell, 1993 ; Hofstadter et le Fluid Analogies Research Group, 1995) fait des analogies dans un micro-domaine de chaînes de lettres : « si \texttt{abc} devient \texttt{abd}, que devient \texttt{xyz} ? ». Pas de \texttt{z} suivant : il faut « glisser » un concept (successeur $\to$ prédécesseur, fin $\to$ début) pour trouver la réponse élégante \texttt{wyz}. L'architecture fait \emph{en même temps} la perception (quelles structures voir dans la chaîne) et l'appariement, par de nombreux petits agents stochastiques ; un réseau de concepts (\emph{slipnet}) dont les distances changent selon le contexte ; une \emph{température} qui mesure la qualité de la structure courante et règle le hasard (haute : exploration ; basse : engagement).
\cours{Ce qu'il démontre et ne démontre pas.} Il montre qu'une architecture où la représentation se construit pendant l'appariement produit la variété des réponses humaines, avec des fréquences plausibles, dans ce micro-domaine. Il ne montre rien sur des domaines scientifiques ; il n'a pas été mis à l'échelle. Mitchell (2021) dresse le bilan récent : l'abstraction et l'analogie restent un point faible des systèmes d'apprentissage.
\cours{Lien.} Le « glissement » conceptuel est la projection de H5 et de la section~\ref{sec:champ-distance} : un changement de contexte rapproche des concepts lointains. La température est l'« agitation » du message 2, avec un rôle précis : plus de hasard quand la structure trouvée est mauvaise.

\subsubsection{Fauconnier et Turner : le mélange conceptuel}
\cours{Ce que dit le travail.} Un \emph{réseau d'intégration} (Fauconnier et Turner, 1998 ; 2002) a deux espaces d'entrée, un espace générique (ce qu'ils ont en commun) et un espace \emph{mélangé} qui reçoit une partie de chacun et développe une structure émergente que n'a aucun des deux. Exemple : l'énigme du moine bouddhiste qui monte la montagne un jour et la redescend le lendemain : y a-t-il un point où il se trouve à la même heure les deux jours ? On superpose les deux journées en une seule où deux moines se croisent : la réponse devient évidente. Autre exemple : « ce chirurgien est un boucher » : l'incompétence n'est ni dans l'entrée « chirurgie » ni dans l'entrée « boucherie », elle naît du mélange (moyens de boucher, fin de chirurgien).
\cours{Ce qu'il démontre et ne démontre pas.} C'est un cadre descriptif très large, qui couvre métaphores, contrefactuels, interfaces (le « bureau » de l'ordinateur). On lui reproche souvent de pouvoir décrire presque tout après coup et de peu prédire quels mélanges se forment (critique générale, référence à vérifier). Les « principes d'optimalité » qu'ils proposent ne sont pas des quantités mesurées.
\cours{Lien.} L'idée fille de H7 est un espace mélangé ; l'espace générique est la « dimension cachée » de H1. Ce que le mélange ajoute à Gentner : la structure \emph{émergente}, qui n'est dans aucun parent — c'est ce qui rend une idée fille plus qu'une correspondance.

\subsection*{B. Découverte assistée par la littérature et science de la science}

\subsubsection{Swanson : la découverte fondée sur la littérature}
\cours{Ce que dit le travail.} Swanson (1986, deux articles) appelle « connaissance publique non découverte » les liens logiques entre littératures qui ne se lisent pas : si $A$ (huile de poisson) agit sur $B$ (viscosité sanguine) et si $B$ est en cause dans $C$ (syndrome de Raynaud), le lien $A\!-\!C$ est « public » mais personne ne l'a écrit. Il en tire des hypothèses, dont celle de l'huile de poisson, ensuite éprouvée en essai contrôlé (DiGiacomo et al., 1989). Avec Smalheiser, il construit un outil interactif (Arrowsmith ; Swanson et Smalheiser, 1997) pour trouver ces littératures complémentaires.
\cours{Ce qu'il démontre et ne démontre pas.} Il démontre l'existence de littératures logiquement liées et disjointes, et qu'on peut en tirer des hypothèses vérifiables. Il ne démontre pas un rendement général : les cas célèbres ont été conduits par un expert qui choisissait les termes $B$ ; l'automatisation produit une explosion de candidats ; l'évaluation des systèmes est notoirement difficile (on ne sait pas ce qu'on a manqué).
\cours{Lien.} C'est H2 au sens (a) ; c'est aussi la base « Swanson ABC » que les trois modèles exigent de battre. La leçon de l'exemple : le pivot doit être une grandeur ou un mécanisme, pas un mot.

\subsubsection{Uzzi, Mukherjee, Stringer et Jones (2013) : combinaisons atypiques et impact}
\cours{Ce que dit le travail.} Pour chaque paire de revues citées ensemble dans la bibliographie d'un article, on compare la fréquence observée de cette paire dans toute la littérature à celle d'un modèle nul qui redistribue les citations au hasard en gardant leurs nombres : un score $z$ élevé = paire conventionnelle, négatif = atypique. Chaque article est décrit par sa conventionnalité médiane et sa nouveauté de queue. Les articles les plus cités ne sont pas les plus nouveaux : ce sont ceux qui combinent un socle très conventionnel et une pointe de combinaisons atypiques ; les équipes y parviennent plus souvent que les auteurs seuls.
\cours{Ce qu'il démontre et ne démontre pas.} Une association robuste, sur l'ensemble du Web of Science, entre ce profil et l'impact mesuré par les citations. Pas de causalité ; une distance au niveau des \emph{revues}, non des idées ; l'impact en citations n'est pas la fécondité ; le résultat dépend du modèle nul choisi.
\cours{Lien.} C'est l'appui le plus cité de H1 et la meilleure nuance de H4 (mélange, non distance moyenne). Base de comparaison B3 chez Carbon.

\subsubsection{Shi et Evans (2023) : surprises de contenus et de contextes}
\cours{Ce que dit le travail.} Sur un très grand corpus d'articles et de brevets (sciences de la vie, sciences physiques, inventions), un modèle d'hypergraphe apprend quelles combinaisons de \emph{contenus} (mots-clés) et de \emph{contextes} (revues citées) sont attendues ; la surprise d'une combinaison réalisée prédit un impact hors norme. Ces avancées surprenantes naissent « à travers » plutôt qu'« à l'intérieur » des chercheurs et des équipes, le plus souvent quand des scientifiques d'un domaine publient des résultats de résolution de problème pour le public d'un domaine lointain (résumé).
\cours{Ce qu'il démontre et ne démontre pas.} Que la surprise \emph{mesurée par un modèle} de ce qui est attendu est associée à l'impact, et que l'origine des auteurs compte. Pas de causalité ; la surprise est relative au modèle ; l'impact reste en citations.
\cours{Lien.} C'est la version moderne et mesurée de H1 (surprise de surface) et de la section~\ref{sec:rencontres} (l'« expédition » d'un chercheur d'un domaine vers un autre est un terme de contact). Il ne mesure pas la structure (H3).

\subsubsection{Sourati et Evans (2023) : une intelligence artificielle « consciente des humains »}
\cours{Ce que dit le travail.} Des modèles non supervisés entraînés sur des inférences simulées « cognitivement accessibles aux experts » (marches aléatoires qui passent par les auteurs autant que par les concepts) prédisent beaucoup mieux les découvertes futures que des modèles fondés sur le seul contenu, surtout quand la littérature pertinente est clairsemée. Ils réussissent « en prédisant les prédictions humaines et les scientifiques qui les feront ». En réglant le modèle pour \emph{éviter} la foule, on obtient des hypothèses « étrangères », scientifiquement prometteuses mais peu susceptibles d'être imaginées avant longtemps (résumé).
\cours{Ce qu'il démontre et ne démontre pas.} Que la distribution des experts est un prédicteur majeur de ce qui sera découvert, et que l'on peut viser hors de cette distribution. Ce qui reste à établir : la valeur réelle des hypothèses « étrangères », jugée par une épreuve indépendante.
\cours{Lien.} Il occupe en partie H5 (densité humaine comme facteur de rencontre) et donne la meilleure réponse au message 6 sur Tshitoyan : une part de ce que « sait » le langage est la géographie des scientifiques. Base de comparaison obligatoire pour toute expérience de prédiction.

\subsubsection{Rzhetsky, Foster, Foster et Evans (2015) : choisir les expériences}
\cours{Ce que dit le travail.} Le savoir biomédical est modélisé comme un réseau (molécules et relations) que la science découvre par une suite d'expériences. Sur un très grand corpus d'articles biomédicaux, les chercheurs suivent des stratégies conservatrices, explorant le voisinage de molécules centrales déjà importantes ; ces stratégies servent sans doute les carrières mais ralentissent la découverte, surtout dans les domaines mûrs, où plus de risque et moins de redondance accéléreraient la découverte du réseau (résumé).
\cours{Ce qu'il démontre et ne démontre pas.} Une description des choix réels et, \emph{dans le modèle}, un gain d'efficacité des stratégies plus risquées. Le gain réel dépend du modèle de coût et de réussite des expériences.
\cours{Lien.} H4 côté production (le « trop près » est le choix dominant), H7 (s'arrimer aux nœuds est conservateur) et le terme de coût que Carbon et Neon ajoutent (partie 6 de leurs réponses).

\subsubsection{Krenn et Zeilinger (2020) ; Krenn et al. (2023) : prédire les liens futurs entre concepts}
\cours{Ce que dit le travail.} SemNet (2020) construit, à partir de la littérature de physique quantique, un réseau évolutif dont les nœuds sont des concepts, reliés quand deux concepts sont étudiés ensemble ; un réseau de neurones prédit les tendances futures et propose des « graines d'idées » personnalisées et surprenantes (résumé). Le concours Science4Cast (2023) en fait un banc d'essai pour l'intelligence artificielle elle-même : prédire l'état futur d'un réseau sémantique de concepts ; des méthodes variées, de la statistique pure à l'apprentissage pur, sont comparées ; « étonnamment », les plus puissantes reposent sur un ensemble de traits de réseau soigneusement choisis plutôt que sur une approche de bout en bout (résumé).
\cours{Ce qu'il démontre et ne démontre pas.} Que la naissance de liens entre concepts est en partie prévisible à partir de la topologie passée ; que des traits simples sont des bases redoutables. Ce n'est pas une mesure de l'intérêt ni de la fécondité des liens prédits.
\cours{Lien.} Ces travaux occupent la question « prédire les mariages » (H1, H2, H8) sans la couche de structure (H3). Toute expérience du programme doit battre leurs traits (degrés, voisins communs, âge des nœuds), en plus de Swanson.

\subsubsection{Tshitoyan et al. (2019) : la connaissance latente dans les plongements}
\cours{Ce que dit le travail.} Voir section~\ref{sec:tshitoyan}. Des plongements de mots appris sans supervision sur des résumés de science des matériaux capturent la structure du tableau périodique et des relations structure–propriété ; une méthode non supervisée recommande des matériaux pour des applications fonctionnelles plusieurs années avant leur découverte ; « la connaissance latente des découvertes futures est dans une large mesure enfouie dans les publications passées » (résumé).
\cours{Ce qu'il démontre et ne démontre pas.} Une validation rétrospective (on entraîne sur le passé, on regarde ce qui a été publié ensuite) dans un domaine où la propriété cible a un vocabulaire stable. Il ne démontre pas la découverte d'un principe nouveau, ni que la structure apprise soit physique plutôt que bibliographique.
\cours{Lien.} H9 sous sa forme faible ; l'exemple de la section~\ref{sec:tshitoyan} propose la forme structurelle.
\subsection*{C. Intérêt et surprise}

\subsubsection{Davis (1971) : « That's interesting! »}
\cours{Ce que dit le travail.} Une théorie est jugée intéressante non parce qu'elle est vraie, mais parce qu'elle \emph{nie une hypothèse tenue par son public}, sans nier tout son cadre. Davis classe les formes de cette négation : ce qui semble désorganisé est organisé (et l'inverse) ; ce qui semble individuel est collectif ; ce qui semble local est général ; ce qui semble stable change ; ce qui semble une cause est un effet ; ce qui semble lié est indépendant ; et, pour plusieurs phénomènes, ce qui semble semblable est différent (et l'inverse). Exemple de ce type (illustration de la note) : soutenir que la déviance, qu'on croyait pathologique, remplit une fonction sociale.
\cours{Ce qu'il démontre et ne démontre pas.} C'est une analyse conceptuelle tirée de théories classiques de la sociologie, non une étude empirique. Elle montre que l'intérêt est \emph{relatif à un public} ; elle ne dit rien de la vérité ni de la fécondité.
\cours{Lien.} C'est la meilleure source pour la question 1 (l'intérêt). Ses deux formes « ce qui semble différent est semblable » et « ce qui semble semblable est différent » sont H1 et H6 : connexion et distinction sont les deux moitiés d'une même liste.

\subsubsection{Schmidhuber : le progrès de compression}
\cours{Ce que dit le travail.} Pour un observateur qui apprend à comprimer ses données (à les prédire), la beauté subjective est la compressibilité, et l'\emph{intérêt} est la \emph{dérivée} de la compressibilité : le progrès que l'observateur fait en apprenant (Schmidhuber, 2009 ; 2010). Un agent curieux maximise ce progrès comme récompense intrinsèque. Une donnée déjà comprise n'apporte pas de progrès (ennui) ; un bruit incompressible non plus (rien à apprendre) ; l'intéressant est entre les deux.
\cours{Ce qu'il démontre et ne démontre pas.} Un cadre formel unifié, mis en œuvre dans des agents d'apprentissage par renforcement curieux. Il n'est pas validé sur des jugements humains d'intérêt, et le progrès de compression d'un apprenant réel est difficile à calculer.
\cours{Lien.} H10 (le bon signal n'est pas l'improbabilité mais le progrès) et H7 (l'unification est une compression ; la surprise après coup est le moment où la compression s'améliore d'un coup).

\subsubsection{Itti et Baldi : la surprise bayésienne}
\cours{Ce que dit le travail.} La surprise d'une donnée est la distance entre les croyances de l'observateur avant et après l'avoir vue : $S=D_{\mathrm{KL}}(P(M\mid D)\,\Vert\,P(M))$ sur un ensemble de modèles $M$. Elle se distingue de l'improbabilité de Shannon : une neige d'écran est très improbable à chaque image mais ne fait plus réviser aucune croyance une fois reconnue comme bruit. Itti et Baldi (2005) l'emploient pour détecter des événements surprenants dans des vidéos ; en 2009, ils montrent que les regards humains sur des vidéos se portent vers les lieux de forte surprise bayésienne, mieux que vers ceux qu'indiquent d'autres mesures.
\cours{Ce qu'il démontre et ne démontre pas.} Une mesure calculable qui prédit l'attention visuelle de bas niveau. Rien sur la surprise sémantique d'une idée, qui suppose de modéliser les croyances d'une communauté.
\cours{Lien.} Version mesurable de la « surprise » de H1 et de H10 ; Carbon et Neon en font le score d'intérêt perçu $I_{\text{perc}}$ (annexes B et C, § 2).

\subsubsection{Berlyne : variables collatives et courbe en U renversé}
\cours{Ce que dit le travail.} Berlyne (1960) nomme « collatives » les propriétés d'un stimulus qui exigent de le comparer à autre chose : nouveauté, complexité, surprise, incongruité, incertitude. Elles élèvent l'éveil ; la valeur hédonique suit une courbe en U renversé (courbe de Wundt) : trop peu, ennui ; trop, aversion. Il distingue curiosité perceptive et épistémique, diversive (fuir l'ennui) et spécifique (combler un manque). En 1970, il étudie expérimentalement nouveauté, complexité et valeur hédonique.
\cours{Ce qu'il démontre et ne démontre pas.} Une organisation durable des recherches sur la curiosité. La forme en U renversé n'est pas trouvée partout (selon le stimulus et la mesure, la préférence peut être monotone) ; elle porte sur l'intérêt, pas sur la fécondité.
\cours{Lien.} Appui de H4 pour la question 1 seulement ; Loewenstein (1994, cité par Neon) en donne une relecture par le « manque d'information ».

\subsection*{D. Sociologie, histoire et philosophie}

\subsubsection{Burt (2004) : trous structuraux et bonnes idées}
\cours{Ce que dit le travail.} Un \emph{trou structural} est l'absence de lien entre deux groupes denses ; le \emph{courtier} qui a des contacts des deux côtés voit plus tôt des façons de faire et des idées que chaque groupe ignore. Chez des cadres de la chaîne d'approvisionnement d'une grande entreprise d'électronique, les idées d'amélioration proposées par ceux dont le réseau franchit des trous structuraux sont jugées plus précieuses par la direction ; Burt parle d'un « avantage de vision ». Mesure : la \emph{contrainte} d'un réseau personnel (faible quand les contacts ne se connaissent pas entre eux).
\cours{Ce qu'il démontre et ne démontre pas.} Une association, dans une organisation, entre position de courtage et valeur \emph{jugée} des idées. Ni en science, ni sur la valeur réalisée, ni causale (les gens qui ont de bonnes idées peuvent aussi construire des réseaux variés).
\cours{Lien.} H8 (surface de contact) et section~\ref{sec:rencontres} : c'est le facteur « rencontres » ; les « deux réseaux denses qui ne se touchent que par les parents » du message 2 sont littéralement des groupes séparés par un trou structural.

\subsubsection{Merton (1961) et Simonton (1979 ; 2004) : les découvertes multiples}
\cours{Ce que dit le travail.} Ogburn et Thomas (1922) dressent une liste d'inventions et de découvertes faites indépendamment par plusieurs personnes. Merton (1961) en tire une thèse : les découvertes multiples sont la règle plutôt que l'exception, et beaucoup de « singletons » sont des multiples devancés (l'autre équipe a abandonné en apprenant la publication). Simonton (1979) traite ces catalogues statistiquement : la distribution du nombre de découvreurs s'ajuste à un modèle de Poisson, ce qui est compatible avec une production au hasard par de nombreux chercheurs exposés aux mêmes prérequis, sans qu'il faille invoquer une inévitabilité de l'« esprit du temps ». En 2004, il développe une théorie de la créativité scientifique comme configuration de hasard sous contraintes.
\cours{Ce qu'il démontre et ne démontre pas.} Que les multiples sont fréquents et que leur distribution est compatible avec un modèle de hasard. Il ne départage pas complètement « hasard » et « esprit du temps » : les deux prédisent des multiples ; et les catalogues historiques mêlent des découvertes qui ne sont pas vraiment identiques.
\cours{Lien.} H5 (b) et (c) : le modèle de Poisson est exactement celui de Carbon et Neon ; le champ $F$ doit faire mieux que le hasard \emph{conditionné aux prérequis}.

\subsubsection{Galison (2003) : les horloges d'Einstein et les cartes de Poincaré}
\cours{Ce que dit le travail.} Une histoire croisée de la technique et de la physique : à la fin du \textsc{xix}\ieme{} siècle, coordonner le temps par signaux électriques est un enjeu de chemins de fer, de télégraphe, de cartographie et d'empire. Poincaré dirige des travaux de détermination des longitudes et réfléchit à la convention de simultanéité ; Einstein, expert à l'office des brevets de Berne, voit passer des dispositifs de synchronisation. Tous deux arrivent à définir la simultanéité par un échange de signaux corrigé du temps de transit : une notion à la fois technique, philosophique et physique.
\cours{Ce qu'il démontre et ne démontre pas.} Une thèse historique documentée : le substrat technique commun rend la reformulation de la simultanéité pensable. Il ne démontre pas un lien causal direct entre un brevet examiné et l'article de 1905, et il traite un cas.
\cours{Lien.} H5 (le cadre technique qui rapproche deux réseaux) et section~\ref{sec:galison} (sa modélisation).

\subsubsection{Kauffman (1995 ; 2000) et Tria, Loreto, Servedio et Strogatz (2014) : le possible adjacent}
\cours{Ce que dit le travail.} Le \emph{possible adjacent} est l'ensemble de ce qui est à un pas de l'actuel : chaque nouveauté réalisée agrandit l'ensemble des nouveautés possibles (Kauffman ; le terme est développé dans \emph{Investigations}, 2000 ; \emph{At Home in the Universe}, 1995, cité par Argon, en porte les idées — lequel l'introduit, à vérifier). Tria et al. (2014) le modélisent par une urne de Pólya généralisée : tirer une couleur nouvelle ajoute à l'urne des couleurs jamais vues. Le modèle prédit la loi de Heaps (rythme d'apparition des nouveautés), la loi de Zipf (distribution de ce qui est exploré) et des corrélations entre nouveautés (l'une prépare l'autre), qu'ils vérifient sur des données d'activité humaine : modifications de Wikipédia, étiquettes d'annotation, suite des mots dans des textes, écoute de morceaux nouveaux (résumé).
\cours{Ce qu'il démontre et ne démontre pas.} Un modèle simple qui reproduit des lois statistiques observées. D'autres modèles reproduisent aussi Heaps et Zipf ; le modèle ne dit pas \emph{quelle} nouveauté vient ensuite ni pourquoi.
\cours{Lien.} H8 (dynamique d'expansion) et la section~\ref{sec:champ-distance} : un cadre nouveau ajoute d'un coup des « couleurs » à l'urne.

\subsubsection{Kitcher (1981) : l'unification explicative}
\cours{Ce que dit le travail.} Expliquer, c'est dériver beaucoup de phénomènes acceptés d'un petit nombre de \emph{schémas d'argument} répétés. La meilleure systématisation (le « magasin explicatif ») est celle qui minimise le nombre de schémas pour un maximum de conclusions. Exemples : les \emph{Principia} de Newton (un schéma pour la chute, les orbites, les marées), Darwin (un schéma pour des faits de répartition, d'anatomie, de fossiles).
\cours{Ce qu'il démontre et ne démontre pas.} Une analyse philosophique de ce que fait une bonne explication, discutée (par exemple sur les explications qui unifient sans expliquer la cause). Pas de mesure opératoire : la compression par longueur de description est une formalisation proposée par d'autres (Carbon, Neon), non par Kitcher.
\cours{Lien.} Critère (a) de H7 ; la Lune et la pomme est son exemple de base.

\subsubsection{Poincaré (1908) : \emph{Science et méthode}, « L'invention mathématique »}
\cours{Ce que dit le travail.} Inventer, c'est choisir parmi les combinaisons possibles ; les combinaisons stériles ne se présentent même pas à l'esprit. Les faits dignes d'étude sont « ceux qui nous révèlent des parentés insoupçonnées entre d'autres faits, connus depuis longtemps, mais qu'on croyait à tort étrangers les uns aux autres. Parmi les combinaisons que l'on choisira, les plus fécondes seront souvent celles qui sont formées d'éléments empruntés à des domaines très éloignés » (p.~51 de l'édition définitive, transcription de Wikisource). Le récit célèbre : à Coutances, « au moment où je mettais le pied sur le marche-pied, l'idée me vint […] que les transformations dont j'avais fait usage pour définir les fonctions fuchsiennes sont identiques à celles de la Géométrie non-euclidienne » (p.~54). Il attribue au travail inconscient un tri par la sensibilité esthétique : les combinaisons utiles sont « les plus belles ».
\cours{Ce qu'il démontre et ne démontre pas.} Un témoignage introspectif d'une grande valeur, non une étude. L'introspection est sujette à la reconstruction après coup.
\cours{Lien.} C'est, mot pour mot, H1 en 1908 (« éléments empruntés à des domaines très éloignés ») ; l'épisode de Coutances est H3 (une \emph{identité de transformations} entre deux domaines, c'est-à-dire un invariant de groupe) ; le tri esthétique est une forme de H7. Le directeur cite Poincaré au message 4 pour la relativité : c'est le même homme et la même idée d'invariance par un groupe.

\subsection*{E. Outils mathématiques}

\subsubsection{La distance de Gromov-Wasserstein}
\cours{Ce que dit le travail.} Pour comparer deux objets qui ne vivent pas dans le même espace (deux réseaux de vocabulaires différents), on ne compare pas leurs points mais leurs \emph{distances internes}. Mémoli (2011) définit, pour deux espaces métriques mesurés $(X,d_X,\mu)$ et $(Y,d_Y,\nu)$,
\[
GW_p(X,Y)=\inf_{\pi\in\Pi(\mu,\nu)}\Big(\iint |d_X(x,x')-d_Y(y,y')|^p\,d\pi(x,y)\,d\pi(x',y')\Big)^{1/p},
\]
où $\pi$ est un couplage (un « appariement flou » de $X$ à $Y$). Elle est nulle quand les deux espaces sont isométriques : elle ignore les étiquettes et ne voit que la forme. Peyré, Cuturi et Solomon (2016) la rendent calculable par régularisation entropique ; Vayer et al. (2019) la « fusionnent » avec une distance entre attributs des nœuds (Gromov-Wasserstein fusionnée), pour graphes étiquetés.
\cours{Ce qu'elle démontre et ne démontre pas.} C'est une vraie distance entre espaces (à isométrie près), avec des propriétés démontrées. Son calcul est un problème quadratique non convexe, de la famille de l'affectation quadratique, difficile en général : les algorithmes trouvent des optimums locaux ; les symétries d'un graphe donnent plusieurs solutions équivalentes ; le résultat dépend de la distance interne choisie (plus court chemin, diffusion) et de la pondération des types de relations.
\cours{Lien.} C'est $d_{\text{struct}}$ chez les trois modèles et la version continue de la « mesure directe des invariances » (section~\ref{sec:mesure-directe}). Elle mesure une ressemblance de forme ; elle ne dit pas quelle relation porte l'analogie, d'où l'intérêt d'extraire ensuite un appariement discret $\phi$.

\subsubsection{Les homomorphismes de graphes typés}
\cours{Ce que dit le travail.} Un graphe $G=(V,E)$ ; un \emph{typage} est un morphisme vers un graphe des types $T$ qui donne à chaque nœud et à chaque arête un type (Ehrig et al., 2006). Un \emph{homomorphisme} $f:G\to H$ envoie les nœuds sur des nœuds en préservant les arêtes et leurs types : $(a,r,b)\in E_G\Rightarrow(f(a),r,f(b))\in E_H$. Variantes : partiel (on n'apparie qu'une partie), injectif (plongement de sous-graphe), gradué (on maximise la part des arêtes préservées, c'est le plus grand sous-graphe commun). Au niveau des catégories, les graphes sont des foncteurs, les homomorphismes des transformations naturelles ; les ologs de Spivak et Kent (2012) utilisent des foncteurs entre schémas de connaissances.
\cours{Ce qu'il démontre et ne démontre pas.} La théorie est riche (Hell et Nešetřil, 2004) ; elle comprend des résultats de complexité : décider s'il existe un homomorphisme vers un graphe fixé $H$ est polynomial si $H$ est biparti et NP-complet sinon, pour les graphes non orientés sans boucle (Hell et Nešetřil, 1990). Ce que l'outil manque d'une analogie réelle : les relations quantitatives (« décroît comme l'inverse du carré »), l'identité stricte des types (il faut un alphabet abstrait commun), le sens de l'inférence (source vers cible), et la causalité.
\cours{Lien.} C'est l'équivalent rigoureux de la « serrure et clé » de H3 proposé par Argon et retenu par le directeur au message 6 ; la question de suivi le met à l'épreuve.

\subsubsection{Complément : des théorèmes de Noether discrets}
\cours{Ce que dit le travail.} Marsden et West (2001) traitent la mécanique par des principes variationnels \emph{discrets} (intégrateurs variationnels), avec un théorème de Noether discret : une symétrie de l'action discrète donne une quantité exactement conservée par le schéma. Baez et Fong (2013) prouvent une version pour les processus de Markov : une observable commute avec le générateur si et seulement si son espérance \emph{et} son écart-type restent constants dans le temps pour tout état initial (résumé).
\cours{Ce qu'il démontre et ne démontre pas.} Que « discret » n'empêche pas Noether. Il faut une dynamique : une action discrète, ou un générateur de Markov. Rien n'établit qu'une dynamique des idées ait une telle forme.
\cours{Lien.} Nuance l'élagage de Noether par les trois modèles (H3) : la question ouverte n'est pas « graphe discret », mais « existe-t-il un générateur mesurable de l'évolution du graphe des concepts dont une symétrie serait observable ? ». Si l'on modélise la naissance des liens par une chaîne de Markov (ce que font les modèles d'urne et de processus ponctuel), la version de Baez et Fong devient au moins énonçable.
%=====================================================================
\section{Accords et désaccords entre Argon, Carbon et Neon ; propositions retenues}
%=====================================================================

\subsection{Ce sur quoi les trois s'accordent}
\begin{enumerate}[itemsep=2pt]
\item \textbf{Séparer intérêt, vérité-fécondité et production} (§ 1 des trois réponses), et donner à l'idée fille deux scores distincts : un intérêt perçu (surprise de surface $\times$ compatibilité) et une fécondité (compression, disruption, reprise indépendante).
\item \textbf{Couper H3 en deux} : l'alignement des réseaux amont est soutenu (Gentner, SME) ; Noether et le lagrangien sont élagués, remplacés par la préservation des relations sous une application partielle.
\item \textbf{Une même architecture} : graphe daté à trois couches ; distance de surface par cosinus et divergence de Jensen-Shannon ; distance de structure par Gromov-Wasserstein (fusionnée) sur des voisinages anonymisés ; score serrure et clé = couverture $\times$ fidélité $\times$ pénalité de conflits ; champ = co-exposition bilinéaire à des instruments en croissance ; intensité de naissance estimée par un modèle de Cox à covariables datées avec échantillonnage cas–témoins ; percolation bipartie avec un seuil.
\item \textbf{Une même validation} : coupure temporelle stricte ; bases simples (attachement préférentiel, Adamic-Adar, Swanson ABC, surface seule) ; biais du survivant traité par les mariages \emph{tentés} et restés stériles (zéro reprise indépendante, essais cliniques négatifs, rétractations) ; découvertes multiples contre singletons appariés ; jugement d'experts en aveugle selon un plan croisé surprise $\times$ compatibilité ; préenregistrement.
\item \textbf{Les mêmes oublis} : les anomalies empiriques (Kuhn) qui fissurent le réseau de l'extérieur ; les coûts matériels et institutionnels (pénalité de l'interdisciplinaire, effet Matthieu, biais contre la nouveauté) ; Carbon et Neon ajoutent la \emph{création} de nouvelles abstractions (l'alphabet des relations grandit) ; Neon ajoute les découvertes non combinatoires (théorèmes d'impossibilité), les « belles au bois dormant » et la causalité interventionnelle.
\end{enumerate}

\paragraph{Pourquoi cet accord pèse peu.} Les trois sont des modèles d'une même maison ; la consigne en sept parties dictait la structure ; et la question servie était jugée non neutre (avertissement 1). Leurs verdicts « soutenue » reprennent largement les hypothèses telles qu'on les leur a présentées : par la règle de symétrie du programme (R4), un accord aussi propre appelle la même méfiance qu'un désaccord. Ce qui a de la valeur est ce qu'ils ajoutent en propre (expériences, contrôles, contre-indices), non le décompte de leurs accords.

\subsection{Où ils divergent}
{\small
\begin{longtable}{@{}>{\raggedright\arraybackslash}p{2.6cm}>{\raggedright\arraybackslash}p{4.25cm}>{\raggedright\arraybackslash}p{4.25cm}>{\raggedright\arraybackslash}p{4.25cm}@{}}
\toprule \textbf{Point} & \textbf{Argon} & \textbf{Carbon} & \textbf{Neon}\\ \midrule\endhead
\bottomrule\endfoot
Ton général & « Soutenue » partout ; « preuve causale directe » (§ 7.3) & Nuances ciblées (H2, H4, H8, H9, H10) & Le plus critique : contre-indices sur H1, H4, H7, H9, H10\\ \midrule
Ferromagnétisme, protéines (H5) & « Métaphore partielle » en § 1, mais \emph{retenus} en § 2.8 (champ d'Ising, complémentarité stérique) : incohérence interne & Champ et polarisation gardés « sous forme statistique », protéines réduites & Écartés tous deux\\ \midrule
H4 & U renversé simple & Mélange socle conventionnel $+$ queue atypique & Le « trop près » reste fécond\\ \midrule
H7 & Soutenue & Soutenue & Proximité aux nœuds contredite par la disruption ; biais rétrospectif\\ \midrule
H9 & Soutenue (« ont prouvé ») & Limite : cooccurrence $\neq$ relation & Partiellement contredite\\ \midrule
H10 & Soutenue & Dopamine = erreur signée de récompense & Surprise brute contredite (bruit)\\ \midrule
Première expérience & Arbres d'opérateurs des équations d'arXiv, coupure 2014, 400 succès contre 400 échecs, seuil d'aire sous la courbe 0,65 (annexe A, § 7) & SemMedDB et Gromov-Wasserstein fusionné, coupure 2017, contrôle par \emph{permutation des types de relations} $M_{\text{perm}}$ (annexe B, § 7) & SemMedDB, coupure 2015, \emph{deux cibles} : naissance (A) et fécondité conditionnelle (B) (annexe C, § 7)\\
\end{longtable}}

\subsection{Erreurs et fragilités relevées par la vérification}
\begin{itemize}[itemsep=2pt]
\item \textbf{Argon}, Tangent-CFT attribué à « Davila et al. » (marqué « à vérifier ») : Crossref donne Mansouri, Rohatgi, Oard, Wu, Giles et Zanibbi (2019), \emph{ICTIR}, DOI 10.1145/3341981.3344235.
\item \textbf{Neon}, « Titouan, V. (2019) » : le prénom pris pour le nom ; l'auteur est Titouan Vayer (arXiv 1805.09114).
\item \textbf{Carbon}, H8 : Bettencourt et al. (2008) modélisent l'émergence des champs par une dynamique de population de type épidémique ; la transition \emph{topologique} des réseaux de collaboration est l'objet de Bettencourt, Kaiser et Kaur (2009), \emph{Journal of Informetrics}. Attribution à revoir.
\item \textbf{Argon}, « possible adjacent » rattaché à Kauffman (1995) : le terme est surtout développé dans \emph{Investigations} (2000) ; lequel l'introduit, à vérifier.
\item \textbf{Tous}, Merton (1961) : aucune trace par DOI ni par Crossref (la revue est sur JSTOR, refusée au robot) : « à vérifier ». Hofstadter et Sander (2013) : aucun identifiant trouvé : « à vérifier ».
\item \textbf{Surinterprétations} : Galison « a démontré » (Argon, Carbon) ; Tshitoyan et al. « ont prouvé » (Argon) ; « preuve causale directe puisque les équations de 2014 précèdent l'idée » (Argon, § 7.3) — la précédence temporelle exclut la fuite du futur, pas les causes communes.
\item \textbf{Noether} : l'argument « réseau discret, donc pas de Noether » est trop rapide (Marsden et West, 2001 ; Baez et Fong, 2013) ; la bonne raison est l'absence de dynamique définie (partie 2, dernier complément).
\item \textbf{Chiffres non vérifiés} : les trois citent des effectifs (articles du Web of Science, cas d'Ogburn et Thomas, projets de l'Australian Research Council) sans page ; laissés tels quels en annexe, non repris ici.
\item \textbf{Forme} : Argon et Carbon dépassent la limite de 3 500 mots (registre, R-5176).
\end{itemize}

\subsection{Ce qu'aucun des trois n'a vu}
Aucun ne cite les travaux qui occupent le plus directement la prédiction des liens futurs entre concepts : Krenn et Zeilinger (2020), Krenn et al. (2023, Science4Cast), Sourati et Evans (2023), Rzhetsky et al. (2015). C'est important pour deux raisons. D'abord, par la règle contre le faux vide (R8) : la question « prédire quels mariages naîtront » n'est pas une zone vierge. Ensuite, Science4Cast apprend que des traits de réseau simples sont des bases difficiles à battre, et Sourati et Evans que la distribution des experts prédit une grande part des découvertes : une couche de structure qui ne bat pas \emph{ces} bases n'apporte rien. Ce qui reste, à ce stade de lecture, non occupé par ces travaux : la couche de \emph{structure} relationnelle (H3) et la \emph{scission} (H6) — à vérifier par la revue bibliographique de la version 2.

\subsection{Propositions retenues par la note (à confirmer par le directeur)}
\begin{enumerate}[itemsep=3pt]
\item \textbf{Une seule première expérience, qui réunit les meilleurs contrôles des trois.} Le cœur H1$\times$H3 testé par l'interaction $d_{\text{surf}}\times K$ (Carbon, Neon) ; le contrôle par permutation des types de relations à topologie fixée (Carbon) ; les deux cibles naissance et fécondité conditionnelle (Neon) ; une réplique sur les équations (Argon). \emph{Raison} : chaque contrôle ferme une lecture parasite différente — la permutation montre que ce sont les \emph{types} qui comptent et non la seule densité ; les deux cibles séparent la plausibilité rhétorique (ce qui se publie) de la fécondité (ce qui est repris) ; les équations suppriment le bruit de l'extraction de relations.
\item \textbf{Bases obligatoires élargies} : en plus de celles des trois modèles, les traits de réseau de Science4Cast et un modèle « conscient des humains » à la Sourati et Evans. \emph{Raison} : partie 3.4.
\item \textbf{Deux définitions concurrentes du champ}, préenregistrées et comparées sur les multiples : la co-exposition aux instruments en croissance (les trois) et l'accroissement du meilleur alignement sous un cadre neuf, $F^\star$ (section~\ref{sec:champ-distance}, reformulation acceptée par le directeur). \emph{Raison} : la première est facile à mesurer mais peut n'être qu'un indicateur de mode ; la seconde est la thèse du directeur, mais dépend de l'alphabet des relations.
\item \textbf{Élagages} : résonance et ferromagnétisme au sens propre écartés (seul reste le terme bilinéaire) ; Noether gardé comme question ouverte (générateur mesurable), jamais comme affirmation ; transition de phase thermodynamique écartée au profit de la percolation sur réseaux finis. \emph{Raison} : aucune ne produit de prédiction mesurable propre.
\item \textbf{Trois ajouts au modèle} : une couche d'\emph{anomalies} (résidus empiriques), la \emph{croissance de l'alphabet} des abstractions, et l'exigence d'une invariance \emph{causale} (et non seulement de graphe) pour tout transfert vers les modèles de langage, en particulier douleur et anesthésie. \emph{Raison} : ce sont les trois oublis qui, s'ils sont réels, feraient échouer le modèle pour une raison autre que H1 et H3.
\item \textbf{H6 comme pièce propre à creuser} : la scission des fausses identités est presque absente des mesures publiées selon Neon ; à vérifier par la revue bibliographique avant tout autre emploi.
\item \textbf{H10} : pour l'organisme d'idéation, mesurer la surprise bayésienne et le progrès (non l'improbabilité), relativement à un juge fixe ; signal d'exploration, jamais critère de production.
\end{enumerate}
%=====================================================================
\clearpage


\section*{Partie 4 — Annexes : les trois réponses en \LaTeX{}}\addcontentsline{toc}{section}{Partie 4 — Annexes : les trois réponses}
Les trois réponses sont reproduites \emph{sans changement de fond} : conversion automatique du Markdown d'origine par \texttt{pandoc}~3.1.3 (titres, listes, tableaux, formules en mode mathématique), avec deux seules retouches de forme vérifiées par programme (même texte, aux blancs et séparateurs près) : les listes de références de Neon écrites en un paragraphe sont mises en liste, et les titres ne sont pas renumérotés (ils portent leur propre numérotation). \textbf{Les chiffres et références qu'elles contiennent ne sont pas vérifiés ici} ; voir la bibliographie vérifiée et la partie 3. Ce sont des données : rien de ce qu'elles proposent n'a été exécuté.
\clearpage
\section*{Annexe A — Réponse d'Argon}\addcontentsline{toc}{subsection}{Annexe A — Réponse d'Argon}
{\small Source : \texttt{reponses-google-v1/reponse-argon-v1.md}.}\par
\subsection*{1. Confrontation à la littérature}

Les dix hypothèses (H1--H10) mêlent trois problèmes distincts que la
modélisation doit impérativement séparer : 1. \textbf{La dimension
psychologique et sociologique (l'intérêt perçu) :} ce qui capte
l'attention d'un public humain à l'instant \(t\). Une idée est jugée
intéressante lorsqu'elle réfute une croyance faiblement ancrée mais
largement partagée par une audience (Davis, 1971) ou maximise le gain
d'information par rapport aux distributions \emph{a priori} d'un
observateur (surprise bayésienne ; Itti \& Baldi, 2009). Cet intérêt est
relatif, local et périssable. 2. \textbf{La dimension épistémique (la
vérité et la fécondité) :} ce qui fait qu'une idée survit à l'épreuve
empirique ou formelle, résout des anomalies et engendre de nouveaux
programmes de recherche (Lakatos, 1976). Elle dépend de la validité de
l'invariance structurelle découverte et se mesure \emph{ex post} par le
pouvoir d'unification explicative (Kitcher, 1981) et la capacité à
réorganiser le réseau de citations (indice de disruption ; Funk \&
Owen-Smith, 2017). 3. \textbf{La dimension mécanique et algorithmique
(la génération) :} le processus d'exploration d'un espace combinatoire
(créativité combinatoire ou transformationnelle ; Boden, 2004)
permettant à un agent --- humain ou artificiel --- d'identifier ces
appariements avant leur validation.

Une idée peut optimiser (1) tout en échouant sur (2) : c'est le cas des
analogies séduisantes mais stériles. Inversement, une découverte
fondamentale (2) peut violer trop radicalement les croyances du moment
pour être jugée intéressante à sa naissance (résistance paradigmatique).

\subsubsection*{Évaluation des hypothèses}

\begin{itemize}
\tightlist
\item
  \textbf{H1 (Mariage éloigné) : Soutenue.} L'idée que la créativité
  scientifique marie des cadres distants en surface mais isomorphes en
  profondeur correspond à la « bisociation » de Koestler (1964) et à la
  théorie de l'alignement structurel (\emph{Structure-Mapping Theory})
  de Gentner (1983), pour qui une analogie féconde ignore les attributs
  de surface d'un objet pour projeter ses relations d'ordre supérieur
  (voir aussi Hofstadter \& Sander, 2013). En scientométrie, Uzzi,
  Mukherjee, Stringer et Jones (2013) ont démontré sur 17,9 millions
  d'articles du \emph{Web of Science} que les articles à fort impact
  combinent une forte conventionnalité avec une queue de combinaisons
  atypiques (journaux rarement co-cités).
\item
  \textbf{H2 (Cascades et transferts) : Soutenue.} L'inférence
  transitive (\(A\text{--}B\) et \(B\text{--}C \implies A\text{--}C\))
  entre littératures disjointes est le fondement de la découverte fondée
  sur la littérature (\emph{Literature-Based Discovery}), formalisée par
  Swanson (1986) avec le modèle ABC (appliqué historiquement au lien
  entre huile de poisson, viscosité sanguine et syndrome de Raynaud). Le
  transfert directionnel d'un domaine source mature vers un domaine
  cible jeune est documenté en sciences cognitives par Holyoak et
  Thagard (1995).
\item
  \textbf{H3 (Serrure et clé, Noether et lagrangien) : Soutenue pour
  l'alignement de graphes ; métaphore sans contenu testable en l'état
  pour Noether et le lagrangien.} La nécessité d'aligner tout le réseau
  de dépendances en amont correspond au \emph{principe de systématicité}
  de Gentner (1983), implémenté dans le \emph{Structure-Mapping Engine}
  (Falkenhainer, Forbus, \& Gentner, 1989) et formalisé en théorie des
  catégories par les \emph{ologs} (foncteurs préservant les diagrammes
  commutatifs entre ontologies ; Spivak \& Kent, 2012). En revanche,
  l'invocation du théorème de Noether et d'un lagrangien est une
  métaphore non opératoire : le théorème de Noether relie les symétries
  continues (groupes de Lie) d'une fonctionnelle d'action différentiable
  \(\mathcal{S} = \int L(q, \dot{q}, t)\,dt\) à des courants conservés.
  Or, un réseau de concepts est discret, orienté et dépourvu de
  dynamique hamiltonienne ou lagrangienne canonique. L'équivalent
  mathématique rigoureux est l'\textbf{homomorphisme de graphes typés}
  (ou l'équivariance d'opérateurs discrets).
\item
  \textbf{H4 (Ni trop près ni trop loin) : Soutenue.} En psychologie de
  la curiosité, ce phénomène suit la courbe en U inversé de Berlyne
  (1960). En économie de l'innovation, Nooteboom, Van Haverbeke,
  Duysters, Gilsing et Van den Oord (2007) ont établi l'existence d'une
  « distance cognitive optimale » entre partenaires : trop faible, elle
  n'apporte aucune nouveauté ; trop grande, elle dépasse la capacité
  d'absorption des agents. En sociologie des sciences, Foster, Rzhetsky
  et Evans (2015) montrent que les chercheurs arbitrent entre
  consolidation locale (faible risque) et sauts distants (taux d'échec
  élevé, gain élevé en cas de succès).
\item
  \textbf{H5 (Champ d'interaction, air du temps, protéines) : Soutenue
  (sociologie des découvertes multiples) / Métaphore partielle
  (ferromagnétisme et protéines).} L'existence des découvertes
  simultanées a été documentée par Ogburn et Thomas (1922), qui en ont
  recensé 148 cas, puis théorisée par Merton (1961), pour qui les
  découvertes multiples sont la norme dès lors que les prérequis
  culturels et instrumentaux sont réunis. Galison (2003) a démontré
  comment un environnement technologique concret (la synchronisation
  électromagnétique des horloges ferroviaires et les bureaux des
  brevets) a agi comme substrat commun à Einstein et Poincaré.
  L'appariement aléatoire après agitation correspond au modèle de
  variation aveugle et sélection sélective de Simonton (2004).
\item
  \textbf{H6 (Connexion et distinction) : Soutenue.} L'alignement
  structurel ne sert pas qu'à identifier des similitudes ; il fait
  émerger des « différences alignables » (\emph{alignable differences} ;
  Gentner \& Markman, 1997) qui révèlent où l'analogie se brise. En
  philosophie des mathématiques, Lakatos (1976) montre que la collision
  entre une conjecture et un contre-exemple force la scission de
  concepts autrefois confondus (\emph{concept-stretching} et
  \emph{monster-barring}).
\item
  \textbf{H7 (Profondeur de l'idée fille) : Soutenue.} En philosophie
  des sciences, Kitcher (1981) définit le pouvoir explicatif d'une
  théorie par sa capacité à dériver un maximum de conclusions à partir
  d'un nombre minimal de schémas d'argumentation (unification). En
  scientométrie, la capacité d'un nœud nouveau à court-circuiter les
  chemins existants et à éclipser ses parents est capturée par l'indice
  de disruption \(CD\) (Funk \& Owen-Smith, 2017 ; Wu, Wang, \& Evans,
  2019).
\item
  \textbf{H8 (Émergence et percolation) : Soutenue.} La dynamique des
  découvertes suit l'expansion dans le « possible adjacent » (Kauffman,
  1995), formalisée par les processus d'urnes de Pólya déclenchées par
  Tria, Loreto, Servedio et Strogatz (2014). Le rôle de la « surface de
  contact » entre communautés isolées est démontré par Burt (2004) : les
  acteurs situés aux « trous structuraux » (\emph{structural holes})
  entre réseaux denses génèrent des idées mieux évaluées.
\item
  \textbf{H9 (Structure cachée du langage) : Soutenue.} Mikolov,
  Sutskever, Chen, Corrado et Dean (2013) ont montré que les plongements
  lexicaux non supervisés transforment des relations sémantiques
  abstraites en translations vectorielles linéaires. Tshitoyan et
  al.~(2019) ont prouvé que des plongements entraînés sans supervision
  sur des résumés de science des matériaux capturent la structure du
  tableau périodique et prédisent des matériaux thermoélectriques des
  années avant leur découverte expérimentale.
\item
  \textbf{H10 (Signal de surprise) : Soutenue.} En neurosciences,
  Schultz, Dayan et Montague (1997) ont établi que les neurones
  dopaminergiques codent une erreur de prédiction de récompense. En
  intelligence artificielle, Schmidhuber (2010) a formalisé la curiosité
  intrinsèque par la dérivée temporelle de la compression des données
  (gain de compression), et Itti et Baldi (2009) ont défini la surprise
  bayésienne par la divergence de Kullback-Leibler entre la distribution
  \emph{a priori} et la distribution \emph{a posteriori} d'un agent.
\end{itemize}

\begin{center}\rule{0.5\linewidth}{0.5pt}\end{center}

\subsection*{2. Modèle formel minimal}

\subsubsection*{2.1. Espace des concepts : graphe multiplexe daté}

Soit \(\mathcal{G}_t = (V_t, \mathcal{E}_t)\) un graphe multiplexe
orienté et pondéré à l'instant discret \(t \in \mathbb{N}\), où \(V_t\)
est l'ensemble des concepts scientifiques attestés jusqu'à \(t\). Le
graphe comporte trois couches
\(\mathcal{E}_t = \{E^{(U)}_t, E^{(S)}_t, E^{(C)}_t\}\) : 1.
\textbf{Couche d'usage et de surface (\(U\)) :} matrice d'adjacence
symétrique \(W^{(U)}_t \in \mathbb{R}_{\ge 0}^{|V_t|\times |V_t|}\)
mesurant la co-occurrence lexicale et bibliométrique des concepts,
associée à un plongement contextuel de surface
\(x_u(t) \in \mathbb{R}^d\) pour chaque \(u \in V_t\). 2. \textbf{Couche
de structure (\(S\)) :} multigraphe orienté et typé
\(G^{(S)}_t = (V_t, E^{(S)}_t, \mathcal{R})\), où chaque arête
\(e = (a, r, b) \in E^{(S)}_t\) porte un type relationnel formel
\(r \in \mathcal{R}\) (inclusion, causalité, dépendance axiomatique,
opérateur mathématique tel que dérivée, proportionnalité ou
conservation). On note \(\mathcal{N}_S(u, t)\) le sous-graphe
relationnel en amont de \(u\) dans un rayon \(k\) (son « cône de
dépendances »). 3. \textbf{Couche des communautés humaines (\(C\)) :} à
chaque concept \(u\) est associée une distribution empirique de
probabilité \(p_u(t) \in \Delta^{M-1}\) sur les \(M\) communautés
scientifiques (définies par les clusters de co-auteurs et de revues).

\subsubsection*{2.2. Les deux distances (surface et structure)}

\begin{itemize}
\item
  \textbf{Distance de surface \(d_U(u, v; t) \in [0, 1]\) :} combine la
  distance cosinus des représentations discursives et la divergence de
  Jensen-Shannon (\(D_{\text{JS}}\)) entre les communautés humaines qui
  manipulent \(u\) et \(v\) :
  \[d_U(u, v; t) = \frac{\alpha}{2}\left(1 - \frac{x_u(t) \cdot x_v(t)}{\|x_u(t)\|_2 \|x_v(t)\|_2}\right) + (1 - \alpha) D_{\text{JS}}\big(p_u(t) \parallel p_v(t)\big), \quad \alpha \in [0,1]\]
\item
  \textbf{Distance de structure \(d_S(u, v; t) \ge 0\) :} doit être
  invariante par renommage des nœuds (indépendante du jargon des
  disciplines). On dépouille les sous-graphes \(\mathcal{N}_S(u, t)\) et
  \(\mathcal{N}_S(v, t)\) de leurs étiquettes lexicales pour ne
  conserver que leur topologie et les types relationnels
  \(r \in \mathcal{R}\). Soient \(A^{(u)}\) et \(A^{(v)}\) les tenseurs
  d'adjacence relationnels de dimensions respectives
  \(n_u \times n_u \times |\mathcal{R}|\) et
  \(n_v \times n_v \times |\mathcal{R}|\). La distance de structure est
  donnée par la \textbf{distance de Gromov-Wasserstein} sur graphes
  typés :
  \[d_S(u, v; t) = \min_{\Pi \in \mathcal{U}(\mu_u, \mu_v)} \sum_{r \in \mathcal{R}} \sum_{i,j=1}^{n_u} \sum_{k,l=1}^{n_v} \left| A^{(u)}_{ij,r} - A^{(v)}_{kl,r} \right|^2 \Pi_{ik} \Pi_{jl}\]
  où \(\mathcal{U}(\mu_u, \mu_v)\) est le polytope de transport des
  matrices de couplage entre les mesures uniformes \(\mu_u\) et
  \(\mu_v\) sur les nœuds de \(\mathcal{N}_S(u, t)\) et
  \(\mathcal{N}_S(v, t)\).
\end{itemize}

\subsubsection*{2.3. Critère serrure-clé et tri des invariances (H3, H6)}

Soit
\(\phi : V(\mathcal{N}_S(u,t)) \rightharpoonup V(\mathcal{N}_S(v,t))\)
une injection partielle telle que \(\phi(u) = v\). On définit trois
grandeurs mesurables : 1. \textbf{La couverture :}
\(\kappa(\phi) = \frac{|\text{dom}(\phi)|}{\min(n_u, n_v)} \in [0, 1]\),
proportion du réseau parent effectivement appariée. 2. \textbf{La
fidélité d'homomorphisme :}
\(\sigma(\phi) = \frac{1}{|E_{\text{dom}}(\phi)|} \sum_{(a,r,b) \in E_{\text{dom}}(\phi)} \mathbf{1}\big\{(\phi(a), r, \phi(b)) \in E^{(S)}_t\big\}\),
où \(E_{\text{dom}}(\phi)\) est l'ensemble des arêtes de
\(\mathcal{N}_S(u,t)\) dont les deux extrémités appartiennent à
\(\text{dom}(\phi)\). 3. \textbf{Les conflits axiomatiques (H6) :} soit
\(\mathcal{I} \subset \mathcal{R} \times \mathcal{R}\) l'ensemble des
paires de relations mutuellement incompatibles (par exemple : causalité
directe vs inverse, monotonie croissante vs décroissante, symétrie vs
antisymétrie). Le nombre de conflits est :
\[\chi(\phi) = \sum_{(a,r,b) \in E_{\text{dom}}(\phi)} \sum_{r' : (r,r') \in \mathcal{I}} \mathbf{1}\big\{(\phi(a), r', \phi(b)) \in E^{(S)}_t\big\}\]

Le \textbf{score serrure-clé} \(\Phi(u,v;t) \in [0,1]\) est l'optimum
sous contrainte de pénalisation des conflits :
\[\Phi(u, v; t) = \max_{\phi : \phi(u)=v} \Big[ \kappa(\phi) \cdot \sigma(\phi) \cdot \exp\big(-\gamma \chi(\phi)\big) \Big], \quad \gamma > 0\]
Un couple présentant une faible distance de surface \(d_U\) mais un
conflit élevé \(\chi(\phi) \gg 0\) est identifié comme une
\textbf{fausse identité} que la collision des réseaux dissout (H6).

\subsubsection*{2.4. Le champ dépendant du temps (H5)}

Soit \(\mathcal{M}_t \subset V_t\) le sous-ensemble des nœuds identifiés
comme \emph{catalyseurs} (nouveaux instruments de mesure, formalismes
mathématiques, technologies expérimentales). Soit
\(a_m(t) = \frac{d}{dt}\ln \deg_U(m, t)\) le taux de croissance de
l'adoption du catalyseur \(m\), et \(h_m(u, t) \in [0,1]\) l'intensité
du couplage entre le concept \(u\) et le catalyseur \(m\) à l'instant
\(t\) (probabilité conditionnelle \(P(m \mid u; t)\)). Le champ
d'interaction induit sur la paire \((u,v)\) par l'« air du temps » est
la forme bilinéaire :
\[H_t(u, v) = \sum_{m \in \mathcal{M}_t} a_m(t)\, h_m(u, t)\, h_m(v, t)\]
Deux concepts très éloignés (\(d_U \to 1\)) subissent un rapprochement
projectif si tous deux commencent simultanément à être reformulés à
travers le même instrument ou formalisme \(m\).

\subsubsection*{2.5. Probabilité de naissance d'une idée (H1, H2, H4,
H5)}

Pour toute paire non encore reliée dans \(E^{(U)}_t\), le taux de hasard
(intensité de processus ponctuel) \(\lambda_{uv}(t)\) de voir naître
l'idée combinée \((u,v)\) dans l'intervalle \([t, t+dt]\) s'écrit :
\[\resizebox{\linewidth}{!}{$\displaystyle \lambda_{uv}(t) = \lambda_0 \, \underbrace{d_U(u,v;t)^{\beta_1} \exp\big(-\beta_2 d_U(u,v;t)\big)}_{\text{H4 : fenêtre en U inversé (surface)}} \cdot \underbrace{\exp\big(-\beta_3 d_S(u,v;t)\big)}_{\text{H1 : proximité cachée}} \cdot \underbrace{\Phi(u,v;t)^{\beta_4}}_{\text{H3 : serrure-clé}} \cdot \underbrace{\exp\big(\beta_5 C_{uv}(t) + \beta_6 H_t(u,v)\big)}_{\text{H2 : cascades + H5 : champ}}$}\]
où \(\beta_1, \dots, \beta_6 \ge 0\) sont les paramètres à estimer et
\(C_{uv}(t) = \sum_{w \in V_t} \tilde{W}^{(U)}_{uw}(t)\tilde{W}^{(U)}_{wv}(t)\)
mesure le flux des cascades transitives à deux sauts
\(u\text{--}w\text{--}v\) (modèle ABC de Swanson, où \(\tilde{W}^{(U)}\)
est la matrice d'adjacence normalisée). Le terme de surface atteint son
maximum en \(d_U^* = \beta_1 / \beta_2\) (H4).

\subsubsection*{2.6. Intérêt subjectif et profondeur épistémique de
l'idée fille (H7, H10)}

Conformément à la distinction posée en section 1, on sépare deux
métriques pour l'idée fille \(z = u \otimes v\) créée à \(t\) : 1.
\textbf{Le signal de surprise immédiate (H1, H10) :} divergence
d'information entre la probabilité structurelle de l'appariement et sa
probabilité de surface \emph{a priori}
\(P_U(u,v;t) \propto \exp(-\eta d_U(u,v;t))\) :
\[S_{\text{surprise}}(z; t) = -\ln P_U(u, v; t) \cdot \Phi(u, v; t) \propto d_U(u, v; t) \cdot \Phi(u, v; t)\]
2. \textbf{La profondeur et le pouvoir d'unification
\(\mathcal{D}(z; t, \tau)\) à l'horizon \(t+\tau\) (H7) :} réduction
pondérée de la longueur des plus courts chemins géodésiques
\(\ell_S(i,j)\) sur la couche structurelle \(G^{(S)}\) entre les
concepts existants, pondérée par leur centralité de vecteur propre
\(c_i(t)\) :
\[\mathcal{D}(z; t, \tau) = \sum_{i, j \in V_t} c_i(t) c_j(t) \Big[ \ell_{G^{(S)}_t}(i, j) - \ell_{G^{(S)}_{t+\tau}}(i, j) \Big]\]
Une idée est profonde si son insertion contracte globalement les
distances entre les nœuds centraux du graphe (unification de la Lune et
de la pomme).

\subsubsection*{2.7. Dynamique d'émergence et percolation (H8)}

Lorsqu'un lien \((u,v)\) est validé entre deux communautés \(A\) et
\(B\), l'application optimale \(\phi^*\) expose
\(|\text{dom}(\phi^*)| - 1\) paires de concepts satellites
\((a, \phi^*(a)) \in A \times B\) comme candidats immédiats. Si chaque
satellite exposé s'apparie à son tour avec une probabilité indépendante
\(p_{\text{trans}} \propto \Phi(u,v;t)\), le nombre de liens
inter-communautés suit un processus de branchement de Galton-Watson
allié à une \textbf{percolation de liens} sur le graphe biparti
\(A \times B\). Une transition de phase (émergence d'une composante
géante hybride, c'est-à-dire d'une nouvelle sous-discipline) survient
lorsque la densité de surface de contact et la qualité de
l'homomorphisme franchissent le seuil critique :
\[R_0 = p_{\text{trans}} \cdot \big(\langle |\text{dom}(\phi^*)| \rangle - 1\big) > 1\]

\subsubsection*{2.8. Traitement explicite des analogies physiques}

\begin{longtable}[]{@{}
  >{\raggedright\arraybackslash}p{(\columnwidth - 4\tabcolsep) * \real{0.3333}}
  >{\raggedright\arraybackslash}p{(\columnwidth - 4\tabcolsep) * \real{0.3333}}
  >{\raggedright\arraybackslash}p{(\columnwidth - 4\tabcolsep) * \real{0.3333}}@{}}
\toprule\noalign{}
\begin{minipage}[b]{\linewidth}\raggedright
Analogie physique
\end{minipage} & \begin{minipage}[b]{\linewidth}\raggedright
Statut dans le modèle
\end{minipage} & \begin{minipage}[b]{\linewidth}\raggedright
Équivalent mathématique exact ou motif d'exclusion
\end{minipage} \\
\midrule\noalign{}
\endhead
\bottomrule\noalign{}
\endlastfoot
\textbf{Puits de potentiel} & \textbf{Retenue} & Fonction d'énergie
libre d'appariement
\(E(u,v;t) = \beta_3 d_S(u,v;t) - \beta_4 \ln \Phi(u,v;t)\). Les minima
locaux correspondent aux isomorphismes de sous-graphes sans conflit
axiomatique (\(\chi = 0\)). \\
\textbf{Champ et polarisation (ferromagnétisme)} & \textbf{Retenue} &
Terme de couplage externe \(H_t(u,v) = \sum_m a_m(t) h_m(u,t) h_m(v,t)\)
de type Hamiltonien d'Ising avec champ local inhomogène, où l'alignement
des spins correspond à l'adoption d'un même formalisme ou instrument
\(m\). \\
\textbf{Appariement de protéines} & \textbf{Retenue} & Complémentarité
géométrique = score d'homomorphisme \(\kappa(\phi)\sigma(\phi)\) ;
encombrement stérique = pénalité exponentielle des conflits axiomatiques
\(\exp(-\gamma\chi(\phi))\). \\
\textbf{Percolation} & \textbf{Retenue} & Percolation de liens sur
graphe biparti inter-communautés avec seuil critique \(R_0 = 1\)
gouverné par la taille moyenne du domaine aligné
\(\langle|\text{dom}(\phi^*)|\rangle\). \\
\textbf{Théorème de Noether et lagrangien} & \textbf{Écartée} & Aucun
espace des phases continu ni intégrale d'action \(\int L\,dt\) n'existe
sur un réseau de concepts discrets. Remplacée par l'invariance
relationnelle discrète sous injection partielle :
\((a,r,b) \in E^{(S)} \implies (\phi(a),r,\phi(b)) \in E^{(S)}\). \\
\textbf{Résonance} & \textbf{Écartée} & Présuppose une équation
d'oscillateur harmonique avec fréquence propre \(\omega_0\) et phase
continue, grandeurs non mesurables sur des concepts sans décisions
arbitraires. Remplacée par \(\Phi(u,v;t)\). \\
\end{longtable}

\begin{center}\rule{0.5\linewidth}{0.5pt}\end{center}

\subsection*{3. Données brutes et extraction des variables}

Pour éviter toute circularité, toutes les variables prédictives à la
date \(t\) doivent être calculées exclusivement sur des archives
antérieures ou égales à \(t\).

\begin{enumerate}
\def\labelenumi{\arabic{enumi}.}
\tightlist
\item
  \textbf{Corpus scientifiques datés en texte intégral et graphes
  bibliométriques :}

  \begin{itemize}
  \tightlist
  \item
    \emph{OpenAlex} (plus de 250 millions de travaux datés, graphes de
    citations, affiliations, concepts hiérarchisés) et \emph{Semantic
    Scholar Open Research Corpus} (S2ORC ; Lo et al., 2020) pour le
    texte intégral.
  \item
    \emph{Sources LaTeX d'arXiv} (physique, mathématiques, informatique,
    biologie quantitative depuis 1991) et \emph{PubMed Central Open
    Access Subset} (biomédecine).
  \end{itemize}
\item
  \textbf{Ontologies et vocabulaires contrôlés :}

  \begin{itemize}
  \tightlist
  \item
    \emph{MeSH / UMLS} (biomédecine), \emph{Physics and Astronomy
    Classification Scheme} (PACS), \emph{Mathematics Subject
    Classification} (MSC), \emph{Computer Science Ontology} (CSO). Ces
    ontologies définissent l'ensemble des nœuds \(V_t\) et servent à
    filtrer le bruit lexical.
  \end{itemize}
\item
  \textbf{Corpus d'équations (pour la couche structurelle \(S\)) :}

  \begin{itemize}
  \tightlist
  \item
    Extraction de toutes les équations des sources LaTeX d'arXiv. Une
    équation est une réalisation pure de la couche \(S\) : une fois les
    symboles de variables anonymisés, son arbre syntaxique d'opérateurs
    décrit une structure mathématique exacte indépendante du domaine
    d'application.
  \end{itemize}
\item
  \textbf{Catalogues de découvertes simultanées et de jalons :}

  \begin{itemize}
  \tightlist
  \item
    Catalogue historique d'Ogburn et Thomas (1922) ; jeux de données des
    articles associés aux prix Nobel (Li, Yin, Fortunato, \& Wang, 2019)
    ; paires de prépublications indépendantes soumises sur arXiv dans
    une fenêtre \(\Delta t < 30\) jours sans auteurs ni références
    récentes en commun et co-citées ultérieurement comme équivalentes.
  \end{itemize}
\item
  \textbf{Catalogues d'analogies et de mariages ayant échoué (contre le
  biais du survivant) :}

  \begin{itemize}
  \tightlist
  \item
    \emph{Mariages mort-nés dans la littérature :} paires de concepts
    \((u,v)\) issues de disciplines distinctes qui co-apparaissent pour
    la première fois dans un article à l'année \(t\), mais ne génèrent
    aucune co-occurrence supplémentaire par d'autres équipes au cours
    des dix années suivantes \([t+1, t+10]\).
  \item
    \emph{Essais translationnels échoués :} base
    \emph{ClinicalTrials.gov} croisée avec \emph{DrugBank}
    (repositionnements de molécules d'une pathologie à une autre ayant
    échoué en phase II ou III pour manque d'efficacité).
  \item
    \emph{Base Retraction Watch} (sous-ensemble des rétractations pour
    erreur méthodologique ou invalidation expérimentale, hors fraude
    délibérée).
  \end{itemize}
\end{enumerate}

\subsubsection*{\texorpdfstring{Table de correspondance Données \(\to\)
Variables}{Table de correspondance Données \textbackslash to Variables}}

\begin{itemize}
\tightlist
\item
  \textbf{\(V_t\) :} entités normalisées des ontologies présentes dans
  le corpus à la date \(\le t\).
\item
  \textbf{\(d_U(u,v;t)\) :} \(x_u(t)\) est la moyenne des vecteurs
  contextuels du concept \(u\) à l'année \(t\) ; \(p_u(t)\) est
  l'histogramme des publications mentionnant \(u\) réparties sur les
  clusters de revues de la couche \(C\) (détectés par l'algorithme de
  Leiden sur le graphe de co-citation à l'instant \(t\)).
\item
  \textbf{\(G^{(S)}_t, d_S(u,v;t)\) et \(\Phi(u,v;t)\) :} construits par
  l'union de deux graphes typés : (a) les triplets sémantiques
  \texttt{(Sujet,\ Prédicat\_Typé,\ Objet)} extraits du texte (par ex.
  \texttt{INHIBITS}, \texttt{INCREASES}, \texttt{IS\_PROPORTIONAL\_TO}),
  et (b) les arbres d'opérateurs mathématiques associés au concept
  \(u\).
\item
  \textbf{\(H_t(u,v)\) :} l'ensemble des catalyseurs \(\mathcal{M}_t\)
  est constitué des nœuds appartenant aux branches « Méthodes,
  Instruments et Équipements » des ontologies (ex. branche \emph{E05
  Investigative Techniques} de MeSH). \(h_m(u,t)\) est l'information
  mutuelle ponctuelle positive (PPMI) entre \(u\) et \(m\) sur la
  fenêtre \([t-3, t]\).
\end{itemize}

\begin{center}\rule{0.5\linewidth}{0.5pt}\end{center}

\subsection*{4. Méthodes et outils}

\begin{enumerate}
\def\labelenumi{\arabic{enumi}.}
\tightlist
\item
  \textbf{Construction de la couche de surface (\(U\)) et alignement
  diachronique :}

  \begin{itemize}
  \tightlist
  \item
    Plongements lexicaux et documentaires entraînés par tranches
    annuelles strictes à l'aide de \textbf{SPECTER2} ou \textbf{SciBERT}
    (Beltagy, Lo, \& Cohan, 2019), ou de modèles \emph{Skip-Gram with
    Negative Sampling} (SGNS) alignés d'une année sur l'autre par
    \textbf{analyse procustéenne orthogonale} (Hamilton, Leskovec, \&
    Jurafsky, 2016) afin d'éviter toute fuite du futur vers le passé.
  \item
    Partitionnement de la couche communautaire \(C\) par
    l'\textbf{algorithme de Leiden} (Traag, Waltman, \& van Eck, 2019)
    sur le réseau bipartite auteurs-revues.
  \end{itemize}
\item
  \textbf{Construction de la couche structurelle (\(S\)) :}

  \begin{itemize}
  \tightlist
  \item
    \emph{Sur le texte :} extraction d'entités et de relations typées
    par \textbf{DyGIE++} (Wadden, Wennberg, Luan, \& Hajishirzi, 2019)
    entraîné sur le corpus \emph{SciERC}, et \textbf{SemRep} (Rindflesch
    \& Fiszman, 2003) pour les prédicats biomédicaux UMLS.
  \item
    \emph{Sur les équations :} conversion des sources LaTeX en
    \emph{Content MathML} via \textbf{LaTeXML}, puis extraction des
    arbres d'opérateurs (\emph{Operator Trees}) via \textbf{Tangent-CFT}
    (Davila et al., 2019 \emph{(à vérifier)}). Toutes les feuilles
    (variables locales \(x, y, \alpha\)) sont remplacées par des jetons
    génériques typés (\texttt{SCALAR}, \texttt{VECTOR}, \texttt{FIELD}),
    ne laissant que le squelette des opérateurs
    (\(\nabla^2, \frac{\partial}{\partial t}, \int, =\)).
  \end{itemize}
\item
  \textbf{Calcul de \(d_S\) et du score serrure-clé \(\Phi\) :}

  \begin{itemize}
  \tightlist
  \item
    \(d_S(u,v;t)\) est calculée par l'algorithme de \textbf{Fused
    Gromov-Wasserstein} (Vayer, Chapel, Flamary, Tavenard, \& Courty,
    2019) implémenté dans la bibliothèque \texttt{POT} (\emph{Python
    Optimal Transport}), avec régularisation entropique de Sinkhorn pour
    passer à l'échelle.
  \item
    \(\Phi(u,v;t)\) (recherche de l'injection partielle maximale
    \(\phi^*\) avec pénalité de conflit \(\chi\)) est résolu exactement
    sur les voisinages de rayon \(k \le 3\) par l'algorithme
    \textbf{VF2++} (Jüttner \& Madarasi, 2018) étendu au \emph{Maximum
    Common Edge Subgraph} (MCES), ou encodé comme programme logique en
    \textbf{Answer Set Programming (solveur Clingo)} où les
    incompatibilités \(\mathcal{I}\) sont déclarées comme contraintes
    d'intégrité souples (\emph{weak constraints}).
  \end{itemize}
\item
  \textbf{Détection du champ \(H_t\) :}

  \begin{itemize}
  \tightlist
  \item
    Identification des pics d'émergence méthodologique \(a_m(t)\) par
    l'\textbf{algorithme de détection de rafales (\emph{burst
    detection}) de Kleinberg} (Kleinberg, 2003) à automate à états finis
    sur les nœuds de \(\mathcal{M}_t\).
  \end{itemize}
\item
  \textbf{Estimation statistique et tests :}

  \begin{itemize}
  \tightlist
  \item
    Estimation de \(\lambda_{uv}(t)\) par un \textbf{modèle à risques
    proportionnels de Cox} avec covariables dépendant du temps et
    échantillonnage cas-témoins imbriqué (\emph{nested case-control} ;
    Borgan, Goldstein, \& Langholz, 1995) pour traiter la taille
    quadratique de \(|V_t \times V_t|\).
  \item
    Mesure de la profondeur \(\mathcal{D}(z)\) et de la fécondité par
    l'indice de disruption \textbf{\(CD_5\)} (Funk \& Owen-Smith, 2017)
    et simulation de la transition de percolation inter-communautés par
    l'\textbf{algorithme de Newman-Ziff} (Newman \& Ziff, 2000).
  \end{itemize}
\end{enumerate}

\begin{center}\rule{0.5\linewidth}{0.5pt}\end{center}

\subsection*{5. Validation hors de l'auteur du modèle}

\subsubsection*{\texorpdfstring{5.1. Prédiction rétrospective
hors-échantillon (\emph{Out-of-time Rolling
Forecast})}{5.1. Prédiction rétrospective hors-échantillon (Out-of-time Rolling Forecast)}}

Le corpus est tronqué à une date charnière \(t_0 = 2010\). Aucun texte,
plongement, graphe ou paramètre postérieur à 2010 n'est accessible lors
de l'entraînement. Le modèle prédit, parmi les paires \((u,v)\) jamais
co-mentionnées jusqu'en 2010 : 1. Leur probabilité d'apparition entre
2011 et 2015 (\(\lambda_{uv}(2010)\)). 2. Conditionnellement à leur
apparition, leur fécondité mesurée en 2025 (\(\mathcal{D}(z; 2015, 10)\)
et indice \(CD_5\)).

\subsubsection*{\texorpdfstring{5.2. Modèles nuls et \emph{baselines} de
référence}{5.2. Modèles nuls et baselines de référence}}

Le modèle complet doit être confronté à quatre modèles de base sur les
mêmes données : 1. \textbf{Attachement préférentiel pur
(Barabási-Albert) :}
\(\lambda_{uv}^{\text{PA}}(t) \propto \deg_U(u,t) \cdot \deg_U(v,t)\).
2. \textbf{Surface et topologie locale seule :} régression logistique
sur \(d_U(u,v;t)\) et l'indice d'Adamic-Adar sur le graphe de
co-occurrence \(W^{(U)}_t\). 3. \textbf{Swanson ABC sans contrôle
structurel :} modèle n'utilisant que la fermeture triadique
\(C_{uv}(t)\) et \(d_U(u,v;t)\). 4. \textbf{Structure seule sans filtre
de surface :} modèle n'utilisant que \(d_S(u,v;t)\) et \(\Phi(u,v;t)\).

\subsubsection*{5.3. Traitement du biais du survivant}

On ne doit pas seulement comparer les découvertes célèbres à des paires
aléatoires non connectées (qui sont trivialement absentes). L'évaluation
porte sur la discrimination entre \textbf{trois cohortes appariées} par
degré \(\deg_U\) et distance de surface \(d_U(u,v;2010)\) : *
\textbf{Cohorte 1 (Mariages féconds) :} paires apparues en 2011--2015
figurant dans le décile supérieur des citations et présentant un indice
de disruption \(CD_5 > 0\) en 2025. * \textbf{Cohorte 2 (Mariages
mort-nés / stériles) :} paires effectivement publiées entre 2011 et 2015
(donc ayant franchi la barrière de l'intérêt subjectif des auteurs et
des pairs à l'instant \(t\)), mais n'ayant reçu aucune co-citation de
suivi par des équipes indépendantes entre 2016 et 2025, ou correspondant
à des essais cliniques de repositionnement échoués. * \textbf{Cohorte 3
(Abstentions structurelles) :} paires non connectées entre 2011 et 2025.

Le modèle n'est validé que si \(\Phi(u,v;2010)\) et \(\chi(\phi^*)\)
séparent nettement la Cohorte 1 de la Cohorte 2 (distinction entre
surprise superficielle et viabilité structurelle).

\subsubsection*{5.4. Test sur les découvertes simultanées (H5)}

Sur le catalogue des découvertes multiples (paires d'articles
indépendants réalisant le même mariage \((u,v)\) dans une fenêtre de 30
jours), on apparie chaque découverte multiple à une découverte «
singleton » de même discipline, même année et même impact final. On
teste si le terme de champ \(H_t(u,v)\) présente une élévation
statistiquement significative dans les 3 années précédant la découverte
pour les multiples par rapport aux singletons (test de Wilcoxon apparié
sur \(\Delta H_t\)).

\subsubsection*{5.5. Évaluation humaine en aveugle et préenregistrement}

\begin{itemize}
\tightlist
\item
  \textbf{Protocole en double aveugle :} le modèle figé à l'année
  courante \(T\) génère 120 propositions d'appariements inédits dans un
  domaine cible (par ex. neurosciences computationnelles ou science des
  matériaux), réparties en quatre bras équilibrés : (1) Haute surprise,
  Haut score \(\Phi\) ; (2) Haute surprise, Faible \(\Phi\) (fausse
  analogie) ; (3) Faible surprise, Haut \(\Phi\) ; (4) Bruit apparié par
  degré. Les propositions sont résumées sous une forme standardisée par
  un gabarit neutre indiquant les concepts parents et le sous-graphe
  aligné \(\phi^*\), sans mentionner le bras d'origine. Un panel
  d'experts du domaine évalue chaque proposition sur trois échelles de
  Likert : \emph{Surprise}, \emph{Plausibilité formelle/mécanistique},
  et \emph{Intention d'investir du temps de recherche}.
\item
  \textbf{Préenregistrement :} l'architecture, la liste des relations
  incompatibles \(\mathcal{I}\), les hyperparamètres
  (\(\alpha, \gamma, k\)) et les tests statistiques sont déposés et
  horodatés sur l'\emph{Open Science Framework} (OSF) avant l'extraction
  de la période de test 2011--2025.
\end{itemize}

\subsubsection*{5.6. Critères explicites de réfutation}

Le modèle (et les hypothèses H1--H5 qu'il formalise) sera considéré
comme \textbf{réfuté} si l'un des résultats suivants est observé sur
l'échantillon de test hors-période : 1. L'ajout des variables
structurelles (\(d_S\) et \(\Phi\)) au modèle de base (Attachement
préférentiel + \(d_U\) + Swanson \(C_{uv}\)) n'apporte aucun gain
statistiquement significatif d'aire sous la courbe Précision-Rappel
(\(\Delta \text{AUC-PR} \le 0,02\) ou \(p > 0,01\) par test de
permutation) pour prédire les mariages féconds. 2. Le score serrure-clé
\(\Phi(u,v;t_0)\) calculé \emph{avant} la publication ne discrimine pas
les mariages féconds (Cohorte 1) des mariages mort-nés (Cohorte 2)
(\(\text{AUC-ROC} < 0,55\)). Cela prouverait que l'isomorphisme
structurel n'est pas la cause de la découverte, mais une rationalisation
\emph{ex post} construite par les chercheurs après le mariage. 3. Les
coefficients estimés \(\beta_1\) ou \(\beta_2\) sont nuls ou négatifs,
réfutant la forme en U inversé de H4 au profit d'une décroissance
monotone pure avec la distance.

\begin{center}\rule{0.5\linewidth}{0.5pt}\end{center}

\subsection*{6. Ce que les hypothèses oublient}

Un modèle complet de la naissance et du succès des idées scientifiques
doit intégrer quatre mécanismes absents de H1--H10 :

\begin{enumerate}
\def\labelenumi{\arabic{enumi}.}
\tightlist
\item
  \textbf{La friction du réel et la rétroaction expérimentale (le filtre
  externe) :} H1--H10 traitent la science comme un jeu purement interne
  de réorganisation de graphes conceptuels. Or, deux réseaux de concepts
  peuvent être mathématiquement isomorphes (\(\Phi = 1, \chi = 0\)) tout
  en étant empiriquement faux lorsque leur prédiction est confrontée à
  la nature (par exemple, la théorie des tourbillons de Descartes ou les
  premières analogies hydrodynamiques de l'éther). Il faut introduire un
  opérateur de \textbf{choc empirique} \(E_{\text{obs}}(t)\) : les
  anomalies expérimentales (Kuhn, 1962) injectent des contraintes de
  conflit forcées dans la couche \(S\) qui invalident des régions
  entières du graphe indépendamment de leur élégance formelle.
\item
  \textbf{L'asymétrie directionnelle du transfert (Donneur vs Récepteur)
  :} Le terme de mariage suggère une symétrie entre les parents
  \((u,v)\). En réalité, le transfert analogique est presque toujours
  orienté d'un domaine \emph{source} à forte densité axiomatique
  (mathématisé, doté d'invariants explicites) vers un domaine
  \emph{cible} riche en données empiriques mais faiblement formalisé
  (Gentner, 1983). Dans l'exemple de H2 (douleur \(\to\) anesthésie
  \(\to\) conscience \(\to\) motifs internes d'un modèle de langage), la
  direction de l'inférence n'est pas symétrique. Le modèle doit donc
  remplacer la mesure symétrique par un score orienté
  \(\Phi(u \to v; t)\) pondéré par le différentiel de profondeur
  axiomatique
  \(\Delta \text{Profondeur}(u,v) = \text{profondeur}(\mathcal{N}_S(u)) - \text{profondeur}(\mathcal{N}_S(v))\).
\item
  \textbf{Les coûts institutionnels et le risque de carrière (la
  pénalité interdisciplinaire) :} Même lorsqu'une idée maximise le
  potentiel épistémique, sa probabilité de publication et de financement
  est freinée par l'organisation sociale des disciplines. Bromham,
  Dinnage et Hua (2016) ont démontré sur 18 476 projets soumis à
  l'\emph{Australian Research Council} que plus la distance
  interdisciplinaire d'un projet est élevée, plus sa probabilité de
  financement diminue. Il faut ajouter à l'équation de naissance
  \(\lambda_{uv}(t)\) un terme de \textbf{friction institutionnelle}
  \(F_{\text{inst}}(p_u, p_v)\) mesuré par l'absence de revues ou de
  comités d'évaluation chevauchant les communautés de \(u\) et de \(v\).
\item
  \textbf{La topologie des agents porteurs (courtiers cognitifs et
  polyglottes) :} Les concepts ne s'apparient pas seuls ; la probabilité
  de collision dépend de l'existence d'agents humains capables de lire
  simultanément les deux sous-graphes \(\mathcal{N}_S(u)\) et
  \(\mathcal{N}_S(v)\). La naissance d'une idée à fort \(d_U\) dépend de
  la densité de chercheurs migrants (ayant changé de discipline entre
  leur doctorat et leur poste actuel) ou d'équipes hybrides positionnées
  sur les trous structuraux de la couche \(C\) (Burt, 2004). Sans
  porteur bilingue, un puits de potentiel structurel reste inexploré.
\end{enumerate}

\begin{center}\rule{0.5\linewidth}{0.5pt}\end{center}

\subsection*{7. Une première expérience peu coûteuse (réfutation du cœur
H1 + H3)}

Cette expérience vise à mettre à l'épreuve le cœur du modèle (H1 :
distance de surface vs proximité structurelle ; H3 : alignement
serrure-clé des réseaux parents) en \textbf{4 jours de travail} avec un
budget informatique inférieur à \textbf{100 €} (une instance cloud
CPU/GPU standard), en éliminant toute ambiguïté liée à l'extraction de
relations en langage naturel grâce à l'usage exclusif des \textbf{arbres
d'opérateurs d'équations mathématiques}.

\subsubsection*{7.1. Pourquoi les équations d'arXiv ?}

Extraire la structure logique d'un texte scientifique par traitement du
langage naturel introduit du bruit et un risque de confusion entre le
vocabulaire (surface) et la relation (structure). À l'inverse, une
équation mathématique exprimée en LaTeX sépare parfaitement la surface
et la structure : * \textbf{La surface (\(U\))} est donnée par le texte
de l'article, ses mots-clés et sa catégorie disciplinaire arXiv (ex.
\texttt{q-bio.PE} {[}Populations et Évolution{]} vs
\texttt{cond-mat.stat-mech} {[}Mécanique Statistique{]}). * \textbf{La
structure pure (\(S\))} est donnée par l'arbre syntaxique abstrait
(\emph{Abstract Syntax Tree}, AST) de l'équation une fois toutes les
lettres de variables remplacées par un symbole générique \(\square\).
Par exemple, l'équation de diffusion de la chaleur et l'équation de
Black-Scholes (après changement de variable) partagent un sous-arbre
d'opérateurs parabolique identique
(\(\frac{\partial \square}{\partial t} = c \frac{\partial^2 \square}{\partial x^2}\)),
bien que leurs vocabulaires de surface soient totalement disjoints.

\subsubsection*{7.2. Protocole étape par étape (4 jours)}

\begin{itemize}
\tightlist
\item
  \textbf{Jour 1 (Collecte et construction des graphes à \(t_0 = 2014\))
  :} Télécharger via l'API OpenAlex et le dépôt bulk d'arXiv un
  échantillon de 100 000 articles publiés entre 2000 et 2014 couvrant
  quatre domaines (\texttt{physics}, \texttt{math}, \texttt{cs},
  \texttt{q-bio}). Extraire les concepts \(V_{2014}\) (termes techniques
  normalisés d'OpenAlex apparaissant dans au moins 20 articles).

  \begin{itemize}
  \tightlist
  \item
    Pour chaque concept \(u \in V_{2014}\), calculer son vecteur de
    surface \(x_u(2014)\) avec \emph{SPECTER2} sur les résumés
    \(\le 2014\), et sa distribution sur les catégories arXiv
    \(p_u(2014)\), d'où \(d_U(u,v;2014)\).
  \item
    Extraire les blocs
    \texttt{\textbackslash{}begin\{equation\}...\textbackslash{}end\{equation\}}
    des articles mentionnant \(u\) (\(\le 2014\)), les convertir en
    arbres d'opérateurs anonymisés via \texttt{SymPy} /
    \texttt{LaTeX2SymPy}, et définir le réseau structurel
    \(\mathcal{N}_S(u, 2014)\) comme la forêt des 10 arbres d'opérateurs
    les plus fréquemment co-mentionnés avec \(u\) (normalisés par TF-IDF
    d'arbres).
  \end{itemize}
\item
  \textbf{Jour 2 (Constitution des cohortes de test sur 2015--2018,
  évaluées en 2024) :} Identifier toutes les paires de concepts
  \((u,v)\) inter-catégories (\(d_U(u,v;2014)\) dans le quartile
  supérieur, c'est-à-dire mariages éloignés en surface) qui n'avaient
  \textbf{jamais} co-existé dans un même article avant 2014, et qui ont
  été réunies pour la première fois dans un article publié entre 2015 et
  2018. Extraire deux cohortes de taille \(N = 400\) chacune, appariées
  exactement par distance de surface \(d_U(u,v;2014)\) et par popularité
  des parents (\(\deg_U(u,2014) + \deg_U(v,2014)\)) :

  \begin{itemize}
  \tightlist
  \item
    \textbf{Groupe Succès (Mariages féconds, \(N=400\)) :} paires dont
    l'article pionnier (2015--2018) figure dans le top 10 \% des
    citations en 2024 \emph{et} qui ont été réutilisées ensemble dans au
    moins 15 articles ultérieurs par d'autres auteurs.
  \item
    \textbf{Groupe Échec (Mariages mort-nés, \(N=400\)) :} paires dont
    l'article pionnier (2015--2018) a été publié, mais qui n'ont généré
    aucune ou au plus une reprise par d'autres auteurs jusqu'en 2024.
  \end{itemize}
\item
  \textbf{Jour 3 (Calcul en aveugle de \(d_S\) et \(\Phi\) à l'horizon
  2014) :} Pour les 800 paires, calculer uniquement à partir des forêts
  d'équations de \textbf{2014} :

  \begin{enumerate}
  \def\labelenumi{\arabic{enumi}.}
  \tightlist
  \item
    La distance de structure \(d_S(u,v;2014)\) via la distance d'édition
    d'arbres (\emph{Tree Edit Distance}, algorithme de Zhang-Shasha)
    entre les arbres d'opérateurs de \(u\) et de \(v\).
  \item
    Le score serrure-clé \(\Phi(u,v;2014)\) défini par la taille
    relative du plus grand sous-arbre d'opérateurs commun (\emph{Maximum
    Common Subtree}) entre \(\mathcal{N}_S(u,2014)\) et
    \(\mathcal{N}_S(v,2014)\), pénalisé si des opérateurs incompatibles
    (ex. opérateur linéaire vs non-linéaire au même nœud de dérivation)
    s'opposent (\(\chi\)).
  \end{enumerate}
\item
  \textbf{Jour 4 (Tests statistiques et décision) :} Comparer la
  distribution de \(\Phi(u,v;2014)\) entre le Groupe Succès et le Groupe
  Échec par un test U de Mann-Whitney, et estimer une régression
  logistique prédisant le succès à partir de \(\Phi(u,v;2014)\), en
  contrôlant pour la fermeture triadique de surface \(C_{uv}(2014)\)
  (modèle ABC de Swanson).
\end{itemize}

\subsubsection*{7.3. Prédiction, seuil et lecture}

\begin{itemize}
\tightlist
\item
  \textbf{Prédiction du modèle (H1 + H3) :} À distance de surface
  \(d_U\) élevée et égale, les mariages interdisciplinaires de
  2015--2018 qui deviendront féconds en 2024 possédaient \emph{déjà en
  2014} un isomorphisme caché significativement plus élevé (\(\Phi\)
  supérieur, \(d_S\) inférieur) entre leurs arbres d'équations
  respectifs que les mariages qui resteront stériles.
\item
  \textbf{Seuil décisionnel :}

  \begin{itemize}
  \tightlist
  \item
    \textbf{Validation :} L'aire sous la courbe ROC séparant les
    mariages féconds des mariages mort-nés sur la seule base de
    \(\Phi(u,v;2014)\) atteint \textbf{\(\text{AUC-ROC} \ge 0,65\)} avec
    \textbf{\(p < 0,005\)}, et le rapport de cotes (\emph{odds ratio})
    de \(\Phi\) standardisé dans la régression logistique contrôlée par
    \(C_{uv}(2014)\) est \textbf{\(\text{OR} \ge 1,50\)}.
  \item
    \textbf{Réfutation :} Si \textbf{\(\text{AUC-ROC} < 0,55\)} (ou
    \(p > 0,05\)), ou si l'ajout de \(\Phi(u,v;2014)\) n'améliore pas la
    log-vraisemblance par rapport au seul terme de transitivité lexicale
    \(C_{uv}(2014)\).
  \end{itemize}
\item
  \textbf{Lecture épistémologique du résultat :}

  \begin{itemize}
  \tightlist
  \item
    Si le seuil est franchi, on dispose d'une preuve causale directe
    (puisque les équations de 2014 précèdent l'idée de 2015--2018) que
    la proximité structurelle cachée préexiste au mariage conceptuel et
    gouverne sa fécondité scientifique. Le modèle peut alors servir de
    moteur de recommandation (\emph{collisionneur}) en cherchant dans le
    graphe de l'année \(T\) les paires à fort \(d_U\) et fort \(\Phi\)
    non encore explorées.
  \item
    Si le seuil n'est pas franchi, H1 et H3 sont réfutées sous leur
    forme réaliste : cela signifie que la compatibilité mathématique («
    serrure et clé ») n'est pas une propriété préexistante découverte
    par les chercheurs, mais une construction plastique façonnée
    \emph{après coup} pour forcer deux domaines à s'ajuster.
  \end{itemize}
\end{itemize}

\begin{center}\rule{0.5\linewidth}{0.5pt}\end{center}

\subsection*{Références bibliographiques}

\begin{itemize}
\tightlist
\item
  Beltagy, I., Lo, K., \& Cohan, A. (2019). SciBERT: A pretrained
  language model for scientific text. \emph{Proceedings of the 2019
  Conference on Empirical Methods in Natural Language Processing
  (EMNLP)}, 3615--3620.
\item
  Berlyne, D. E. (1960). \emph{Conflict, Arousal, and Curiosity}.
  McGraw-Hill.
\item
  Boden, M. A. (2004). \emph{The Creative Mind: Myths and Mechanisms}
  (2nd ed.). Routledge.
\item
  Borgan, Ø., Goldstein, L., \& Langholz, B. (1995). Methods for the
  analysis of sampled cohort data in the Cox proportional hazards model.
  \emph{The Annals of Statistics}, 23(5), 1749--1778.
\item
  Bromham, L., Dinnage, R., \& Hua, X. (2016). Interdisciplinary
  research has consistently lower funding success. \emph{Nature},
  534(7609), 684--687.
\item
  Burt, R. S. (2004). Structural holes and good ideas. \emph{American
  Journal of Sociology}, 110(2), 349--399.
\item
  Davila, K., Joshi, R., Setlur, S., Govindaraju, V., \& Zanibbi, R.
  (2019). Tangent-CFT: An embedding model for mathematical formulas.
  \emph{Proceedings of the 2019 ACM SIGIR International Conference on
  Theory of Information Retrieval}, 41--44. \emph{(à vérifier pagination
  exacte)}
\item
  Davis, M. S. (1971). That's interesting! Towards a phenomenology of
  sociology and a sociology of phenomenology. \emph{Philosophy of the
  Social Sciences}, 1(2), 309--344.
\item
  Falkenhainer, B., Forbus, K. D., \& Gentner, D. (1989). The
  structure-mapping engine: Algorithm and examples. \emph{Artificial
  Intelligence}, 41(1), 1--63.
\item
  Foster, J. G., Rzhetsky, A., \& Evans, J. A. (2015). Tradition and
  innovation in scientists' research strategies. \emph{American
  Sociological Review}, 80(5), 875--908.
\item
  Funk, R. J., \& Owen-Smith, J. (2017). A dynamic network measure of
  technological change. \emph{Management Science}, 63(3), 791--817.
\item
  Galison, P. (2003). \emph{Einstein's Clocks, Poincaré's Maps: Empires
  of Time}. W. W. Norton \& Company.
\item
  Gentner, D. (1983). Structure-mapping: A theoretical framework for
  analogy. \emph{Cognitive Science}, 7(2), 155--170.
\item
  Gentner, D., \& Markman, A. B. (1997). Structure mapping in analogy
  and similarity. \emph{American Psychologist}, 52(1), 45--56.
\item
  Hamilton, W. L., Leskovec, J., \& Jurafsky, D. (2016). Diachronic word
  embeddings reveal statistical laws of semantic change.
  \emph{Proceedings of the 54th Annual Meeting of the Association for
  Computational Linguistics (ACL)}, 1489--1501.
\item
  Hofstadter, D., \& Sander, E. (2013). \emph{Surfaces and Essences:
  Analogy as the Fuel and Fire of Thinking}. Basic Books.
\item
  Holyoak, K. J., \& Thagard, P. (1995). \emph{Mental Leaps: Analogy in
  Creative Thought}. MIT Press.
\item
  Itti, L., \& Baldi, P. (2009). Bayesian surprise attracts human
  attention. \emph{Vision Research}, 49(10), 1295--1306.
\item
  Jüttner, A., \& Madarasi, P. (2018). VF2++---An improved subgraph
  isomorphism algorithm. \emph{Discrete Applied Mathematics}, 242,
  69--81.
\item
  Kauffman, S. A. (1995). \emph{At Home in the Universe: The Search for
  the Laws of Self-Organization and Complexity}. Oxford University
  Press.
\item
  Kitcher, P. (1981). Explanatory unification. \emph{Philosophy of
  Science}, 48(4), 507--531.
\item
  Kleinberg, J. (2003). Bursty and hierarchical structure in streams.
  \emph{Data Mining and Knowledge Discovery}, 7(4), 373--397.
\item
  Koestler, A. (1964). \emph{The Act of Creation}. Hutchinson.
\item
  Kuhn, T. S. (1962). \emph{The Structure of Scientific Revolutions}.
  University of Chicago Press.
\item
  Lakatos, I. (1976). \emph{Proofs and Refutations: The Logic of
  Mathematical Discovery}. Cambridge University Press.
\item
  Li, J., Yin, Y., Fortunato, S., \& Wang, D. (2019). A dataset of
  publication records for Nobel laureates. \emph{Scientific Data}, 6,
  33.
\item
  Lo, K., Wang, L. L., Neumann, M., Kinney, R., \& Weld, D. S. (2020).
  S2ORC: The Semantic Scholar Open Research Corpus. \emph{Proceedings of
  the 58th Annual Meeting of the Association for Computational
  Linguistics (ACL)}, 4969--4983.
\item
  Merton, R. K. (1961). Singletons and multiples in scientific
  discovery: A chapter in the sociology of science. \emph{Proceedings of
  the American Philosophical Society}, 105(5), 470--486.
\item
  Mikolov, T., Sutskever, I., Chen, K., Corrado, G. S., \& Dean, J.
  (2013). Distributed representations of words and phrases and their
  compositionality. \emph{Advances in Neural Information Processing
  Systems (NeurIPS)}, 26, 3111--3119.
\item
  Newman, M. E. J., \& Ziff, R. M. (2000). Efficient Monte Carlo
  algorithm and high-precision results for percolation. \emph{Physical
  Review Letters}, 85(19), 4104--4107.
\item
  Nooteboom, B., Van Haverbeke, W., Duysters, G., Gilsing, V., \& Van
  den Oord, A. (2007). Optimal cognitive distance and absorptive
  capacity. \emph{Research Policy}, 36(7), 1016--1034.
\item
  Ogburn, W. F., \& Thomas, D. (1922). Are inventions inevitable? A note
  on social evolution. \emph{Political Science Quarterly}, 37(1),
  83--98.
\item
  Rindflesch, T. C., \& Fiszman, M. (2003). The interaction of domain
  knowledge and linguistic structure in natural language processing:
  Interpreting hypernymic propositions in biomedical text. \emph{Journal
  of Biomedical Informatics}, 36(6), 462--477.
\item
  Schmidhuber, J. (2010). Formal theory of creativity, fun, and
  intrinsic motivation (1990--2010). \emph{IEEE Transactions on
  Autonomous Mental Development}, 2(3), 230--247.
\item
  Schultz, W., Dayan, P., \& Montague, P. R. (1997). A neural substrate
  of prediction and reward. \emph{Science}, 275(5306), 1593--1599.
\item
  Simonton, D. K. (2004). \emph{Creativity in Science: Chance, Logic,
  Genius, and Zeitgeist}. Cambridge University Press.
\item
  Spivak, D. I., \& Kent, R. E. (2012). Ologs: A categorical framework
  for knowledge representation. \emph{PLOS ONE}, 7(1), e24274.
\item
  Swanson, D. R. (1986). Undiscovered public knowledge. \emph{The
  Library Quarterly}, 56(2), 103--118.
\item
  Traag, V. A., Waltman, L., \& van Eck, N. J. (2019). From Louvain to
  Leiden: Guaranteeing well-connected communities. \emph{Scientific
  Reports}, 9, 5233.
\item
  Tria, F., Loreto, V., Servedio, V. D. P., \& Strogatz, S. H. (2014).
  The dynamics of correlated novelties. \emph{Scientific Reports}, 4,
  5890.
\item
  Tshitoyan, V., Dagdelen, J., Weston, L., Dunn, A., Rong, Z., Kononova,
  O., Persson, K. A., Ceder, G., \& Jain, A. (2019). Unsupervised word
  embeddings capture latent knowledge from materials science literature.
  \emph{Nature}, 571(7763), 95--98.
\item
  Uzzi, B., Mukherjee, S., Stringer, M., \& Jones, B. (2013). Atypical
  combinations and scientific impact. \emph{Science}, 342(6157),
  468--472.
\item
  Vayer, T., Chapel, L., Flamary, R., Tavenard, R., \& Courty, N.
  (2019). Optimal transport for structured data with application on
  graphs. \emph{Proceedings of the 36th International Conference on
  Machine Learning (ICML)}, 97, 6275--6284.
\item
  Wadden, D., Wennberg, U., Luan, Y., \& Hajishirzi, H. (2019). Entity,
  relation, and event extraction with contextualized span
  representations. \emph{Proceedings of the 2019 Conference on Empirical
  Methods in Natural Language Processing (EMNLP)}, 5784--5789.
\item
  Wu, L., Wang, D., \& Evans, J. A. (2019). Large teams develop and
  small teams disrupt science and technology. \emph{Nature}, 566(7744),
  378--382.
\end{itemize}
\clearpage
\section*{Annexe B — Réponse de Carbon}\addcontentsline{toc}{subsection}{Annexe B — Réponse de Carbon}
{\small Source : \texttt{reponses-google-v1/reponse-carbon-v1.md}.}\par
\subsection*{1. Confrontation à la littérature}

Les dix hypothèses entremêlent trois questions analytiquement distinctes
qu'un modèle testable doit séparer : 1. \textbf{L'intérêt perçu
(psychosociologique)} : ce qui capte l'attention d'un public humain
donné à l'instant \(t\), relativement à ses croyances antérieures
(Davis, 1971 ; Itti \& Baldi, 2009). Une idée fausse peut être
fascinante. 2. \textbf{La validité et la fécondité épistémique} : ce qui
rend une proposition vraie, unificatrice et génératrice de résultats
ultérieurs vérifiés (Lakatos, 1970 ; Kitcher, 1981 ; Uzzi, Mukherjee,
Stringer, \& Jones, 2013). 3. \textbf{Le mécanisme de production
(heuristique)} : le processus cognitif, social ou algorithmique qui fait
émerger l'idée candidate (Koestler, 1964 ; Swanson, 1986 ; Simonton,
2004).

Confondre (1) et (2) conduit à sélectionner des métaphores séduisantes
mais stériles ; confondre (2) et (3) fait prendre un filtre rétrospectif
pour un générateur tractable.

\begin{itemize}
\tightlist
\item
  \textbf{H1 (mariage éloigné) --- Soutenue (questions 1, 2 et 3).} La
  \emph{bisociation} de matrices de pensée indépendantes a été théorisée
  par Koestler (1964, \emph{The Act of Creation}, Hutchinson) et
  formalisée en sciences cognitives par la théorie de l'appariement
  structurel (\emph{structure-mapping}), qui oppose la similarité des
  attributs de surface à l'isomorphisme relationnel abstrait (Gentner,
  1983, « Structure-mapping: A theoretical framework for analogy »,
  \emph{Cognitive Science}, 7(2), 155--170 ; Hofstadter \& Sander, 2013,
  \emph{L'Analogie : cœur de la pensée}, Odile Jacob). En science de la
  science, les stratégies reliant des concepts séparés à travers les
  frontières disciplinaires sont rares et risquées, mais concentrent les
  percées majeures (Foster, Rzhetsky, \& Evans, 2015, « Tradition and
  innovation in scientists' research strategies », \emph{American
  Sociological Review}, 80(5), 875--908).
\item
  \textbf{H2 (cascades et transferts) --- Soutenue (question 3), mais
  insuffisante seule (question 2).} Le chaînage transitif
  \(A\text{--}B\) et \(B\text{--}C\) engendrant l'hypothèse neuve
  \(A\text{--}C\) constitue le modèle ABC de découverte fondée sur la
  littérature (Swanson, 1986, « Fish oil, Raynaud's syndrome, and
  undiscovered public knowledge », \emph{Perspectives in Biology and
  Medicine}, 30(1), 7--18). Le transfert vers un domaine cible jeune
  relève du transfert analogique contraint par le but (Holyoak \&
  Thagard, 1995, \emph{Mental Leaps: Analogy in Creative Thought}, MIT
  Press). Toutefois, sans filtre structurel (H3), le chaînage ABC
  souffre d'une explosion de faux positifs par polysémie ou glissement
  conceptuel --- écueil manifeste dans l'exemple de l'anesthésie
  transposée aux modèles de langage si l'on ne spécifie pas quel
  invariant computationnel remplace les boucles thalamocorticales
  biologiques.
\item
  \textbf{H3 (serrure et clé, symétrie de Noether) --- Scindée :
  soutenue pour l'alignement amont ; métaphore sans contenu testable
  pour Noether et le lagrangien.} La nécessité d'aligner des systèmes
  connectés de relations d'ordre supérieur plutôt que des prédicats
  isolés est le \emph{principe de systématicité} de Gentner (1983),
  formalisable par les foncteurs préservant la structure en théorie des
  catégories (Mac Lane, 1971, \emph{Categories for the Working
  Mathematician}, Springer ; Spivak, 2014, \emph{Category Theory for the
  Sciences}, MIT Press). En revanche, l'invocation du théorème de
  Noether et d'un lagrangien est une métaphore non testable : Noether
  requiert l'action d'un groupe de Lie continu de symétries sur une
  intégrale d'action différentiable, structure absente d'un réseau
  discret de concepts.
\item
  \textbf{H4 (ni trop près ni trop loin) --- Soutenue avec nuance
  asymétrique (questions 1 et 2).} En psychologie de la curiosité,
  l'intérêt suit une courbe en U inversé selon l'écart aux attentes
  (Berlyne, 1960, \emph{Conflict, Arousal, and Curiosity}, McGraw-Hill).
  En bibliométrie sur les 17,9 millions d'articles du Web of Science
  étudiés par Uzzi et al.~(2013, « Atypical combinations and scientific
  impact », \emph{Science}, 342(6157), 468--472), la fécondité maximale
  n'est pas atteinte à une distance uniforme intermédiaire, mais par
  l'association simultanée d'un socle hautement conventionnel et d'une
  queue d'appariements très atypiques.
\item
  \textbf{H5 (champ d'interaction, air du temps, découvertes
  simultanées) --- Soutenue en histoire et sociologie des sciences ;
  métaphore réductible pour le ferromagnétisme et les protéines.}
  L'apparition quasi simultanée de découvertes indépendantes sous
  l'effet d'une maturation technique et culturelle commune est
  solidement établie (Ogburn \& Thomas, 1922, « Are inventions
  inevitable? A note on social evolution », \emph{Political Science
  Quarterly}, 37(1), 83--98, qui recensent 148 cas historiques ; Merton,
  1961, « Singletons and multiples in scientific discovery »,
  \emph{Proceedings of the American Philosophical Society}, 105(5),
  470--486). Le rôle exact de la synchronisation ferroviaire et des
  réseaux d'horloges dans la relativité restreinte est démontré par
  Galison (2003, \emph{Einstein's Clocks, Poincaré's Maps}, W. W.
  Norton), et l'analogie hydrodynamique de l'évaporation des trous noirs
  provient d'Unruh (1981, « Experimental black-hole evaporation? »,
  \emph{Physical Review Letters}, 46(21), 1351--1353). L'appariement
  stochastique après agitation rejoint le modèle de variation aveugle et
  rétention sélective (Campbell, 1960, « Blind variation and selective
  retention in creative thought », \emph{Psychological Review}, 67(6),
  380--400 ; Simonton, 2004, \emph{Creativity in Science}, Cambridge
  University Press). Les images ferromagnétique et protéique n'ont de
  valeur qu'une fois réduites à un processus ponctuel à facteur latent
  partagé.
\item
  \textbf{H6 (connexion et distinction) --- Soutenue (questions 1 et
  2).} Davis (1971, « That's interesting! », \emph{Philosophy of the
  Social Sciences}, 1, 309--344) montre qu'une thèse jugée intéressante
  soit unifie des phénomènes crus hétérogènes, soit dissocie des
  phénomènes crus identiques. Cognitivement, l'alignement structurel
  fait émerger des \emph{différences alignables} qui corrigent les
  fausses équivalences de surface (Gentner \& Markman, 1997, « Structure
  mapping in analogy and similarity », \emph{American Psychologist},
  52(1), 45--56).
\item
  \textbf{H7 (profondeur de l'idée fille) --- Soutenue (question 2).} En
  épistémologie, la profondeur explicative --- illustrée par la synthèse
  newtonienne --- réside dans l'unification qui réduit le nombre de
  schémas argumentatifs indépendants requis pour dériver les phénomènes
  (Kitcher, 1981, « Explanatory unification », \emph{Philosophy of
  Science}, 48(4), 507--531). Dans les réseaux de citations, une idée
  profonde qui supplante ses antécédents en devenant elle-même le nœud
  de référence se mesure par l'indice de disruption (Funk \& Owen-Smith,
  2017, « A dynamic network measure of technological change »,
  \emph{Management Science}, 63(3), 791--817 ; Wu, Wang, \& Evans, 2019,
  « Large teams develop and small teams disrupt science and technology
  », \emph{Nature}, 566(7744), 378--382).
\item
  \textbf{H8 (émergence et percolation) --- Partiellement soutenue
  (question 3).} L'avantage créatif des courtiers situés à l'interface
  de « trous structurels » entre communautés est établi (Burt, 2004, «
  Structural holes and good ideas », \emph{American Journal of
  Sociology}, 110(2), 349--399). La naissance d'une discipline présente
  une densification topologique assimilable à une transition de
  percolation dans le réseau de coauteurs (Bettencourt, Kaiser, Kaur,
  Castillo-Chávez, \& Wojick, 2008, « Population modeling of the
  emergence and development of scientific fields »,
  \emph{Scientometrics}, 75(3), 495--518), bien que la preuve d'une
  percolation conceptuelle pure reste ouverte.
\item
  \textbf{H9 (structure cachée du langage) --- Soutenue avec limite
  critique.} Les plongements lexicaux non supervisés capturent des
  régularités relationnelles linéaires (Mikolov, Sutskever, Chen,
  Corrado, \& Dean, 2013, « Distributed representations of words and
  phrases and their compositionality », \emph{NeurIPS}, 26, 3111--3119)
  et permettent même des prédictions rétrospectives de matériaux
  fonctionnels (Tshitoyan et al., 2019, « Unsupervised word embeddings
  capture latent knowledge from materials science literature »,
  \emph{Nature}, 571(7763), 95--98). Cependant, les vecteurs de
  cooccurrence confondent proximité thématique de surface et
  isomorphisme relationnel abstrait : extraire ce dernier exige des
  représentations explicites des relations (Lu, Wu, \& Holyoak, 2019, «
  Emergence of analogy from relation learning », \emph{PNAS}, 116(10),
  4176--4181).
\item
  \textbf{H10 (signal de surprise dopaminergique) --- Soutenue sous
  réserve d'une distinction neurobiologique (questions 1 et 3).} Les
  neurones dopaminergiques mésencéphaliques codent canoniquement une
  \emph{erreur de prédiction de récompense} signée plutôt qu'une
  surprise informationnelle pure non signée (Schultz, Dayan, \&
  Montague, 1997, « A neural substrate of prediction and reward »,
  \emph{Science}, 275(5306), 1593--1599). Chez un agent artificiel
  explorateur, le signal pertinent est la surprise bayésienne (Itti \&
  Baldi, 2009, « Bayesian surprise attracts human attention »,
  \emph{Vision Research}, 49(10), 1295--1306) ou le gain de compression
  intrinsèque (Schmidhuber, 2010, « Formal theory of creativity, fun,
  and intrinsic motivation », \emph{IEEE Transactions on Autonomous
  Mental Development}, 2(3), 230--247).
\end{itemize}

\begin{center}\rule{0.5\linewidth}{0.5pt}\end{center}

\subsection*{2. Modèle formel minimal}

\subsubsection*{2.1 Espace des concepts multicouche daté}

À chaque date \(t \in \mathbb{R}_+\), l'état du savoir publié est
représenté par un graphe multicouche
\(\mathcal{G}(t) = \big(V(t), \mathcal{L}_{\text{surf}}(t), \mathcal{L}_{\text{struct}}(t), \mathcal{L}_{\text{soc}}(t)\big)\),
où \(V(t)\) est l'ensemble fini des concepts attestés avant \(t\) : *
\textbf{Couche d'usage \(\mathcal{L}_{\text{surf}}(t)\)} : associe à
tout \(u \in V(t)\) un plongement distributionnel unitaire
\(x_u(t) \in \mathbb{S}^{d-1}\) appris sur le texte antérieur à \(t\),
et une distribution empirique de publication
\(p_{\text{rev}}(\cdot \mid u, t)\) sur les revues/séminaires. *
\textbf{Couche structurelle \(\mathcal{L}_{\text{struct}}(t)\)} :
multigraphe orienté typé \((V(t), E(t), \mathcal{R})\), où une arête
\((z_1, r, z_2) \in E(t)\) exprime une dépendance formelle, mathématique
ou causale attestée de type \(r \in \mathcal{R}\). Une projection
ontologique \(\pi : \mathcal{R} \to \mathcal{R}_{\text{abs}}\) envoie
les relations disciplinaires vers un alphabet relationnel abstrait
(monotonie croissante/décroissante, dérivation, loi de conservation,
rétroaction négative, etc.). Pour un entier fixé \(k \ge 1\), on note
\(G_u(t) = (V_u, E_u)\) le sous-graphe relationnel amont de rayon \(k\)
convergeant vers \(u\). * \textbf{Couche sociale
\(\mathcal{L}_{\text{soc}}(t)\)} : matrice bipartite auteurs--concepts
\(M(t) \in \mathbb{R}_{\ge 0}^{|A(t)| \times |V(t)|}\) sur la fenêtre
glissante \([t-\tau, t]\), induisant les distributions
\(p_{\text{soc}}(a \mid u, t)\) et \(p_{\text{soc}}(u \mid a, t)\).

\subsubsection*{2.2 Deux distances inter-domaines et critère serrure-clé}

La \textbf{distance de surface} combine la séparation des communautés
humaines et l'orthogonalité lexicale, bornée dans \([0,1]\) quelle que
soit la paire de disciplines :
\[d_{\text{surf}}(u, v; t) = \frac{1}{2}\sqrt{\text{JSD}_2\!\big(p_{\text{soc}}(\cdot \mid u, t) \parallel p_{\text{soc}}(\cdot \mid v, t)\big)} + \frac{1 - x_u(t)^\top x_v(t)}{4}\]
où \(\text{JSD}_2\) désigne la divergence de Jensen-Shannon en
logarithme de base 2.

Pour comparer la structure à travers deux domaines sans vocabulaire
commun, soit \(\Phi(u,v)\) l'ensemble des applications partielles
injectives \(\phi : S \to V_v\) définies sur \(S \subseteq V_u\) telles
que \(u \in S\) et \(\phi(u) = v\). Soit
\(\mathcal{I} \subset \mathcal{R}_{\text{abs}} \times \mathcal{R}_{\text{abs}}\)
la relation symétrique d'incompatibilité logique ou fonctionnelle (par
exemple monotonies de signes opposés). On définit la \textbf{couverture}
\(\kappa(\phi) \in [0,1]\) et le taux de \textbf{conflits}
\(\chi(\phi) \in [0,1]\) :
\[\kappa(\phi) = \frac{\sum_{(z_1, r, z_2) \in E_u(S)} \omega(z_1, z_2)\,\mathbf{1}\!\big[(\phi(z_1), \pi(r), \phi(z_2)) \in \pi(E_v)\big]}{\sum_{(z_1, r, z_2) \in E_u} \omega(z_1, z_2)}\]
\[\chi(\phi) = \frac{1}{|E_u(S)|} \sum_{(z_1, r, z_2) \in E_u(S)} \mathbf{1}\!\big[\exists r' \in \mathcal{R}_{\text{abs}} : (\phi(z_1), r', \phi(z_2)) \in \pi(E_v) \land (\pi(r), r') \in \mathcal{I}\big]\]
où \(\omega(z_1, z_2) > 0\) décroît avec la distance géodésique à \(u\)
(et \(\chi(\phi)=0\) si \(E_u(S)=\varnothing\)). Le \textbf{score
serrure-clé} et la \textbf{distance de structure} valent
(\(\gamma > 0\)) :
\[\mathcal{K}(u, v; t) = \max_{\phi \in \Phi(u,v)} \kappa(\phi)\exp\!\big(-\gamma \chi(\phi)\big) \in [0, 1], \qquad d_{\text{struct}}(u, v; t) = 1 - \mathcal{K}(u, v; t)\]
Ce formalisme capture aussi \textbf{H6} : deux concepts \(u, u'\)
confondus en surface (\(d_{\text{surf}}(u, u'; t) \approx 0\)) forment
une \emph{fausse identité} révélée par collision avec un concept sonde
\(v\) dès que \(|\mathcal{K}(u, v; t) - \mathcal{K}(u', v; t)| \gg 0\).

\subsubsection*{2.3 Champ temporel et probabilité de naissance d'une
idée}

Soit \(\mathcal{M}_{\text{inst}}(t)\) le répertoire des instruments
expérimentaux, bases de données et formalismes mathématiques actifs à
\(t\). Chaque outil \(m \in \mathcal{M}_{\text{inst}}(t)\) possède une
impulsion d'adoption globale
\(\alpha_m(t) = \max\!\big(0, \frac{d}{dt}\log N_m(t)\big)\) (\(N_m(t)\)
étant son volume cumulé de publications) et un taux de pénétration
locale \(h_m(u, t) \in [0,1]\) parmi les travaux traitant de \(u\). Le
\textbf{couplage de champ} entre \(u\) et \(v\) s'écrit :
\[F(u, v; t) = \sum_{m \in \mathcal{M}_{\text{inst}}(t)} \alpha_m(t)\, h_m(u, t)\, h_m(v, t)\]
Soit
\(C_{\text{soc}}(u, v; t) = \sum_{a \in A(t)} p_{\text{soc}}(u \mid a, t)\,p_{\text{soc}}(v \mid a, t)\)
la surface de contact directe par chercheurs passeurs. Pour deux
concepts non encore substantiellement liés avant \(t\), l'intensité
instantanée (taux de hasard) de naissance de l'idée hybride
\(w = (u, v, \phi^*)\) dans \([t, t+dt)\) est :
\[\lambda_{uv}(t) = \big(\eta_0 C_{\text{soc}}(u, v; t) + \eta_1 F(u, v; t)\big) \cdot \sigma\!\Big(\beta_1 \mathcal{K}(u, v; t) - \beta_2 \big(d_{\text{surf}}(u, v; t) - \mu\big)^2 - \beta_0\Big)\]
où \(\sigma(z) = (1+e^{-z})^{-1}\),
\(\eta_0, \eta_1, \beta_1, \beta_2 > 0\), et \(\mu \in (0,1)\) est la
distance de surface optimale (\textbf{H4}). Si \(F(u,v;t)\) croît
brusquement, le nombre \(N_{uv}(\Delta t)\) d'équipes indépendantes
formant le mariage avant diffusion mutuelle suit une loi de Poisson de
moyenne
\(\Lambda_{uv}(\Delta t) = \int_t^{t+\Delta t} \lambda_{uv}(s)\,ds\),
engendrant des découvertes simultanées avec probabilité
\(1 - (1 + \Lambda_{uv})e^{-\Lambda_{uv}}\) (\textbf{H5}).

\subsubsection*{2.4 Intérêt de l'idée fille et dynamique d'émergence}

Le modèle sépare les deux sens de la valeur de \(w\) : 1.
\textbf{Intérêt perçu / surprise cognitive pour une communauté \(c\)
(H1, H10)} : surprise bayésienne
\(S_c(w; t) = D_{\text{KL}}\big(P_c(\Theta \mid \mathcal{G}(t), w) \parallel P_c(\Theta \mid \mathcal{G}(t))\big)\)
sur les croyances relationnelles latentes \(\Theta\) de la communauté
\(c\). Sous un a priori où la probabilité attendue d'un lien décroît
exponentiellement avec \(d_{\text{surf}}\) et où la révision effective
d'ordre supérieur est proportionnelle à \(\mathcal{K}\), son estimateur
observable d'ordre un est :
\[\mathcal{I}_{\text{perçu}}(w; t) = d_{\text{surf}}(u, v; t) \cdot \mathcal{K}(u, v; t)\]
2. \textbf{Profondeur explicative et fécondité (H7)} : * \emph{Ex ante
(unification structurelle)} : gain de longueur de description minimale
(Rissanen, 1978, « Modeling by shortest data description »,
\emph{Automatica}, 14(5), 465--471) obtenu en factorisant \(G_u\) et
\(G_v\) par le schéma abstrait commun \(G_{\phi^*}^{\text{abs}}\) :
\[\Delta \mathcal{L}(w; t) = \mathcal{L}(G_u \cup G_v) - \big[\mathcal{L}(G_{\phi^*}^{\text{abs}}) + \mathcal{L}(G_u \mid G_{\phi^*}^{\text{abs}}) + \mathcal{L}(G_v \mid G_{\phi^*}^{\text{abs}})\big]\]
* \emph{Ex post à horizon \(\Delta T\)} :
\(\mathcal{F}_{\text{obs}}(w; t+\Delta T) = \text{PR}(w, t+\Delta T) \cdot \frac{1 + \text{CD}(w, t+\Delta T)}{2}\),
produit de la centralité PageRank de \(w\) et de son indice de
disruption \(\text{CD} \in [-1, 1]\) (Funk \& Owen-Smith, 2017).

Pour l'\textbf{émergence inter-domaines (H8)} entre deux ensembles
\(D_1, D_2 \subset V(t)\), soit \(G_{12}(t)\) le graphe bipartite des
paires structurellement viables
(\(\mathcal{K}(u,v;t) \ge \tau_{\mathcal{K}}\)), où chaque arête
s'active sur \([t, t+\Delta t]\) avec probabilité
\(p_{uv} = 1 - \exp(-\Lambda_{uv}(\Delta t))\). Une composante géante
hybride émerge lorsque le rayon spectral \(\rho(B_{12} \odot P) > 1\),
où \(B_{12}\) est la matrice non-rétropropagée de Hashimoto de
\(G_{12}\) et \(P\) l'opérateur des probabilités d'arêtes (Karrer,
Newman, \& Zdeborová, 2014, « Percolation on sparse networks »,
\emph{Physical Review Letters}, 113(20), 208702).

\subsubsection*{2.5 Statut mathématique exact des analogies physiques}

\begin{itemize}
\tightlist
\item
  \textbf{Puits de potentiel --- Conservée sous forme combinatoire
  discrète.} Équivalent exact : fonction d'énergie d'alignement
  \(E(\phi) = -\log \kappa(\phi) + \gamma \chi(\phi)\) sur \(\Phi(u,v)\)
  ; la « profondeur » sans ambiguïté d'une analogie se mesure par
  l'écart spectral d'énergie
  \(\Delta E = E(\phi_{(2)}) - E(\phi_{(1)})\) entre le deuxième
  meilleur alignement et l'optimum. Toute dynamique inertielle continue
  dans ce puits est écartée.
\item
  \textbf{Champ et polarisation --- Conservés sous forme statistique.}
  Le champ externe est le vecteur mesurable
  \(\mathbf{h}(u,t) = (h_m(u,t))_m\) pondéré par \(\alpha_m(t)\) ; la «
  polarisation » d'une discipline \(D\) vers l'outil \(m\) est le
  paramètre d'ordre empirique
  \(M_m(D, t) = \frac{1}{|D|}\sum_{u \in D} h_m(u, t) \in [0, 1]\).
\item
  \textbf{Résonance --- Écartée explicitement.} Une résonance exige un
  oscillateur entretenu caractérisé par une fréquence propre et une
  phase continue. Un graphe conceptuel n'ayant pas de phase périodique
  canonique mesurable, cette métaphore est redondante avec
  \(\mathcal{K}(u,v;t)\) et \(F(u,v;t)\).
\item
  \textbf{Théorème de Noether et lagrangien --- Écartés sous leur forme
  différentielle, remplacés par l'homomorphisme fonctoriel.} Faute de
  variété différentiable, de variables conjuguées et d'intégrale
  d'action falsifiable, l'analogie lagrangienne n'ajoute aucune
  prédiction. L'équivalent exact et mesurable de la « conservation d'un
  invariant sous transformation » est la préservation des relations
  abstraites \(\pi(E_u)\) par l'injection partielle \(\phi\), mesurée
  par \(\kappa(\phi)\), \(\chi(\phi)\) et \(\Delta \mathcal{L}(w; t)\).
\item
  \textbf{Percolation --- Conservée exactement.} Percolation de liens
  hétérogène sur le graphe bipartite \(G_{12}(t)\), avec paramètre de
  contrôle \(\rho(B_{12} \odot P)\) et paramètre d'ordre mesurable
  \(S_{\max}(t)/|D_1 \cup D_2|\) (fraction de concepts dans la plus
  grande composante connexe inter-domaines).
\end{itemize}

\begin{center}\rule{0.5\linewidth}{0.5pt}\end{center}

\subsection*{3. Données brutes}

\begin{enumerate}
\def\labelenumi{\arabic{enumi}.}
\tightlist
\item
  \textbf{Corpus datés en texte intégral} : sources LaTeX
  d'\textbf{arXiv}, sous-ensemble en libre accès de \textbf{PubMed
  Central (PMC OA)} et \textbf{S2ORC} (Lo, Wang, Neumann, Kinney, \&
  Weld, 2020, « S2ORC: The Semantic Scholar Open Research Corpus »,
  \emph{Proceedings of ACL}, 4969--4983). \emph{Variables tirées} :
  vocabulaire contrôlé \(V(t)\), plongements \(x_u(t)\), mentions
  d'instruments/méthodes \(h_m(u,t)\) et impulsions \(\alpha_m(t)\),
  dates exactes de première co-mobilisation substantielle (dans une même
  section méthodologique ou démonstrative, et non en simple revue
  bibliographique).
\item
  \textbf{Graphes de citation et de coauteurs} : \textbf{OpenAlex}
  (Priem, Piwowar, \& Orr, 2022, \emph{arXiv:2205.01833}) et
  \textbf{SciSciNet} (Lin, Yin, Liu, \& Wang, 2023, « SciSciNet: A
  large-scale open data lake for the science of science research »,
  \emph{Scientific Data}, 10, 315). \emph{Variables tirées} : couche
  \(\mathcal{L}_{\text{soc}}(t)\), \(d_{\text{surf}}(u,v;t)\), courtage
  \(C_{\text{soc}}(u,v;t)\), centralité \(\text{PR}(w, t+\Delta T)\) et
  disruption \(\text{CD}(w, t+\Delta T)\).
\item
  \textbf{Ontologies et graphes relationnels typés} : \textbf{MeSH} et
  \textbf{SemMedDB} pour les prédications causales biomédicales
  horodatées (Kilicoglu, Shin, Fiszman, Rosemblat, \& Rindflesch, 2012,
  « SemMedDB: A PubMed-scale repository of biomedical semantic
  predications », \emph{Bioinformatics}, 28(23), 3158--3160),
  classifications \textbf{MSC} et \textbf{ACM CCS}, et historique
  horodaté des révisions de \textbf{Wikidata} (filtré avant \(t\) pour
  proscrire toute fuite du futur). \emph{Variables tirées} :
  \(\mathcal{L}_{\text{struct}}(t)\), projection abstraite \(\pi\),
  voisinages \(G_u(t)\).
\item
  \textbf{Corpus d'équations} : environnements mathématiques extraits
  des sources LaTeX d'arXiv et convertis en arbres d'opérateurs typés
  (suivant le protocole d'ARQMath ; Mansouri et al., 2022 --- référence
  exacte du rapport CLEF à vérifier, l'outil sous-jacent
  \textbf{LaTeXML} étant standard), en reliant chaque symbole à sa
  définition textuelle locale. \emph{Variables tirées} : homomorphismes
  exacts d'équations pour \(\mathcal{K}(u,v;t)\) et compression
  \(\Delta \mathcal{L}(w;t)\) en sciences exactes.
\item
  \textbf{Catalogues de découvertes et de découvertes simultanées} :
  articles associés aux prix Nobel (Li, Yin, Fortunato, \& Wang, 2019, «
  A dataset of publication records for Nobel laureates »,
  \emph{Scientific Data}, 6, 33), paires but/mécanisme en ingénierie
  (Hope, Chan, Kittur, \& Shahaf, 2017, « Accelerating innovation
  through analogy mining », \emph{Proceedings of KDD}, 235--243), liste
  historique d'Ogburn \& Thomas (1922), et \emph{articles jumeaux}
  contemporains identifiés par co-citation symétrique et dates de
  soumission séparées par un court intervalle sans auteurs ni références
  croisés (méthode des \emph{paper twins} de Bikard, Murray, \& Gans,
  2015, « Exploring trade-offs in the organization of scientific work »,
  \emph{Management Science}, 61(7), 1473--1495).
\item
  \textbf{Catalogues d'analogies ayant échoué} : (i) toutes les paires
  \((u,v)\) co-mobilisées pour la première fois à \(t\) dans arXiv/PMC
  mais ne recevant aucune reprise indépendante (hors autocitations) sur
  \([t, t+\Delta T]\) ; (ii) transferts mécanistes précliniques
  débouchant sur des essais négatifs enregistrés dans
  \textbf{ClinicalTrials.gov} ; (iii) articles rétractés pour invalidité
  scientifique dans la base \textbf{Retraction Watch / Crossref}.
\end{enumerate}

\begin{center}\rule{0.5\linewidth}{0.5pt}\end{center}

\subsection*{4. Méthodes et outils}

\begin{itemize}
\tightlist
\item
  \textbf{Extraction de \(V(t)\), \(E(t)\) et \(\mathbf{h}(u,t)\)} :
  reconnaissance d'entités (\emph{Task, Method, Material, Metric}) et
  extraction de relations typées par \textbf{SciBERT} (Beltagy, Lo, \&
  Cohan, 2019, \emph{Proceedings of EMNLP-IJCNLP}, 3615--3620) affiné
  sur \textbf{SciERC} (Luan, He, Ostendorf, \& Hajishirzi, 2018,
  \emph{Proceedings of EMNLP}, 3219--3232), ré-entraîné uniquement sur
  des textes antérieurs à la coupure temporelle. Pour les équations :
  conversion LaTeX vers arbres d'opérateurs via \textbf{LaTeXML}, puis
  canonisation algébrique sous \textbf{SymPy} (normalisation des
  variables muettes et de l'ordre commutatif).
\item
  \textbf{Plongements diachroniques et \(d_{\text{surf}}\)} :
  entraînement annuel de modèles \textbf{Word2Vec Skip-gram avec
  échantillonnage négatif (SGNS)} alignés d'une année sur l'autre par
  \textbf{procuste orthogonal} (Hamilton, Leskovec, \& Jurafsky, 2016,
  \emph{Proceedings of ACL}, 1489--1501), complétés au niveau
  documentaire par \textbf{SPECTER} (Cohan, Feldman, Beltagy, Downey, \&
  Weld, 2020, \emph{Proceedings of ACL}, 2270--2282) ré-entraîné avant
  \(t\). Partitionnement communautaire de contrôle par
  l'\textbf{algorithme de Leiden} (Traag, Waltman, \& van Eck, 2019,
  \emph{Scientific Reports}, 9, 5233).
\item
  \textbf{Calcul de \(\mathcal{K}(u,v;t)\) et \(d_{\text{struct}}\)} :

  \begin{itemize}
  \tightlist
  \item
    \emph{Régime exact} (\(|V_u| \le 25\), choix d'ingénierie) :
    formulation du maximum de \(\kappa(\phi) - \gamma' \chi(\phi)\) en
    programmation linéaire en nombres entiers mixtes (\textbf{MILP})
    résolue par \textbf{CP-SAT (Google OR-Tools)}.
  \item
    \emph{Passage à l'échelle} : transport optimal structuré
    \textbf{Fused Gromov-Wasserstein (FGW)} (Vayer, Chapel, Flamary,
    Tavenard, \& Courty, 2020, \emph{Algorithms}, 13(9), 212) via la
    bibliothèque \textbf{POT}, en définissant la matrice de coût
    exclusivement sur les types d'arêtes abstraits \(\pi(r)\) et les
    pénalités d'incompatibilité \(\mathcal{I}\), \emph{sans aucun accès
    aux étiquettes lexicales des nœuds}.
  \end{itemize}
\item
  \textbf{Champ \(F(u,v;t)\) et estimation de \(\lambda_{uv}(t)\)} :
  estimation causale de \(\alpha_m(t)\) par filtrage de Kalman
  unidirectionnel sur \([t_0, t]\). Vu la taille quadratique
  \(|V(t)|^2\), estimation de \(\lambda_{uv}(t)\) par \textbf{régression
  à risques proportionnels de Cox dépendante du temps avec
  échantillonnage cas-témoins imbriqué} (\emph{nested case-control}) et
  test de la courbure \textbf{H4} par \textbf{modèles additifs
  généralisés (GAM)} à splines cubiques pénalisées.
\item
  \textbf{Unification \(\Delta \mathcal{L}(w;t)\), percolation et
  surprise d'agent} : compression MDL sur graphes par l'algorithme
  \textbf{Subdue} ; valeur propre dominante de la matrice creuse de
  Hashimoto \(B_{12} \odot P\) par itération d'\textbf{Arnoldi} ;
  estimation de la surprise bayésienne \(S_c(w;t)\) par approximation de
  Laplace sur les poids de la dernière couche d'un \textbf{R-GCN}
  (Schlichtkrull et al., 2018, \emph{ESWC}, LNCS 10843, 593--607)
  entraîné à prédire les liens de la communauté \(c\) avant \(t\).
\end{itemize}

\begin{center}\rule{0.5\linewidth}{0.5pt}\end{center}

\subsection*{5. Validation hors de l'auteur du modèle}

\begin{enumerate}
\def\labelenumi{\arabic{enumi}.}
\tightlist
\item
  \textbf{Prédiction rétrospective hors échantillon à étanchéité
  temporelle stricte} : fixer une date de coupure \(T_0\) (par exemple
  par protocole : fin 2015). Construire vocabulaires, plongements,
  graphes relationnels et paramètres sur \(t \le T_0\) --- en
  interdisant tout grand modèle de langage pré-entraîné sur des textes
  postérieurs à \(T_0\). Prédire sur \([T_0+1, T_1]\) : (a) la naissance
  de nouvelles paires \((u,v)\) ; (b) parmi toutes les paires nées en
  \(T_0+1\), leur fécondité \(\mathcal{F}_{\text{obs}}\) à \(T_1\).
  Évaluer par l'aire sous la courbe précision-rappel (\textbf{AUPRC}),
  le \textbf{Recall@K} et la corrélation de rang de Spearman.
\item
  \textbf{Bases de comparaison simples} :

  \begin{itemize}
  \tightlist
  \item
    \emph{B1 (Topologie locale)} : attachement préférentiel et score
    d'Adamic-Adar dans le graphe de cooccurrence.
  \item
    \emph{B2 (Cascade non structurée, H2 seule)} : score ABC de Swanson
    pondéré et similarité cosinus \(x_u(T_0)^\top x_v(T_0)\).
  \item
    \emph{B3 (Atypicité bibliométrique)} : z-score de co-citation des
    journaux d'Uzzi et al.~(2013) sans couche relationnelle
    \(\mathcal{L}_{\text{struct}}\).
  \item
    \emph{B4 (Apprentissage profond sans invariant)} : prédiction de
    liens par R-GCN standard dépourvu de contrainte d'homomorphisme
    \(\mathcal{K}\) et du terme de champ \(F\).
  \end{itemize}
\item
  \textbf{Traitement du biais du survivant} : l'évaluation de la
  fécondité porte sur l'intégralité des mariages \emph{tentés} en
  \(T_0+1\) (incluant prépublications non publiées en revue, articles
  sans citation externe, essais cliniques négatifs et rétractations), et
  non sur un florilège de succès historiques. On sépare ainsi l'effet de
  \(\mathcal{K}\) sur la \emph{tentative} (\(\lambda_{uv}\)) de son
  effet sur la \emph{survie féconde} (\(\mathcal{F}_{\text{obs}}\)).
\item
  \textbf{Test sur les découvertes simultanées} : apparier chaque paire
  d'\emph{articles jumeaux} née dans \([T_0+1, T_1]\) à un ensemble de
  découvertes \emph{singletons} de même année, même discipline et même
  fréquence marginale des concepts parents \((u,v)\). Tester par
  régression logistique conditionnelle si \(F(u,v; T_0)\) prédit
  positivement le caractère multiple de la découverte à
  \(C_{\text{soc}}(u,v; T_0)\) contrôlé.
\item
  \textbf{Évaluation humaine en double aveugle} : extraire à \(T_0\) des
  hypothèses candidates rédigées sous un gabarit syntaxique uniforme
  (masquant auteurs et prestige institutionnel), réparties selon un plan
  factoriel croisant distance de surface (\(d_{\text{surf}}\) haute vs
  basse) et alignement (\(\mathcal{K}\) haut vs bas, plus un contrôle
  aléatoire). Des chercheurs du domaine cible évaluent à l'aveugle la
  \emph{surprise}, la \emph{plausibilité mécaniste} et la \emph{priorité
  d'investigation}, testant directement la question 1 (intérêt perçu)
  indépendamment du bruit bibliométrique.
\item
  \textbf{Préenregistrement et critères de réfutation} : archiver sur
  \textbf{OSF} avant toute inspection des données post-\(T_0\) le
  conteneur logiciel figé, la liste des hyperparamètres et les règles de
  décision. Le modèle est \textbf{réfuté} si :

  \begin{itemize}
  \tightlist
  \item
    \emph{(Réfutation du cœur H1--H3)} Sur la cohorte complète des
    mariages nés en \(T_0+1\), l'ajout de \(\mathcal{K}\) et de
    l'interaction \(d_{\text{surf}} \cdot \mathcal{K}\) ne produit aucun
    gain positif statistiquement significatif d'AUPRC par rapport à
    \(B1\text{--}B3\) pour prédire la fécondité, ou si les experts en
    aveugle préfèrent le quadrant (\(d_{\text{surf}}\) haute,
    \(\mathcal{K}\) bas) au quadrant (\(d_{\text{surf}}\) haute,
    \(\mathcal{K}\) haut).
  \item
    \emph{(Réfutation de H4)} Le terme quadratique estimé vérifie
    \(\hat{\beta}_2 \le 0\) (absence de zone intermédiaire optimale).
  \item
    \emph{(Réfutation de H5)} Le champ \(F(u,v; T_0)\) ne discrimine pas
    les découvertes simultanées des singletons appariés au-delà du
    hasard.
  \end{itemize}
\end{enumerate}

\begin{center}\rule{0.5\linewidth}{0.5pt}\end{center}

\subsection*{6. Ce que les hypothèses oublient}

\begin{enumerate}
\def\labelenumi{\arabic{enumi}.}
\tightlist
\item
  \textbf{Coût d'épreuve expérimentale et tractabilité empirique.} Une
  analogie dotée d'un score \(\mathcal{K}\) élevé demeure spéculative si
  le coût matériel et instrumental \(C_{\text{exp}}(w, t)\) d'une
  expérience cruciale excède les ressources disponibles (Stephan, 2012,
  \emph{How Economics Shapes Science}, Harvard University Press). Il
  faut pondérer la fécondité attendue par un facteur mesurable
  d'accessibilité expérimentale.
\item
  \textbf{Traction par les anomalies empiriques.} H1--H10 décrivent une
  dynamique poussée par la recombinaison de concepts existants, en
  omettant le rôle moteur des \emph{anomalies} --- échecs persistants
  d'un cadre établi face à une mesure précise (Kuhn, 1962, \emph{The
  Structure of Scientific Revolutions}, University of Chicago Press ;
  Lakatos, 1970, in \emph{Criticism and the Growth of Knowledge},
  p.~91--196). Une couche explicite de problèmes ouverts et de
  contradictions empiriques doit polariser la recherche d'analogies.
\item
  \textbf{Asymétries institutionnelles, aversion au risque et effet
  Matthieu.} À structure identique, l'adoption d'une combinaison
  atypique dépend du statut institutionnel des auteurs et subit un biais
  conservateur dans l'évaluation par les pairs et le financement
  (Merton, 1968, « The Matthew effect in science », \emph{Science},
  159(3810), 56--63 ; Wang, Veugelers, \& Stephan, 2017, « Bias against
  novelty in science », \emph{Research Policy}, 46(8), 1416--1436).
\item
  \textbf{Genèse de nouvelles abstractions.} H3 suppose que l'invariant
  abstrait préexiste dans \(\mathcal{R}_{\text{abs}}\). Or les
  révolutions conceptuelles majeures construisent des objets formels
  inédits (calcul infinitésimal, entropie informationnelle) qui
  enrichissent l'alphabet relationnel lui-même
  (\(\mathcal{R}_{\text{abs}}(t) \subsetneq \mathcal{R}_{\text{abs}}(t+1)\)
  ; Carey, 2009, \emph{The Origin of Concepts}, Oxford University
  Press).
\item
  \textbf{Régime de réfutation rapide.} Engendrer des mariages lointains
  est peu coûteux ; la valeur scientifique dépend d'un mécanisme de
  sélection négative éliminant rapidement les alignements
  sous-contraints avant qu'ils ne saturent l'attention collective
  (Campbell, 1960).
\end{enumerate}

\begin{center}\rule{0.5\linewidth}{0.5pt}\end{center}

\subsection*{7. Une première expérience peu coûteuse}

Cet essai, exécutable en quelques jours sur une station de travail
standard avec un budget quasi nul (données ouvertes uniquement), met
directement en péril \textbf{H1} et \textbf{H3} sans biais du survivant.

\begin{itemize}
\tightlist
\item
  \textbf{Périmètre et données} : croiser \textbf{SemMedDB},
  \textbf{MeSH} et \textbf{OpenAlex}. Extraire toutes les paires de
  concepts \((u,v)\) appartenant à deux arbres MeSH majeurs distincts
  qui apparaissent conjointement pour la première fois dans le titre ou
  le résumé d'un article PubMed au cours d'une année de référence fixée
  par protocole (\(T_{\text{naiss}} = 2018\)). Toutes les variables
  explicatives sont calculées exclusivement sur les publications
  \(\le 2017\).
\item
  \textbf{Variables (\(\le 2017\))} :

  \begin{enumerate}
  \def\labelenumi{\arabic{enumi}.}
  \tightlist
  \item
    \(d_{\text{surf}}(u, v; 2017)\) : combinaison de la divergence de
    Jensen-Shannon des revues et de la distance cosinus SGNS entraînée
    sur les résumés PubMed \(\le 2017\).
  \item
    Contrôle \textbf{H2} : score ABC de Swanson sur les intermédiaires
    partagés dans SemMedDB \(\le 2017\).
  \item
    Cœur \textbf{H3} : score \(\mathcal{K}(u, v; 2017)\) calculé par
    \textbf{Fused Gromov-Wasserstein} (\texttt{POT}) entre les
    sous-graphes causaux amont de rayon \(k=2\) (\emph{CAUSES,
    STIMULATES, INHIBITS, DISRUPTS, PREVENTS}), avec pénalité de conflit
    \(\chi\) dès qu'une relation excitatrice est alignée sur une
    relation inhibitrice.
  \end{enumerate}
\item
  \textbf{Variable cible (2019--2024)} : parmi toutes les paires
  effectivement nées en 2018, binaire de fécondité \(Y=1\) si la paire
  \((u,v)\) est reprise sur 2019--2024 par des équipes indépendantes
  (aucun coauteur commun avec l'article de 2018, seuil de protocole fixé
  avant analyse, par exemple au décile supérieur des reprises
  indépendantes), contre \(Y=0\) en l'absence totale de reprise
  indépendante.
\item
  \textbf{Comparaison critique} : évaluer par validation croisée
  stratifiée par couple de branches MeSH trois régressions logistiques
  régularisées :

  \begin{itemize}
  \tightlist
  \item
    \(M_0\) (Surface + H2) : fréquences marginales, attachement
    préférentiel, score ABC de Swanson, \(d_{\text{surf}}\) et
    \(d_{\text{surf}}^2\).
  \item
    \(M_1\) (H1 + H3) :
    \(M_0 + \mathcal{K} + d_{\text{surf}} \cdot \mathcal{K}\).
  \item
    \(M_{\text{perm}}\) (Contrôle de serrure-clé) : identique à \(M_1\),
    mais où \(\mathcal{K}\) est recalculé après permutation aléatoire
    des étiquettes causales (\emph{STIMULATES} \(\leftrightarrow\)
    \emph{INHIBITS}) à topologie de graphe inchangée.
  \end{itemize}
\item
  \textbf{Prédiction, seuil décisionnel (fixé \emph{a priori} par
  protocole) et lecture} :

  \begin{itemize}
  \tightlist
  \item
    \emph{Prédiction} :
    \(\hat{\beta}_{d_{\text{surf}} \cdot \mathcal{K}} > 0\) dans
    \(M_1\), et
    \(\text{AUPRC}(M_1) > \text{AUPRC}(M_{\text{perm}}) > \text{AUPRC}(M_0)\).
  \item
    \emph{Seuil d'acceptation a priori} : gain relatif d'AUPRC d'au
    moins \(5\,\%\) de \(M_1\) sur \(M_0\), avec borne inférieure de
    l'intervalle de confiance bootstrap à \(99\,\%\) de
    \(\text{AUPRC}(M_1) - \text{AUPRC}(M_{\text{perm}})\) strictement
    positive.
  \item
    \emph{Lecture} : si \(M_1\) ne bat pas \(M_0\), le cœur
    \textbf{H1+H3} est réfuté sur le savoir biomédical extrait ; si
    \(M_1\) bat \(M_0\) mais pas \(M_{\text{perm}}\), seule la densité
    du voisinage compte et l'hypothèse \textbf{H3} de conservation
    d'invariants sans conflits est réfutée ; si les deux seuils sont
    franchis, \textbf{H1} et \textbf{H3} survivent à l'épreuve et
    justifient le passage à l'échelle sur les corpus d'équations et le
    champ \(F(t)\).
  \end{itemize}
\end{itemize}
\clearpage
\section*{Annexe C — Réponse de Neon}\addcontentsline{toc}{subsection}{Annexe C — Réponse de Neon}
{\small Source : \texttt{reponses-google-v1/reponse-neon-v1.md}.}\par
\subsection*{1. Confrontation à la littérature}

Avant d'évaluer H1--H10, il faut dissocier trois questions que ces
hypothèses confondent : 1. \textbf{L'intérêt perçu par un public humain}
: propriété épistémique et sociale, relative aux croyances \emph{a
priori} d'une communauté à l'instant \(t\). Comme l'a montré Davis
(1971), une proposition est jugée intéressante lorsqu'elle contredit
certains présupposés faiblement tenus du public sans rejeter l'ensemble
de son cadre ; une idée fausse peut donc être très intéressante (ex. la
fusion froide à son annonce) et une idée vraie paraître triviale ou
incompréhensible. 2. \textbf{La vérité et la fécondité scientifique} :
adéquation prédictive au réel, cohérence explicative, capacité
d'unification et de génération de prédictions testables. 3. \textbf{Les
mécanismes de production} : processus cognitifs, algorithmiques et
sociaux par lesquels une hypothèse émerge.

Les dix hypothèses mélangent ces trois registres : H1, H4 et H10 visent
surtout l'intérêt et la génération ; H3, H6 et H7 touchent la validité
structurelle et la fécondité ; H2, H5, H8 et H9 portent sur la
production et la diffusion.

\begin{itemize}
\tightlist
\item
  \textbf{H1 (mariage éloigné) --- Soutenue (avec nuance forte).} En
  psychologie de l'analogie, Gentner (1983) montre que les analogies
  scientifiques fécondes reposent sur une dissimilarité de surface
  couplée à un isomorphisme relationnel profond. En bibliométrie, Uzzi
  et al.~(2013) et Shi \& Evans (2023) montrent que les articles à fort
  impact associent un socle très conventionnel à un petit nombre de
  combinaisons atypiques ; Foster et al.~(2015) et Wang et al.~(2017)
  confirment cependant que les sauts lointains ont une variance élevée,
  un fort risque d'échec et une reconnaissance retardée, et Kaplan \&
  Vakili (2015) montrent que la recombinaison locale produit aussi des
  percées. H1 est donc incomplète sans l'ancrage conventionnel.
\item
  \textbf{H2 (cascades et transferts) --- Soutenue pour la production,
  contredite comme critère suffisant de validité.} Le modèle ABC de
  découverte fondée sur la littérature de Swanson (1986) et Swanson \&
  Smalheiser (1997) prouve que des chaînes \(A\text{–}B\text{–}C\)
  relient des corpus disjoints et génèrent des hypothèses biologiques
  testables. Mais la transitivité naïve explose combinatoirement en faux
  positifs (polysémie du terme pivot \(B\), biais des nœuds très
  connectés). L'exemple \emph{douleur--anesthésie--conscience--LLM}
  illustre précisément ce piège : sans contraintes mécanistiques (H3),
  le transfert risque la fausse homologie.
\item
  \textbf{H3 (serrure et clé ; Noether et lagrangien) --- Soutenue pour
  l'alignement structurel ; métaphore sans contenu testable pour
  Noether/lagrangien.} Le principe de systématicité de Gentner (1983),
  l'algorithme SME (Falkenhainer et al., 1989) et la théorie de la
  cohérence explicative de Thagard (1989) soutiennent l'idée qu'une
  analogie viable aligne des réseaux entiers de relations en évitant les
  contradictions. En revanche, le théorème de Noether (1918) exige une
  symétrie continue différentiable d'une intégrale d'action gouvernant
  des trajectoires extrémales : sans principe variationnel mesurable
  dans un graphe conceptuel discret et dissipatif, la référence au
  lagrangien est purement métaphorique.
\item
  \textbf{H4 (ni trop près ni trop loin) --- Soutenue pour l'intérêt
  perçu ; partiellement soutenue pour la fécondité.} La courbe en U
  inversé entre nouveauté et valeur hédonique ou curiosité est bien
  établie (Berlyne, 1970 ; Loewenstein, 1994), tout comme l'équilibre
  entre familier et surprenant (Davis, 1971 ; Uzzi et al., 2013).
  Toutefois, pour la fécondité scientifique, des sauts locaux restent
  productifs (Kaplan \& Vakili, 2015) : le « trop près » est moins
  pénalisé en vérité qu'en intérêt perçu.
\item
  \textbf{H5 (champ, air du temps, ferromagnétisme, appariement
  aléatoire) --- Soutenue pour le rôle du contexte
  matériel/méthodologique dans les découvertes simultanées ; métaphore
  sans contenu testable pour le ferromagnétisme et les protéines.}
  Ogburn \& Thomas (1922) et Merton (1961) documentent l'omniprésence
  des découvertes multiples dès que les prérequis conceptuels et
  techniques sont réunis ; Galison (2003) confirme le rôle matériel des
  réseaux d'horloges dans la genèse de la relativité. Simonton (1979)
  montre que la distribution des multiples s'ajuste bien à un processus
  stochastique de Poisson (compatible avec l'idée de collisions
  aléatoires conditionnées par un contexte partagé). L'analogie
  ferromagnétique (brisure spontanée de symétrie) et celle de
  l'appariement protéique n'ajoutent aucune prédiction mesurable propre.
\item
  \textbf{H6 (connexion et distinction) --- Soutenue historiquement et
  cognitivement ; inconnue à grande échelle bibliométrique.} Hesse
  (1966) souligne le rôle heuristique des \emph{analogies négatives}
  (points de rupture). Wiser \& Carey (1983) montrent que le progrès
  conceptuel procède autant par différenciation de fausses identités
  (séparer chaleur et température) que par unification. En revanche, les
  modèles quantitatifs sur graphes scientifiques se concentrent presque
  exclusivement sur la recombinaison et testent très rarement la
  scission conceptuelle.
\item
  \textbf{H7 (profondeur : proximité aux hubs et unification
  rétrospective) --- Soutenue pour l'unification ; partiellement
  contredite pour la proximité aux hubs.} L'unification comme
  compression explicative est défendue par Kitcher (1981) et formalisée
  en théorie de l'information par le principe de longueur de description
  minimale, MDL (Grünwald, 2007). En revanche, Funk \& Owen-Smith (2017)
  et Wu et al.~(2019) montrent que les idées les plus disruptives
  \emph{éclipsent} les anciens hubs citationnels au lieu de s'y arrimer.
  En outre, l'impression de comprendre \emph{après coup} pourquoi l'idée
  surprenait est fortement contaminée par le biais rétrospectif
  (Fischhoff, 1975).
\item
  \textbf{H8 (émergence : transition de phase, percolation, surface de
  contact) --- Soutenue pour la surface de contact et les seuils de
  diffusion ; métaphore sans contenu testable pour la transition de
  phase thermodynamique stricte.} Burt (2004) montre que les courtiers
  placés aux interfaces entre communautés (« trous structuraux »)
  produisent des idées mieux évaluées. Bettencourt et al.~(2008)
  modélisent l'émergence de champs scientifiques par des dynamiques
  épidémiques à seuil. Cependant, les réseaux scientifiques finis et
  hétérogènes ne présentent pas d'exposants critiques thermodynamiques
  vérifiés.
\item
  \textbf{H9 (structure cachée du langage via méthodes non supervisées /
  ACP) --- Partiellement soutenue, partiellement contredite.} Tshitoyan
  et al.~(2019) montrent que des plongements lexicaux non supervisés en
  science des matériaux prédisent des composés fonctionnels avant leur
  découverte. Mais les vecteurs distributionnels échouent sur les
  analogies relationnelles complexes (Rogers et al., 2017 ; Mitchell,
  2021) car ils encodent d'abord la co-occurrence de surface ---
  exactement la dimension que H1 demande de dépasser.
\item
  \textbf{H10 (signal de surprise analogue dopaminergique) --- Soutenue
  pour les signaux d'information intrinsèque ; contredite pour la
  surprise brute.} En neurosciences, les neurones dopaminergiques codent
  surtout l'erreur de prédiction de récompense (Schultz et al., 1997),
  non la simple improbabilité. En apprentissage par renforcement, la
  surprise bayésienne (Itti \& Baldi, 2009) attire l'attention, mais
  maximiser la surprise brute échoue sur le bruit aléatoire ; il faut
  maximiser le \emph{progrès de compression} ou d'apprentissage (Oudeyer
  et al., 2007 ; Schmidhuber, 2010).
\end{itemize}

\textbf{Références de la partie 1 :} - Berlyne, D. E. (1970).
\emph{Novelty, complexity, and hedonic value}. Perception \&
Psychophysics, 8(5), 279--286. - Bettencourt, L. M. A., Kaiser, D. I.,
Kaur, J., Castillo-Chávez, C., \& Wojick, D. E. (2008). \emph{Population
modeling of the emergence and development of scientific fields}.
Scientometrics, 75(3), 495--518. - Burt, R. S. (2004). \emph{Structural
holes and good ideas}. American Journal of Sociology, 110(2), 349--399.
- Davis, M. S. (1971). \emph{That's interesting! Towards a phenomenology
of sociology and a sociology of phenomenology}. Philosophy of the Social
Sciences, 1(2), 309--344. - Falkenhainer, B., Forbus, K. D., \& Gentner,
D. (1989). \emph{The Structure-Mapping Engine: Algorithm and examples}.
Artificial Intelligence, 41(1), 1--63. - Fischhoff, B. (1975).
\emph{Hindsight \ensuremath{\neq} foresight: The effect of outcome knowledge on judgment
under uncertainty}. Journal of Experimental Psychology: Human Perception
and Performance, 1(3), 288--299. - Foster, J. G., Rzhetsky, A., \&
Evans, J. A. (2015). \emph{Tradition and innovation in scientists'
research strategies}. American Sociological Review, 80(5), 875--908. -
Funk, R. J., \& Owen-Smith, J. (2017). \emph{A dynamic network measure
of technological change}. Management Science, 63(3), 791--817. -
Galison, P. (2003). \emph{Einstein's Clocks, Poincaré's Maps: Empires of
Time}. New York : W. W. Norton. - Gentner, D. (1983).
\emph{Structure-mapping: A theoretical framework for analogy}. Cognitive
Science, 7(2), 155--170. - Grünwald, P. D. (2007). \emph{The Minimum
Description Length Principle}. Cambridge, MA : MIT Press. - Hesse, M. B.
(1966). \emph{Models and Analogies in Science}. Notre Dame : University
of Notre Dame Press. - Itti, L., \& Baldi, P. (2009). \emph{Bayesian
surprise attracts human attention}. Vision Research, 49(10), 1295--1306.
- Kaplan, S., \& Vakili, K. (2015). \emph{The double-edged sword of
recombination in breakthrough innovation}. Strategic Management Journal,
36(10) (pagination à vérifier). - Kitcher, P. (1981). \emph{Explanatory
unification}. Philosophy of Science, 48(4), 507--531. - Loewenstein, G.
(1994). \emph{The psychology of curiosity: A review and
reinterpretation}. Psychological Bulletin, 116(1), 75--98. - Merton, R.
K. (1961). \emph{Singletons and multiples in scientific discovery}.
Proceedings of the American Philosophical Society, 105(5), 470--486. -
Mitchell, M. (2021). \emph{Abstraction and analogy-making in artificial
intelligence}. Annals of the New York Academy of Sciences, 1505(1),
79--101. - Noether, E. (1918). \emph{Invariante Variationsprobleme}.
Nachrichten von der Gesellschaft der Wissenschaften zu Göttingen,
Mathematisch-Physikalische Klasse, 235--257. - Ogburn, W. F., \& Thomas,
D. (1922). \emph{Are inventions inevitable? A note on social evolution}.
Political Science Quarterly, 37(1), 83--98. - Oudeyer, P.-Y., Kaplan,
F., \& Hafner, V. V. (2007). \emph{Intrinsic motivation systems for
autonomous mental development}. IEEE Transactions on Evolutionary
Computation, 11(2), 265--286. - Rogers, A., Drozd, A., \& Li, B. (2017).
\emph{The (too many) problems of analogical reasoning with word
vectors}. Proceedings of \emph{SEM 2017 (pagination à vérifier). -
Schmidhuber, J. (2010). }Formal theory of creativity, fun, and intrinsic
motivation (1990--2010)\emph{. IEEE Transactions on Autonomous Mental
Development, 2(3), 230--247. - Schultz, W., Dayan, P., \& Montague, P.
R. (1997). }A neural substrate of prediction and reward\emph{. Science,
275(5306), 1593--1599. - Shi, F., \& Evans, J. (2023). }Surprising
combinations of research contents and contexts are related to impact and
emerge with scientific outsiders from distant disciplines\emph{. Nature
Communications, 14 (n° d'article à vérifier). - Simonton, D. K. (1979).
}Multiple discovery and invention: Zeitgeist, genius, or chance?*
Journal of Personality and Social Psychology, 37(9), 1603--1616. -
Swanson, D. R. (1986). \emph{Fish oil, Raynaud's syndrome, and
undiscovered public knowledge}. Perspectives in Biology and Medicine,
30(1), 7--18. - Swanson, D. R., \& Smalheiser, N. R. (1997). \emph{An
interactive system for finding complementary literatures}. Artificial
Intelligence, 91(2), 183--203. - Thagard, P. (1989). \emph{Explanatory
coherence}. Behavioral and Brain Sciences, 12(3), 435--467. - Tshitoyan,
V., et al.~(2019). \emph{Unsupervised word embeddings capture latent
knowledge from materials science literature}. Nature, 571, 95--98. -
Uzzi, B., Mukherjee, S., Stringer, M., \& Jones, B. (2013).
\emph{Atypical combinations and scientific impact}. Science, 342(6157),
468--472. - Wang, J., Veugelers, R., \& Stephan, P. (2017). \emph{Bias
against novelty in science}. Research Policy, 46(8), 1416--1436. -
Wiser, M., \& Carey, S. (1983). \emph{When heat and temperature were
one}. In D. Gentner \& A. L. Stevens (dir.), \emph{Mental Models}
(pp.~267--297). Hillsdale, NJ : Erlbaum. - Wu, L., Wang, D., \& Evans,
J. A. (2019). \emph{Large teams develop and small teams disrupt science
and technology}. Nature, 566, 378--382.

\begin{center}\rule{0.5\linewidth}{0.5pt}\end{center}

\subsection*{2. Modèle formel minimal}

\subsubsection*{2.1 Espace des concepts}

À chaque date discrète \(t \in \mathbb{N}\) (trimestre ou année), l'état
des connaissances est un graphe multicouche daté
\(\mathcal{G}_t = \big(V_t, E_t^{(U)}, E_t^{(S)}, E_t^{(C)}, \mathcal{A}_t\big)\)
construit \textbf{strictement} sur les documents publiés jusqu'à \(t\) :
- \(V_t\) : ensemble des concepts actifs (entités, phénomènes,
opérateurs). - \textbf{Couche d'usage} \(E_t^{(U)}\) : graphe pondéré de
co-occurrence textuelle muni de plongements lexicaux
\(u_i(t) \in \mathbb{R}^{d_U}\). - \textbf{Couche structurelle}
\(E_t^{(S)} \subseteq V_t \times \mathcal{R} \times V_t\) : multigraphe
dirigé et typé, où \(\mathcal{R}\) est un vocabulaire fini de relations
\emph{indépendantes du domaine} (ex. \emph{cause}, \emph{inhibe},
\emph{conserve}, \emph{rétroagit négativement sur}, \emph{dérive de}),
muni d'une table d'incompatibilités \((r, \neg r)\). On note
\(G_i(t) = N_k^{(S)}(i,t)\) le sous-graphe relationnel local de rayon
\(k\) (typiquement \(k=2\)) autour de \(i\). - \textbf{Couche
communautaire} \(E_t^{(C)}\) : graphe chercheurs--concepts induisant,
pour chaque concept \(i\), une distribution d'appartenance disciplinaire
\(c_i(t) \in \Delta^{M-1}\) sur \(M\) communautés et une surface de
contact inter-communautés \(B_{ij}(t) \ge 0\) (densité de chercheurs ou
de chemins citationnels à deux sauts reliant les communautés portant
\(i\) et \(j\)). - \(\mathcal{A}_t\) : distribution probabiliste de
référence modélisant les attentes du corpus jusqu'à \(t\).

\subsubsection*{2.2 Deux distances comparables d'un domaine à l'autre}

\begin{enumerate}
\def\labelenumi{\arabic{enumi}.}
\tightlist
\item
  \textbf{Distance de surface} \(d_{\mathrm{surf}}(i,j;t) \in [0,1]\) :
  \[d_{\mathrm{surf}}(i,j;t) = \lambda_1 \frac{1 - \cos\!\big(u_i(t), u_j(t)\big)}{2} + \lambda_2 \, d_{\mathrm{JS}}\!\big(c_i(t), c_j(t)\big) + \lambda_3 \, \frac{\mathrm{dist}_{E_t^{(U)}}(i,j)}{D_t}\]
  avec \(\lambda_1+\lambda_2+\lambda_3 = 1\), \(d_{\mathrm{JS}}\) la
  distance de Jensen--Shannon entre profils communautaires, et \(D_t\)
  le diamètre effectif de \(E_t^{(U)}\).
\item
  \textbf{Distance structurelle inter-domaines}
  \(d_{\mathrm{struct}}(i,j;t) \in [0,1]\) : pour comparer deux domaines
  sans fuite lexicale, on supprime tout nom de nœud dans \(G_i(t)\) et
  \(G_j(t)\) et l'on ne conserve que les tenseurs d'adjacence
  relationnels \(R_i(a,a'), R_j(b,b') \in \{0,1\}^{|\mathcal{R}|}\). La
  distance est donnée par le transport optimal relationnel de
  Gromov--Wasserstein (enraciné en \(\pi_{ij}>0\)) :
  \[d_{\mathrm{struct}}(i,j;t) = \min_{\pi \in \Pi(\mu_i,\mu_j),\, \pi_{ij} > 0} \sum_{a,a' \in G_i} \sum_{b,b' \in G_j} \ell\!\big(R_i(a,a'), R_j(b,b')\big)\, \pi_{ab}\,\pi_{a'b'}\]
  où \(\ell\) est une perte normalisée sur le type de relation
  uniquement. Deux systèmes à vocabulaires disjoints partageant un même
  motif causal ou différentiel obtiennent
  \(d_{\mathrm{struct}} \approx 0\) malgré
  \(d_{\mathrm{surf}} \approx 1\).
\end{enumerate}

\subsubsection*{2.3 Critère serrure--clé (H3 et H6)}

Pour une bijection partielle
\(\phi : U_i \subseteq V(G_i) \to U_j \subseteq V(G_j)\) telle que
\(\phi(i)=j\) (extraite du couplage optimal \(\pi^*\) ou par alignement
discret), on définit trois grandeurs mesurables : - \textbf{Préservation
relationnelle} :
\[\mathrm{Pres}(\phi) = \frac{|\{(a,r,a') \in E(G_i|_{U_i}) : (\phi(a), r, \phi(a')) \in E(G_j)\}|}{|E(G_i|_{U_i})|}\]
- \textbf{Couverture} (évitant les alignements triviaux d'un nœud isolé)
: \[\mathrm{Cov}(\phi) = \frac{|U_i| + |U_j|}{|V(G_i)| + |V(G_j)|}\] -
\textbf{Taux de conflit} (incompatibilités relationnelles explicites) :
\[\mathrm{Conf}(\phi) = \frac{|\{(a,r,a') \in E(G_i|_{U_i}) : (\phi(a), \neg r, \phi(a')) \in E(G_j)\}|}{|E(G_i|_{U_i})|}\]
Le \textbf{score serrure--clé} vaut :
\[K(i,j;t) = \max_{\phi:\,\phi(i)=j} \Big[\mathrm{Pres}(\phi)^{\alpha}\,\mathrm{Cov}(\phi)^{\beta}\,\big(1 - \mathrm{Conf}(\phi)\big)^{\gamma}\Big] \in [0,1]\]
Symétriquement, H6 (tri des fausses identités) donne un \textbf{score de
différenciation} pour deux concepts proches en surface mais
structurellement conflictuels :
\(S_{\mathrm{split}}(i,j;t) = \big(1 - d_{\mathrm{surf}}(i,j;t)\big)\,\mathrm{Conf}(\phi^*)\).

\subsubsection*{2.4 Champ dépendant du temps (H5)}

On associe à un ensemble de \(d_M\) instruments, méthodes ou cadres
formels un vecteur de saillance exogène \(h(t) \in \mathbb{R}_+^{d_M}\)
(intensité de poussée temporelle de chaque outil à \(t\)). Chaque
concept \(i\) possède une susceptibilité méthodologique
\(m_i(t) \in \mathbb{R}_+^{d_M}\) (co-occurrence ou applicabilité de
l'outil \(q\) au concept \(i\)). Le champ induit une
\textbf{co-exposition projective} :
\[\Phi_t(i,j) = m_i(t)^\top \mathrm{diag}\!\big(h(t)\big)\, m_j(t)\]
L'arrivée d'une technologie transversale (horloges synchronisées,
cryo-microscopie, transformateurs) fait croître \(h_q(t)\), rapprochant
brusquement des couples \((i,j)\) jusque-là disjoints.

\subsubsection*{2.5 Probabilité de naissance d'une idée}

Pour une paire inédite \((i,j) \notin E_t^{(U)}\), la naissance d'un
lien dans \([t, t+\Delta t]\) suit un \textbf{processus de Poisson
inhomogène} de taux d'intensité :
\[\lambda_{ij}(t) = \lambda_0 \, B_{ij}(t)^{\rho} \, \exp\!\big(\eta\,\Phi_t(i,j)\big) \, \exp\!\left(-\frac{\big(d_{\mathrm{surf}}(i,j;t) - d^*\big)^2}{2\sigma^2}\right) \, K(i,j;t)^{\theta}\]
La probabilité d'apparition vaut
\(P(N_{ij}(t) \ge 1) = 1 - e^{-\lambda_{ij}(t)\Delta t}\), et celle
d'une \textbf{découverte simultanée indépendante} dans une fenêtre
\(\Delta t\) avant diffusion vaut
\(P(N_{ij}(t) \ge 2) = 1 - \big(1 + \lambda_{ij}(t)\Delta t\big)e^{-\lambda_{ij}(t)\Delta t}\).

\subsubsection*{2.6 Intérêt perçu et fécondité de l'idée fille}

Conformément à la Partie 1, l'idée fille \(z = (i \bowtie_{\phi^*} j)\)
reçoit \textbf{deux scores distincts} : - \textbf{Intérêt perçu à \(t\)}
(surprise bayésienne tempérée par la recevabilité structurelle et la
courbe en cloche) :
\[I_{\mathrm{perc}}(z;t) = D_{\mathrm{KL}}\!\big(P(\cdot \mid \mathcal{A}_t \cup \{z\}) \,\|\, P(\cdot \mid \mathcal{A}_t)\big) \cdot K(i,j;t) \cdot \exp\!\left(-\frac{(d_{\mathrm{surf}}(i,j;t) - d^*)^2}{2\sigma^2}\right)\]
- \textbf{Fécondité unificatrice potentielle} (compression MDL de la
couche structurelle + nombre de prédictions transférables non encore
vérifiées, soit l'analogie neutre de Hesse) :
\[F(z;t) = \Delta L_{\mathrm{MDL}}\!\big(E_t^{(S)} \mid z\big) + \omega\,\Big|\big\{(a,r,a') \in E(G_i|_{U_i}) : (\phi^*(a), r, \phi^*(a')) \notin E(G_j) \text{ et non en conflit}\big\}\Big|\]

\subsubsection*{2.7 Dynamique d'émergence (H8)}

L'adoption de \(z\) dans la couche humaine \(E_t^{(C)}\) est modélisée
par une percolation par liens (ou processus de branchement) sur le
sous-graphe des communautés en contact. Avec une probabilité moyenne
d'adoption \(p_{\mathrm{adopt}}(z,t)\) (fonction croissante de
\(I_{\mathrm{perc}}\) puis des premiers résultats empiriques) et la
distribution des degrés inter-communautés \(k_{\mathrm{inter}}\), le
critère de formation d'une composante géante d'adoptants (critère de
Molloy--Reed ; Newman, 2010) est :
\[R_{\mathrm{eff}}(z,t) = p_{\mathrm{adopt}}(z,t)\,\frac{\langle k_{\mathrm{inter}}^2 \rangle - \langle k_{\mathrm{inter}} \rangle}{\langle k_{\mathrm{inter}} \rangle} > 1\]

\subsubsection*{2.8 Sort exact des analogies physiques}

\begin{longtable}[]{@{}
  >{\raggedright\arraybackslash}p{(\columnwidth - 4\tabcolsep) * \real{0.3333}}
  >{\raggedright\arraybackslash}p{(\columnwidth - 4\tabcolsep) * \real{0.3333}}
  >{\raggedright\arraybackslash}p{(\columnwidth - 4\tabcolsep) * \real{0.3333}}@{}}
\toprule\noalign{}
\begin{minipage}[b]{\linewidth}\raggedright
Analogie
\end{minipage} & \begin{minipage}[b]{\linewidth}\raggedright
Décision
\end{minipage} & \begin{minipage}[b]{\linewidth}\raggedright
Équivalent mathématique mesurable ou motif d'écartement
\end{minipage} \\
\midrule\noalign{}
\endhead
\bottomrule\noalign{}
\endlastfoot
\textbf{Puits (de potentiel)} & \textbf{Traduite} & Minimum local d'une
énergie de cohérence \(E(\mathbf{s}) = -\sum w_{ab}s_a s_b\) (Thagard,
1989) ou cluster à faible conductance dans \(E_t^{(U)}\) mesurant le
taux de rétention des chercheurs. \\
\textbf{Champ (et ferromagnétisme)} & \textbf{Champ traduit ;
ferromagnétisme écarté} & Champ traduit par la covariable temporelle
exogène \(h(t)\) dans \(\Phi_t(i,j)\). Le ferromagnétisme (aimantation
spontanée sans champ externe) est écarté : aucune symétrie de spin
n'existe dans \(V_t\). \\
\textbf{Polarisation} & \textbf{Traduite (sens restreint)} & Alignement
directionnel des vecteurs de susceptibilité \(m_i(t)\) sur l'axe
dominant \(h(t)\), mesuré par \(\cos(m_i(t), h(t))\), ou modularité
signée de \(E_t^{(C)}\). Tout sens électromagnétique est écarté. \\
\textbf{Résonance} & \textbf{Écartée} & Aucun oscillateur ni fréquence
propre dans un graphe conceptuel ; l'effet de sélectivité est
entièrement absorbé par le pic de \(K(i,j;t)\). \\
\textbf{Noether et lagrangien} & \textbf{Écartés} & Aucun espace de
configuration lisse ni principe variationnel \(\delta\mathcal{S}=0\).
L'invariance utile sous changement de domaine est formalisée directement
par les morphismes partiels de graphes
(\(\mathrm{Pres}, \mathrm{Cov}, \mathrm{Conf}\)). \\
\textbf{Percolation} & \textbf{Traduite (réseaux finis)} & Transition de
composante géante \(S_1(t)/|V_t|\) autour du seuil
\(R_{\mathrm{eff}} > 1\) sur \(E_t^{(C)}\). Les exposants critiques
thermodynamiques en taille infinie sont écartés. \\
\textbf{Appariement de protéines} & \textbf{Écarté} & Réduit sans perte
au taux stochastique de Poisson \(\lambda_{ij}(t)\) conditionné par le
contact \(B_{ij}\) et la compatibilité géométrique \(K(i,j)\). \\
\end{longtable}

\emph{Référence ajoutée :} Newman, M. E. J. (2010). \emph{Networks: An
Introduction}. Oxford : Oxford University Press.

\begin{center}\rule{0.5\linewidth}{0.5pt}\end{center}

\subsection*{3. Données brutes et extraction des variables}

Le modèle exige des sources intégralement datées pour garantir la
causalité temporelle :

\begin{enumerate}
\def\labelenumi{\arabic{enumi}.}
\tightlist
\item
  \textbf{Corpus datés en texte intégral} : \textbf{PubMed Central Open
  Access (PMC OA)}, \textbf{arXiv} (sources LaTeX datées au jour près),
  \textbf{Semantic Scholar Open Research Corpus (S2ORC)} et
  \textbf{OpenAlex}. \emph{Variables tirées} : vocabulaire \(V_t\),
  couche d'usage \(E_t^{(U)}\), plongements diachroniques \(u_i(t)\),
  vecteur méthodologique \(m_i(t)\), champ \(h(t)\) (fréquences et
  poussées de méthodes/instruments), et corpus \(\mathcal{A}_t\) pour la
  surprise bayésienne.
\item
  \textbf{Graphes de citations et de coauteurs} : \textbf{OpenAlex},
  \textbf{Crossref}, \textbf{INSPIRE-HEP} / \textbf{NASA ADS}.
  \emph{Variables tirées} : couche communautaire \(E_t^{(C)}\), profils
  disciplinaires \(c_i(t)\), distance \(d_{\mathrm{JS}}\), surface de
  contact \(B_{ij}(t)\), distribution \(k_{\mathrm{inter}}\), et
  indicateurs avals (citations normalisées, indice de disruption
  \(CD_5\), taille de composante géante \(S_1(t)\)).
\item
  \textbf{Ontologies et bases relationnelles structurées} :
  \textbf{SemMedDB} (prédications biomédicales typées et datées
  extraites de PubMed ; Kilicoglu et al., 2012), \textbf{UMLS},
  \textbf{Gene Ontology}, \textbf{Wikidata}. \emph{Variables tirées} :
  couche structurelle \(E_t^{(S)}\), vocabulaire \(\mathcal{R}\) et
  table d'incompatibilités \((r, \neg r)\) (ex. \emph{STIMULATES} vs
  \emph{INHIBITS}, \emph{CAUSES} vs \emph{PREVENTS}).
\item
  \textbf{Corpus d'équations} : équations extraites des sources LaTeX
  d'\textbf{arXiv} (\textbf{arXMLiv}). \emph{Variables tirées} : couche
  structurelle \(E_t^{(S)}\) à très haute précision en physique,
  mathématiques appliquées et biologie quantitative, par arbres
  d'opérateurs anonymisés (invariants au renommage des variables
  physiques).
\item
  \textbf{Catalogues de découvertes et de multiples} : corpus
  historiques d'Ogburn \& Thomas (1922) et de Simonton (1979) pour la
  validation qualitative ; surtout, \textbf{multiples modernes
  objectifs} extraits automatiquement d'arXiv et bioRxiv (dépôts
  indépendants introduisant la même paire \((i,j)\) dans une fenêtre
  étroite, sans coauteurs communs ni citations mutuelles).
  \emph{Variables tirées} : comptages observés \(N_{ij}(t)\) pour tester
  \(\lambda_{ij}(t)\) et \(\Phi_t(i,j)\).
\item
  \textbf{Catalogues d'analogies échouées et de pistes stériles} : pour
  corriger le biais du survivant, on mobilise (i) l'ensemble exhaustif
  des mariages nouveaux publiés restés sans suite (zéro ou une citation
  à cinq ans), (ii) les préprints arXiv/bioRxiv jamais publiés en revue,
  (iii) la base \textbf{Retraction Watch}, (iv) les essais cliniques
  négatifs ou interrompus (\textbf{ClinicalTrials.gov}) pour tester les
  transferts biomédicaux (H2), et (v) un recueil de contrôle d'analogies
  historiques réfutées (ex. solide élastique pour l'éther luminifère).
  \emph{Variables tirées} : exemples négatifs réalistes permettant de
  dissocier naissance, intérêt initial et fécondité.
\end{enumerate}

\emph{Référence ajoutée :} Kilicoglu, H., Shin, D., Fiszman, M.,
Rosemblat, G., \& Rindflesch, T. C. (2012). \emph{SemMedDB: A
PubMed-scale repository of biomedical semantic predications}.
Bioinformatics, 28(23), 3158--3160.

\begin{center}\rule{0.5\linewidth}{0.5pt}\end{center}

\subsection*{4. Méthodes et outils}

Chaque variable et chaque test correspond à un algorithme nommé,
entraîné \textbf{exclusivement} sur les données antérieures ou égales à
\(t\) :

\begin{itemize}
\tightlist
\item
  \textbf{Normalisation de \(V_t\) et couche d'usage \(E_t^{(U)}\)} :
  extraction d'entités par dictionnaires contrôlés (UMLS/MeSH, OpenAlex
  Concepts) et analyse syntaxique ; matrice \textbf{PPMI} factorisée par
  \textbf{SVD tronquée} (ou \textbf{Skip-gram with Negative Sampling}
  aligné d'une période à l'autre par \textbf{Procrustes orthogonal} ;
  Hamilton et al., 2016) pour \(u_i(t)\). Aucun modèle de langage
  pré-entraîné sur le futur ne peut fournir \(u_i(t)\). Distance
  géodésique dans \(E_t^{(U)}\) par \textbf{Dijkstra} (ou \textbf{BFS}
  sur graphes non pondérés).
\item
  \textbf{Couche communautaire \(E_t^{(C)}\)} : partitionnement des
  coauteurs et citations par l'algorithme de \textbf{Leiden} (Traag et
  al., 2019) ou \textbf{Stochastic Block Model hiérarchique}
  (\textbf{graph-tool}) ; calcul de \(d_{\mathrm{JS}}\) et de la
  contrainte structurale de \textbf{Burt} / nombre de chemins
  inter-communautés à deux sauts pour \(B_{ij}(t)\).
\item
  \textbf{Couche structurelle \(E_t^{(S)}\)} : extraction directe depuis
  \textbf{SemMedDB} (\textbf{SemRep}) en biomédecine ; pour les
  équations d'arXiv, conversion par \textbf{LaTeXML} et normalisation
  symbolique en arbres d'opérateurs anonymisés via \textbf{SymPy} ; en
  texte général, extraction locale phrase-par-phrase de triplets typés
  (ex. \textbf{OpenIE} ou parseur relationnel supervisé sur schémas
  fixes), auditée sur un échantillon annoté manuellement.
\item
  \textbf{Distance structurelle \(d_{\mathrm{struct}}\) et score
  serrure--clé \(K(i,j;t)\)} :

  \begin{itemize}
  \tightlist
  \item
    Criblage rapide par \textbf{Fused Gromov--Wasserstein} (Titouan et
    al., 2019) résolu par gradient conditionnel (\textbf{Frank--Wolfe}
    avec projections de \textbf{Sinkhorn--Knopp}, bibliothèque
    \textbf{POT}).
  \item
    Sur les paires candidates, alignement discret exact par
    \textbf{Structure-Mapping Engine (SME)} (Falkenhainer et al., 1989)
    ou recherche du plus grand sous-graphe commun étiqueté
    (\textbf{algorithme McSplit}), donnant directement
    \(\mathrm{Pres}\), \(\mathrm{Cov}\) et \(\mathrm{Conf}\).
  \item
    Pour les arbres d'équations, distance d'édition d'arbres par
    l'algorithme \textbf{APTED}.
  \end{itemize}
\item
  \textbf{Champ \(h(t)\)} : détection de poussées temporelles sur les
  méthodes et instruments par l'automate de \textbf{Kleinberg (2003)}
  (\emph{burst detection}), combinée à une \textbf{factorisation en
  matrices non négatives (NMF)} pour extraire des cadres méthodologiques
  latents.
\item
  \textbf{Surprise bayésienne et compression MDL} : modèle génératif
  léger entraîné uniquement sur \([t_0, t]\) (plongement de graphe de
  connaissances \textbf{RotatE} + modèle de langage causal petit format
  entraîné \emph{from scratch} jusqu'à \(t\)) pour estimer la
  log-surprise et \(D_{\mathrm{KL}}\) après une étape de mise à jour
  variationnelle ; compression structurelle \(\Delta L_{\mathrm{MDL}}\)
  estimée par découverte de motifs compressifs (\textbf{SUBDUE}).
\item
  \textbf{Estimation et tests statistiques} :

  \begin{itemize}
  \tightlist
  \item
    Taux de naissance \(\lambda_{ij}(t)\) : \textbf{modèle additif
    généralisé (GAM)} binomial négatif (pour surdispersion) ou
    \textbf{modèle de Cox} à covariables dépendantes du temps
    (\textbf{lifelines} / \textbf{mgcv}), avec splines lissées sur
    \(d_{\mathrm{surf}}\) afin de tester H4 sans imposer de forme
    gaussienne \emph{a priori}.
  \item
    Comparaison de modèles nichés par \textbf{test du rapport de
    vraisemblance}, \textbf{AIC/BIC} et \textbf{bootstrap en blocs} par
    concept.
  \item
    Dynamique d'adoption et \(R_{\mathrm{eff}}\) : \textbf{processus de
    Hawkes} multivarié sur les dates d'adoption et suivi empirique de la
    plus grande composante connexe.
  \end{itemize}
\end{itemize}

\emph{Références ajoutées :}

\begin{itemize}
\tightlist
\item
  Hamilton, W. L., Leskovec, J., \& Jurafsky, D. (2016).
  \emph{Diachronic word embeddings reveal statistical laws of semantic
  change}. Proceedings of ACL, 1489--1501.
\item
  Kleinberg, J. (2003). \emph{Bursty and hierarchical structure in
  streams}. Data Mining and Knowledge Discovery, 7(4), 373--397.
\item
  Titouan, V., Courty, N., Tavenard, R., Chapel, L., \& Flamary, R.
  (2019). \emph{Optimal transport for structured data with application
  on graphs}. Proceedings of ICML (PMLR 97).
\item
  Traag, V. A., Waltman, L., \& van Eck, N. J. (2019). \emph{From
  Louvain to Leiden: Guaranteeing well-connected communities}.
  Scientific Reports, 9, 5233.
\end{itemize}

\begin{center}\rule{0.5\linewidth}{0.5pt}\end{center}

\subsection*{5. Validation hors de l'auteur du modèle}

\begin{enumerate}
\def\labelenumi{\arabic{enumi}.}
\tightlist
\item
  \textbf{Prédiction rétrospective strictement causale contre des bases
  simples.} On entraîne et calibre tous les hyperparamètres sur une
  fenêtre passée \([t_0, T_1]\) (ex. 1990--2010), on gèle le modèle,
  puis on prédit sur une période future tenue à l'écart \([T_1+1, T_2]\)
  (ex. 2011--2020) : (i) l'apparition de liens nouveaux \((i,j)\) et
  (ii) la fécondité ultérieure (top décile des citations normalisées ou
  de l'indice de disruption \(CD_5\) à cinq ans). Aucune ressource
  entraînée après \(T_1\) n'est autorisée. Le modèle complet doit
  surpasser cinq bases simples : (B1) attachement préférentiel
  \(\deg(i)\deg(j)\) ; (B2) proximité de surface seule (Adamic--Adar et
  cosinus \(u_i, u_j\)) ; (B3) modèle ABC de Swanson pondéré sans
  critère serrure--clé ; (B4) courbe en cloche sur \(d_{\mathrm{surf}}\)
  seule (H1+H4 sans H3) ; (B5) \textbf{Node2Vec} + classifieur sur la
  seule couche d'usage. Vu le fort déséquilibre des classes, on évalue
  par \textbf{PR-AUC}, \textbf{Precision@K} et \textbf{Brier score}
  (calibration).
\item
  \textbf{Traitement du biais du survivant.} On bannit les palmarès
  rétrospectifs de « grandes idées » comme unique jeu de test. On
  modélise explicitement deux étapes (à la manière d'une correction de
  sélection) : \emph{Étape 1}, naissance/proposition d'une combinaison
  sur l'ensemble des paires possibles ; \emph{Étape 2}, fécondité
  (reprise par des équipes indépendantes, réplication, absence de
  rétractation, succès clinique le cas échéant)
  \emph{conditionnellement} à la publication, en incluant tous les
  articles publiés à zéro citation et les préprints abandonnés. On
  audite également les « faux positifs » à haut score \(K\) jamais
  publiés : s'ils s'avèrent trivialement faux à l'examen humain, le
  modèle est pris en défaut.
\item
  \textbf{Test sur les découvertes simultanées.} Sur un échantillon
  cas--témoins issu d'arXiv et bioRxiv, on compare chaque multiple
  indépendant (deux équipes sans coauteurs ni citations mutuelles dans
  une fenêtre courte) à des découvertes uniques (\emph{singletons})
  appariées en époque, discipline et distance \(d_{\mathrm{surf}}\). H5
  prédit un saut significativement plus élevé de la variation du champ
  \(\Delta \Phi_t(i,j)\) dans les mois précédant le multiple, ainsi
  qu'un ajustement de la distribution empirique de la multiplicité
  \(N_{ij}\) à la loi de Poisson prédite par
  \(\lambda_{ij}(t)\Delta t\).
\item
  \textbf{Évaluation humaine en aveugle.} Un protocole neutre et
  standardisé traduit en fiches anonymisées de longueur uniforme un
  échantillon stratifié de combinaisons inédites (haut score complet,
  haut score surface seule, haut score structure seule, aléatoires). Des
  experts indépendants recrutés hors de l'équipe notent chaque fiche en
  double aveugle sur quatre échelles séparées : \emph{surprise},
  \emph{intérêt}, \emph{plausibilité scientifique} et \emph{priorité de
  test} (accord inter-juges mesuré par l'\(\alpha\) de Krippendorff). Le
  modèle prédit une dissociation : \(I_{\mathrm{perc}}\) doit mieux
  corréler avec la surprise et l'intérêt, tandis que \(K\) et \(F\)
  doivent mieux corréler avec la plausibilité et la priorité
  expérimentale.
\item
  \textbf{Préenregistrement.} Avant d'accéder aux données de test
  \([T_1+1, T_2]\) et avant tout jugement humain, le code d'extraction,
  le dictionnaire \(\mathcal{R}\), les poids gelés, les baselines, les
  métriques et les seuils de réfutation sont déposés avec horodatage sur
  \textbf{OSF (Open Science Framework)}, accompagnés d'une liste
  prospective de prédictions pour les années à venir.
\item
  \textbf{Résultats qui réfuteraient le modèle.} Le modèle est réfuté si
  l'un des résultats suivants survient sur \([T_1+1, T_2]\) :

  \begin{itemize}
  \tightlist
  \item
    \textbf{Réfutation du cœur (H1+H3)} : le modèle complet n'apporte
    aucun gain statistiquement significatif en PR-AUC sur la meilleure
    baseline sans couche structurelle, ou le paramètre estimé
    \(\hat{\theta}\) de \(K(i,j;t)\) est nul ou négatif.
  \item
    \textbf{Réfutation de H4} : la relation estimée par GAM entre
    \(d_{\mathrm{surf}}\) et naissance/fécondité est strictement
    monotone décroissante (pas d'optimum intermédiaire).
  \item
    \textbf{Réfutation de H5} : \(\Delta \Phi_t(i,j)\) ne distingue pas
    les multiples des \emph{singletons} appariés.
  \item
    \textbf{Réfutation de H3 en aveugle} : \(K(i,j;t)\) ne prédit pas
    mieux la plausibilité jugée par les experts que la simple surprise
    lexicale.
  \end{itemize}
\end{enumerate}

\begin{center}\rule{0.5\linewidth}{0.5pt}\end{center}

\subsection*{6. Ce que les hypothèses oublient}

\begin{enumerate}
\def\labelenumi{\arabic{enumi}.}
\tightlist
\item
  \textbf{La résistance empirique du réel et les anomalies.} H1--H10
  décrivent la découverte comme une combinatoire interne au réseau
  conceptuel. Or une grande part des percées naît d'anomalies
  expérimentales qui fissurent le réseau de l'extérieur (Kuhn, 1962) ;
  une analogie parfaitement cohérente en interne (haut score \(K\)) peut
  être empiriquement fausse. Il faut ajouter une \textbf{couche de
  données et de résidus prédictifs} (anomalies non expliquées).
\item
  \textbf{La création de concepts primitifs et le changement de
  représentation.} Le modèle suppose que les parents \(i, j\) existent
  déjà dans \(V_t\). Or des idées majeures inventent de nouveaux objets
  ou espaces d'états absents des parents (entropie, spin, espace de
  Hilbert). Il manque un opérateur de \textbf{synthèse de variables
  latentes}.
\item
  \textbf{Les coûts matériels et la valeur expérimentale de
  l'information.} Une idée féconde sur le papier mais intestable faute
  d'instrument ou de budget reste stérile. Le choix des chercheurs obéit
  aussi à un ratio bénéfice épistémique / coût expérimental.
\item
  \textbf{Les incitations institutionnelles et biais sociaux.}
  Financements, effets de prestige (effet Matthieu ; Merton, 1968),
  conservatisme de l'évaluation par les pairs (Wang et al., 2017) et
  barrières linguistiques filtrent fortement ce qui émerge et survit.
  Sans covariables institutionnelles, le modèle confondra absence d'idée
  et rejet social.
\item
  \textbf{Les découvertes non combinatoires.} Contre-exemples
  destructeurs, théorèmes d'impossibilité (Gödel, Bell, Arrow),
  simplifications axiomatiques et approfondissements verticaux ne
  relèvent pas d'un mariage entre domaines éloignés.
\item
  \textbf{Les délais de maturation (\emph{Sleeping Beauties}).} Une idée
  vraie et féconde peut rester ignorée durant des décennies avant son «
  réveil » (Ke et al., 2015) : toute validation à fenêtre courte
  sous-estime ces cas.
\item
  \textbf{La causalité interventionnelle.} Un isomorphisme de graphe
  relationnel ne garantit pas l'invariance d'un modèle causal structurel
  (Pearl, 2009) : deux systèmes peuvent partager une topologie
  corrélationnelle sans partager de mécanisme transférable (point
  critique pour le transfert neurosciences--LLM évoqué en H2).
\end{enumerate}

\emph{Références ajoutées :}

\begin{itemize}
\tightlist
\item
  Ke, Q., Ferrara, E., Radicchi, F., \& Flammini, A. (2015).
  \emph{Defining and identifying Sleeping Beauties in science}. PNAS,
  112(24), 7426--7431.
\item
  Kuhn, T. S. (1962). \emph{The Structure of Scientific Revolutions}.
  Chicago : University of Chicago Press.
\item
  Merton, R. K. (1968). \emph{The Matthew Effect in Science}. Science,
  159(3810), 56--63.
\item
  Pearl, J. (2009). \emph{Causality: Models, Reasoning, and Inference}
  (2e éd.). Cambridge : Cambridge University Press.
\end{itemize}

\begin{center}\rule{0.5\linewidth}{0.5pt}\end{center}

\subsection*{7. Une première expérience peu coûteuse pour réfuter le cœur
(H1 et H3)}

\textbf{Objectif.} Tester en quatre à cinq jours, pour un budget
inférieur à quelques centaines d'euros (une station de travail standard
ou quelques heures de GPU cloud, données 100 \% publiques), si
l'interaction entre distance de surface (H1) et alignement serrure--clé
(H3) prédit mieux la fécondité des nouvelles combinaisons scientifiques
que la surface et les chemins ABC seuls.

\textbf{Protocole (zéro annotation préalable).} On utilise
\textbf{SemMedDB} (couche structurelle \(E_t^{(S)}\) déjà typée et datée
avec prédicats et paires d'antonymes) couplée à \textbf{PubMed /
OpenAlex} (couches \(E_t^{(U)}\) et \(E_t^{(C)}\)). 1. \textbf{Gel
temporel à \(T_1 = 2015\)} : construire \(\mathcal{G}_{2015}\)
uniquement sur les publications \(\le 2015\) (restreint aux concepts
biomédicaux actifs ayant un voisinage relationnel non trivial). 2.
\textbf{Période de test tenue à l'écart (2016--2020)} : recenser toutes
les paires \((i,j)\) jamais co-apparues jusqu'en 2015 mais mariées dans
au moins un article publié en 2016--2017. Pour neutraliser le biais du
survivant, on définit deux cibles : - \emph{Cible A (naissance)} :
paires nouvellement publiées en 2016--2017 contre paires non publiées
appariées sur \(\deg(i)\deg(j)\). - \emph{Cible B (fécondité
conditionnelle --- cible principale)} : parmi les \textbf{seules}
nouvelles paires publiées en 2016--2017, opposer celles qui sont
reprises par plusieurs équipes indépendantes et atteignent le décile
supérieur de citations normalisées à cinq ans (2018--2022) à celles qui
restent stériles (aucune reprise indépendante). 3. \textbf{Calcul des
prédicteurs sur \(\mathcal{G}_{2015}\)} : \(d_{\mathrm{surf}}\) (cosinus
SVD/Word2Vec \(\le 2015\) + \(d_{\mathrm{JS}}\) sur branches
MeSH/revues) et \(K(i,j;2015)\) sur les voisinages relationnels
anonymisés de rayon 2 via \textbf{Fused Gromov--Wasserstein}
(\textbf{POT}) + pénalité de conflits d'antonymes. 4.
\textbf{Comparaison préenregistrée} en validation croisée groupée par
concept (empêchant toute fuite par nœud) : - \(M_0\) :
\(\deg(i)\deg(j)\) + nombre de chemins ABC de Swanson +
\(d_{\mathrm{surf}}\) linéaire. - \(M_1\) (H1/H4 sans H3) : \(M_0\) +
spline quadratique de \(d_{\mathrm{surf}}\). - \(M_2\) (cœur H1+H3) :
\(M_1\) + \(K(i,j)\) + interaction \(d_{\mathrm{surf}} \times K(i,j)\).

\textbf{Prédictions et seuils décisionnels (fixés \emph{a priori}).} -
\textbf{P1 (test de H3)} : sur la Cible B, \(M_2\) bat \(M_1\) avec un
gain \(\Delta\text{PR-AUC}(M_2 - M_1) \ge +0{,}03\) (IC bootstrap 95 \%
excluant \(0\), \(p < 0{,}01\) par permutation en blocs sur les
concepts) et \(\hat{\theta} > 0\). - \textbf{P2 (interaction H1
\(\times\) H3)} : le terme d'interaction \(d_{\mathrm{surf}} \times K\)
est strictement positif (\(p < 0{,}01\)) : dans le tercile supérieur de
\(d_{\mathrm{surf}}\), une hausse d'un écart-type de \(K\) accroît les
cotes de fécondité d'un facteur \(\text{OR} \ge 1{,}20\), alors que la
distance de surface seule à \(K\) faible n'augmente pas la fécondité.

\textbf{Lecture des résultats.} - \emph{P1 et P2 atteintes} : le cœur
H1+H3 passe son premier test de réfutation ; on réplique aussitôt sur
les arbres d'équations LaTeX d'arXiv (couche structurelle non bruitée
par le NLP) avant d'ajouter le champ \(h(t)\) et l'évaluation humaine en
aveugle. - \emph{P1 échoue (\(M_2 \approx M_1\) ou
\(\hat{\theta} \le 0\))} : avant d'abandonner définitivement H3, auditer
manuellement un petit échantillon de voisinages SemMedDB ; si la
précision des triplets est satisfaisante mais \(K\) reste inopérant,
\textbf{H3 est réfutée} sous sa forme de morphisme de graphe local. -
\emph{P1 réussit mais P2 échoue (effet négatif de \(d_{\mathrm{surf}}\)
même à \(K\) élevé)} : \textbf{H1 est réfutée} ; l'alignement structurel
favorise la fécondité \emph{locale}, non les mariages éloignés. -
\emph{\(K\) prédit la Cible A (publication) mais pas la Cible B
(fécondité)} : le critère serrure--clé capture la \textbf{plausibilité
rhétorique} perçue par les auteurs et relecteurs, pas la fécondité
scientifique réelle.
%=====================================================================
\section*{Bibliographie vérifiée}\addcontentsline{toc}{section}{Bibliographie vérifiée}
Chaque référence a été cherchée en ligne le 23 novembre 2026. \emph{Mode} : source de la vérification (Crossref par DOI ou recherche bibliographique ; interface d'arXiv ; OpenLibrary pour les livres sans DOI ; page de l'éditeur ; Wikisource). « Résumé lu » : le résumé a été lu en ligne ; sinon seule l'existence et les métadonnées sont vérifiées. \emph{Cité par} : A = Argon, C = Carbon, N = Neon, n = la présente note. Les métadonnées (volume, pages) sont celles des réponses des services, non des PDF.
{\footnotesize\renewcommand{\arraystretch}{1.1}
\begin{longtable}{@{}>{\raggedright\arraybackslash}p{8.7cm}>{\raggedright\arraybackslash}p{3.6cm}>{\raggedright\arraybackslash}p{2.1cm}>{\raggedright\arraybackslash}p{0.9cm}@{}}
\toprule \textbf{Référence} & \textbf{Identifiant} & \textbf{Mode} & \textbf{Cité par}\\ \midrule\endhead \bottomrule\endfoot
Baez, J. C. et Fong, B. (2013). A Noether theorem for Markov processes. & \nolinkurl{arXiv:1203.2035} & arXiv & n\\
Beltagy, I., Lo, K. et Cohan, A. (2019). SciBERT: A pretrained language model for scientific text. \emph{EMNLP-IJCNLP}. & \nolinkurl{doi:10.18653/v1/D19-1371} & Crossref & A C\\
Berlyne, D. E. (1960). \emph{Conflict, Arousal, and Curiosity}. McGraw-Hill. & OpenLibrary OL7597818W & OpenLibrary & A C n\\
Berlyne, D. E. (1970). Novelty, complexity, and hedonic value. \emph{Perception \& Psychophysics} 8(5), 279--286. & \nolinkurl{doi:10.3758/BF03212593} & Crossref & N n\\
Bettencourt, L. M. A., Kaiser, D. I., Kaur, J., Castillo-Chávez, C. et Wojick, D. E. (2008). Population modeling of the emergence and development of scientific fields. \emph{Scientometrics} 75(3), 495--518. & \nolinkurl{doi:10.1007/s11192-007-1888-4} & Crossref & C N\\
Bettencourt, L. M. A., Kaiser, D. I. et Kaur, J. (2009). Scientific discovery and topological transitions in collaboration networks. \emph{Journal of Informetrics} 3, 210--221. & \nolinkurl{doi:10.1016/j.joi.2009.03.001} & Crossref & n\\
Bikard, M., Murray, F. et Gans, J. S. (2015). Exploring trade-offs in the organization of scientific work: Collaboration and scientific reward. \emph{Management Science} 61(7), 1473--1495. & \nolinkurl{doi:10.1287/mnsc.2014.2052} & Crossref & C\\
Boden, M. A. (2004). \emph{The Creative Mind: Myths and Mechanisms} (2\textsuperscript{e} éd. ; 1\textsuperscript{re} éd. 1990). Routledge. & OpenLibrary OL4484307W & OpenLibrary (1\textsuperscript{re} éd.) & A\\
Borgan, Ø., Goldstein, L. et Langholz, B. (1995). Methods for the analysis of sampled cohort data in the Cox proportional hazards model. \emph{The Annals of Statistics} 23(5), 1749--1778. & \nolinkurl{doi:10.1214/aos/1176324322} & Crossref & A\\
Bromham, L., Dinnage, R. et Hua, X. (2016). Interdisciplinary research has consistently lower funding success. \emph{Nature} 534(7609), 684--687. & \nolinkurl{doi:10.1038/nature18315} & Crossref & A\\
Burt, R. S. (2004). Structural holes and good ideas. \emph{American Journal of Sociology} 110(2), 349--399. & \nolinkurl{doi:10.1086/421787} & Crossref & A C N n\\
Campbell, D. T. (1960). Blind variation and selective retentions in creative thought as in other knowledge processes. \emph{Psychological Review} 67(6), 380--400. & \nolinkurl{doi:10.1037/h0040373} & Crossref & C\\
Carey, S. (2009). \emph{The Origin of Concepts}. Oxford University Press. & \nolinkurl{doi:10.1093/acprof:oso/9780195367638.001.0001} & Crossref & C\\
Chalmers, D. J., French, R. M. et Hofstadter, D. R. (1992). High-level perception, representation, and analogy: A critique of artificial intelligence methodology. \emph{Journal of Experimental \& Theoretical Artificial Intelligence} 4, 185--211. & \nolinkurl{doi:10.1080/09528139208953747} & Crossref & n\\
Cohan, A., Feldman, S., Beltagy, I., Downey, D. et Weld, D. S. (2020). SPECTER: Document-level representation learning using citation-informed transformers. \emph{ACL}. & \nolinkurl{doi:10.18653/v1/2020.acl-main.207} & Crossref & C\\
Davis, M. S. (1971). That's interesting! Towards a phenomenology of sociology and a sociology of phenomenology. \emph{Philosophy of the Social Sciences} 1(2), 309--344. & \nolinkurl{doi:10.1177/004839317100100211} & Crossref & A C N n\\
DiGiacomo, R. A., Kremer, J. M. et Shah, D. M. (1989). Fish-oil dietary supplementation in patients with Raynaud's phenomenon: A double-blind, controlled, prospective study. \emph{The American Journal of Medicine} 86, 158--164. & \nolinkurl{doi:10.1016/0002-9343(89)90261-1} & Crossref & n\\
Ehrig, H., Ehrig, K., Prange, U. et Taentzer, G. (2006). \emph{Fundamentals of Algebraic Graph Transformation}. Springer. & \nolinkurl{doi:10.1007/3-540-31188-2} & Crossref & n\\
Falkenhainer, B., Forbus, K. D. et Gentner, D. (1989). The structure-mapping engine: Algorithm and examples. \emph{Artificial Intelligence} 41(1), 1--63. & \nolinkurl{doi:10.1016/0004-3702(89)90077-5} & Crossref & A N n\\
Fauconnier, G. et Turner, M. (1998). Conceptual integration networks. \emph{Cognitive Science} 22(2), 133--187. & \nolinkurl{doi:10.1207/s15516709cog2202_1} & Crossref & n\\
Fauconnier, G. et Turner, M. (2002). \emph{The Way We Think}. Basic Books. & OpenLibrary OL16979359W & OpenLibrary & n\\
Fischhoff, B. (1975). Hindsight is not equal to foresight: The effect of outcome knowledge on judgment under uncertainty. \emph{Journal of Experimental Psychology: Human Perception and Performance} 1(3), 288--299. & \nolinkurl{doi:10.1037/0096-1523.1.3.288} & Crossref & N\\
Foster, J. G., Rzhetsky, A. et Evans, J. A. (2015). Tradition and innovation in scientists' research strategies. \emph{American Sociological Review} 80(5), 875--908. & \nolinkurl{doi:10.1177/0003122415601618} & Crossref & A C N\\
Funk, R. J. et Owen-Smith, J. (2017). A dynamic network measure of technological change. \emph{Management Science} 63(3), 791--817. & \nolinkurl{doi:10.1287/mnsc.2015.2366} & Crossref & A C N\\
Galison, P. (2003). \emph{Einstein's Clocks, Poincaré's Maps: Empires of Time}. W. W. Norton. & OpenLibrary OL8063285W & OpenLibrary & A C N n\\
Gentner, D. (1983). Structure-mapping: A theoretical framework for analogy. \emph{Cognitive Science} 7(2), 155--170. & \nolinkurl{doi:10.1207/s15516709cog0702_3} & Crossref & A C N n\\
Gentner, D. et Markman, A. B. (1997). Structure mapping in analogy and similarity. \emph{American Psychologist} 52(1), 45--56. & \nolinkurl{doi:10.1037/0003-066X.52.1.45} & Crossref & A C\\
Gentner, D. et Stevens, A. L. (dir.) (1983). \emph{Mental Models}. Erlbaum. (Contient Wiser et Carey ; chapitre non vérifié.) & OpenLibrary OL16641020W & OpenLibrary (livre) & N\\
Grünwald, P. D. (2007). \emph{The Minimum Description Length Principle}. MIT Press. & \nolinkurl{doi:10.7551/mitpress/4643.001.0001} & Crossref & N\\
Hamilton, W. L., Leskovec, J. et Jurafsky, D. (2016). Diachronic word embeddings reveal statistical laws of semantic change. \emph{ACL}. & \nolinkurl{doi:10.18653/v1/P16-1141} & Crossref & A C N\\
Hell, P. et Nešetřil, J. (1990). On the complexity of H-coloring. \emph{Journal of Combinatorial Theory, Series B} 48, 92--110. & \nolinkurl{doi:10.1016/0095-8956(90)90132-J} & Crossref & n\\
Hell, P. et Nešetřil, J. (2004). \emph{Graphs and Homomorphisms}. Oxford University Press. & \nolinkurl{doi:10.1093/acprof:oso/9780198528173.001.0001} & Crossref & n\\
Hesse, M. B. (1966). \emph{Models and Analogies in Science}. University of Notre Dame Press (1\textsuperscript{re} éd. 1963). & OpenLibrary OL1419101W & OpenLibrary & N\\
Hofstadter, D. et Sander, E. (2013). \emph{Surfaces and Essences: Analogy as the Fuel and Fire of Thinking}. Basic Books (éd. française : \emph{L'Analogie, cœur de la pensée}, Odile Jacob). & aucun & \textbf{à vérifier} & A C\\
Hofstadter, D. R. et Fluid Analogies Research Group (1995). \emph{Fluid Concepts and Creative Analogies}. Basic Books. & OpenLibrary OL716837W & OpenLibrary & n\\
Holyoak, K. J. et Thagard, P. (1995). \emph{Mental Leaps: Analogy in Creative Thought}. MIT Press. & \nolinkurl{doi:10.7551/mitpress/4549.001.0001} & Crossref & A C\\
Hope, T., Chan, J., Kittur, A. et Shahaf, D. (2017). Accelerating innovation through analogy mining. \emph{KDD}, 235--243. & \nolinkurl{doi:10.1145/3097983.3098038} & Crossref & C\\
Itti, L. et Baldi, P. (2005). A principled approach to detecting surprising events in video. \emph{CVPR}. & \nolinkurl{doi:10.1109/CVPR.2005.40} & Crossref & n\\
Itti, L. et Baldi, P. (2009). Bayesian surprise attracts human attention. \emph{Vision Research} 49(10), 1295--1306. & \nolinkurl{doi:10.1016/j.visres.2008.09.007} & Crossref & A C N n\\
Jüttner, A. et Madarasi, P. (2018). VF2++ --- An improved subgraph isomorphism algorithm. \emph{Discrete Applied Mathematics} 242, 69--81. & \nolinkurl{doi:10.1016/j.dam.2018.02.018} & Crossref & A\\
Kaplan, S. et Vakili, K. (2015). The double-edged sword of recombination in breakthrough innovation. \emph{Strategic Management Journal} 36(10), 1435--1457. & \nolinkurl{doi:10.1002/smj.2294} & Crossref & N\\
Karrer, B., Newman, M. E. J. et Zdeborová, L. (2014). Percolation on sparse networks. \emph{Physical Review Letters} 113(20), 208702. & \nolinkurl{doi:10.1103/PhysRevLett.113.208702} & Crossref & C\\
Kauffman, S. (1995). \emph{At Home in the Universe}. Oxford University Press. & OpenLibrary OL8003172W & OpenLibrary & A\\
Kauffman, S. A. (2000). \emph{Investigations}. Oxford University Press. & OpenLibrary OL3493380W & OpenLibrary & n\\
Ke, Q., Ferrara, E., Radicchi, F. et Flammini, A. (2015). Defining and identifying Sleeping Beauties in science. \emph{PNAS} 112(24), 7426--7431. & \nolinkurl{doi:10.1073/pnas.1424329112} & Crossref & N n\\
Kilicoglu, H., Shin, D., Fiszman, M., Rosemblat, G. et Rindflesch, T. C. (2012). SemMedDB: A PubMed-scale repository of biomedical semantic predications. \emph{Bioinformatics} 28(23), 3158--3160. & \nolinkurl{doi:10.1093/bioinformatics/bts591} & Crossref & C N\\
Kitcher, P. (1981). Explanatory unification. \emph{Philosophy of Science} 48(4), 507--531. & \nolinkurl{doi:10.1086/289019} & Crossref & A C N n\\
Kleinberg, J. (2003). Bursty and hierarchical structure in streams. \emph{Data Mining and Knowledge Discovery} 7(4), 373--397. & \nolinkurl{doi:10.1023/A:1024940629314} & Crossref & A N\\
Koestler, A. (1964). \emph{The Act of Creation}. Hutchinson. & OpenLibrary OL15860364W & OpenLibrary & A C n\\
Krenn, M. et Zeilinger, A. (2020). Predicting research trends with semantic and neural networks with an application in quantum physics. \emph{PNAS} 117(4), 1910--1916. & \nolinkurl{doi:10.1073/pnas.1914370116} & Crossref & n\\
Krenn, M., Buffoni, L., Coutinho, B., Eppel, S., Foster, J. G., Gritsevskiy, A. et al. (2023). Forecasting the future of artificial intelligence with machine learning-based link prediction in an exponentially growing knowledge network. \emph{Nature Machine Intelligence} 5(11), 1326--1335. & \nolinkurl{doi:10.1038/s42256-023-00735-0} ; \nolinkurl{arXiv:2210.00881} & Crossref, arXiv & n\\
Kuhn, T. S. (1962). \emph{The Structure of Scientific Revolutions}. University of Chicago Press. & OpenLibrary OL3259254W & OpenLibrary & A C N\\
Lakatos, I. (1970). Falsification and the methodology of scientific research programmes. In \emph{Criticism and the Growth of Knowledge}, 91--196. Cambridge University Press. & \nolinkurl{doi:10.1017/CBO9781139171434.009} & Crossref & C n\\
Lakatos, I. (1976). \emph{Proofs and Refutations}. Cambridge University Press. & \nolinkurl{doi:10.1017/CBO9781139171472} & Crossref & A\\
Li, J., Yin, Y., Fortunato, S. et Wang, D. (2019). A dataset of publication records for Nobel laureates. \emph{Scientific Data} 6, 33. & \nolinkurl{doi:10.1038/s41597-019-0033-6} & Crossref & A C\\
Lin, Z., Yin, Y., Liu, L. et Wang, D. (2023). SciSciNet: A large-scale open data lake for the science of science research. \emph{Scientific Data} 10, 315. & \nolinkurl{doi:10.1038/s41597-023-02198-9} & Crossref & C\\
Lo, K., Wang, L. L., Neumann, M., Kinney, R. et Weld, D. S. (2020). S2ORC: The Semantic Scholar Open Research Corpus. \emph{ACL}. & \nolinkurl{doi:10.18653/v1/2020.acl-main.447} & Crossref & A C\\
Loewenstein, G. (1994). The psychology of curiosity: A review and reinterpretation. \emph{Psychological Bulletin} 116(1), 75--98. & \nolinkurl{doi:10.1037/0033-2909.116.1.75} & Crossref & N\\
Lu, H., Wu, Y. N. et Holyoak, K. J. (2019). Emergence of analogy from relation learning. \emph{PNAS} 116(10), 4176--4181. & \nolinkurl{doi:10.1073/pnas.1814779116} & Crossref & C n\\
Luan, Y., He, L., Ostendorf, M. et Hajishirzi, H. (2018). Multi-task identification of entities, relations, and coreference for scientific knowledge graph construction. \emph{EMNLP}. & \nolinkurl{doi:10.18653/v1/D18-1360} & Crossref & C\\
Mac Lane, S. (1971). \emph{Categories for the Working Mathematician}. Springer. & \nolinkurl{doi:10.1007/978-1-4612-9839-7} & Crossref & C\\
Mansouri, B., Novotný, V., Agarwal, A., Oard, D. W. et Zanibbi, R. (2022). Overview of ARQMath-3 (2022): Third CLEF lab on answer retrieval for questions on math. \emph{CLEF}, LNCS. & \nolinkurl{doi:10.1007/978-3-031-13643-6_20} & Crossref & C\\
Mansouri, B., Rohatgi, S., Oard, D. W., Wu, J., Giles, C. L. et Zanibbi, R. (2019). Tangent-CFT: An embedding model for mathematical formulas. \emph{ICTIR}. (Argon l'attribue à « Davila et al. ».) & \nolinkurl{doi:10.1145/3341981.3344235} & Crossref & A\\
Marsden, J. E. et West, M. (2001). Discrete mechanics and variational integrators. \emph{Acta Numerica} 10, 357--514. & \nolinkurl{doi:10.1017/S096249290100006X} & Crossref & n\\
Mémoli, F. (2011). Gromov--Wasserstein distances and the metric approach to object matching. \emph{Foundations of Computational Mathematics} 11(4), 417--487. & \nolinkurl{doi:10.1007/s10208-011-9093-5} & Crossref & n\\
Merton, R. K. (1961). Singletons and multiples in scientific discovery: A chapter in the sociology of science. \emph{Proceedings of the American Philosophical Society} 105(5), 470--486. & aucun DOI trouvé ; JSTOR refusé au robot & \textbf{à vérifier} & A C N n\\
Merton, R. K. (1968). The Matthew effect in science. \emph{Science} 159(3810), 56--63. & \nolinkurl{doi:10.1126/science.159.3810.56} & Crossref & C N\\
Mikolov, T., Sutskever, I., Chen, K., Corrado, G. et Dean, J. (2013). Distributed representations of words and phrases and their compositionality. \emph{NeurIPS} 26. & \nolinkurl{arXiv:1310.4546} & arXiv ; page NeurIPS & A C n\\
Mitchell, M. (1993). \emph{Analogy-Making as Perception}. MIT Press. & OpenLibrary (ISBN 9780262515443) & OpenLibrary & n\\
Mitchell, M. (2021). Abstraction and analogy-making in artificial intelligence. \emph{Annals of the New York Academy of Sciences} 1505(1), 79--101. & \nolinkurl{doi:10.1111/nyas.14619} & Crossref & N n\\
Mitchell, M. et Hofstadter, D. R. (1990). The emergence of understanding in a computer model of concepts and analogy-making. \emph{Physica D}. & \nolinkurl{doi:10.1016/0167-2789(90)90086-5} & Crossref & n\\
Newman, M. E. J. (2010). \emph{Networks: An Introduction}. Oxford University Press. & \nolinkurl{doi:10.1093/acprof:oso/9780199206650.001.0001} & Crossref & N\\
Newman, M. E. J. et Ziff, R. M. (2000). Efficient Monte Carlo algorithm and high-precision results for percolation. \emph{Physical Review Letters} 85(19), 4104--4107. & \nolinkurl{doi:10.1103/PhysRevLett.85.4104} & Crossref & A\\
Noether, E. (1918). Invariante Variationsprobleme. \emph{Nachrichten der Gesellschaft der Wissenschaften zu Göttingen}. Traduction anglaise par M. A. Tavel ; repris dans les \emph{Gesammelte Abhandlungen} (Springer). & \nolinkurl{arXiv:physics/0503066} ; \nolinkurl{doi:10.1007/978-3-642-61761-4_15} & arXiv, Crossref (reprise) & N n\\
Nooteboom, B., Van Haverbeke, W., Duysters, G., Gilsing, V. et van den Oord, A. (2007). Optimal cognitive distance and absorptive capacity. \emph{Research Policy} 36(7), 1016--1034. & \nolinkurl{doi:10.1016/j.respol.2007.04.003} & Crossref & A\\
Ogburn, W. F. et Thomas, D. (1922). Are inventions inevitable? A note on social evolution. \emph{Political Science Quarterly} 37(1), 83--98. & \nolinkurl{doi:10.2307/2142320} & Crossref & A C N n\\
Oudeyer, P.-Y., Kaplan, F. et Hafner, V. V. (2007). Intrinsic motivation systems for autonomous mental development. \emph{IEEE Transactions on Evolutionary Computation} 11(2), 265--286. & \nolinkurl{doi:10.1109/TEVC.2006.890271} & Crossref & N n\\
Pearl, J. (2009). \emph{Causality} (2\textsuperscript{e} éd.). Cambridge University Press. & \nolinkurl{doi:10.1017/CBO9780511803161} & Crossref & N\\
Peyré, G., Cuturi, M. et Solomon, J. (2016). Gromov--Wasserstein averaging of kernel and distance matrices. \emph{ICML}, PMLR 48. & \nolinkurl{proceedings.mlr.press/v48/peyre16.html} & page PMLR & n\\
Poincaré, H. (1908). \emph{Science et méthode}. Flammarion. Livre I, chap. III, « L'invention mathématique » ; édition définitive transcrite sur Wikisource, p.~51 et 54. & fr.wikisource.org, « Science et méthode/Livre premier, § III » & Wikisource (texte lu) & n\\
Priem, J., Piwowar, H. et Orr, R. (2022). OpenAlex: A fully-open index of scholarly works, authors, venues, institutions, and concepts. & \nolinkurl{arXiv:2205.01833} & arXiv & C\\
Rindflesch, T. C. et Fiszman, M. (2003). The interaction of domain knowledge and linguistic structure in natural language processing: Interpreting hypernymic propositions in biomedical text. \emph{Journal of Biomedical Informatics} 36(6), 462--477. & \nolinkurl{doi:10.1016/j.jbi.2003.11.003} & Crossref & A\\
Rissanen, J. (1978). Modeling by shortest data description. \emph{Automatica} 14(5), 465--471. & \nolinkurl{doi:10.1016/0005-1098(78)90005-5} & Crossref & C\\
Rogers, A., Drozd, A. et Li, B. (2017). The (too many) problems of analogical reasoning with word vectors. \emph{*SEM 2017}. & \nolinkurl{doi:10.18653/v1/S17-1017} & Crossref & N n\\
Rzhetsky, A., Foster, J. G., Foster, I. T. et Evans, J. A. (2015). Choosing experiments to accelerate collective discovery. \emph{PNAS} 112(47), 14569--14574. & \nolinkurl{doi:10.1073/pnas.1509757112} & Crossref & n\\
Schlichtkrull, M. et al. (2018). Modeling relational data with graph convolutional networks. \emph{ESWC}, LNCS, 593--607. & \nolinkurl{doi:10.1007/978-3-319-93417-4_38} & Crossref & C\\
Schmidhuber, J. (2009). Driven by compression progress: A simple principle explains essential aspects of subjective beauty, novelty, surprise, interestingness, attention, curiosity, creativity, art, science, music, jokes. & \nolinkurl{arXiv:0812.4360} & arXiv & n\\
Schmidhuber, J. (2010). Formal theory of creativity, fun, and intrinsic motivation (1990--2010). \emph{IEEE Transactions on Autonomous Mental Development} 2(3), 230--247. & \nolinkurl{doi:10.1109/TAMD.2010.2056368} & Crossref & A C N n\\
Schultz, W., Dayan, P. et Montague, P. R. (1997). A neural substrate of prediction and reward. \emph{Science} 275(5306), 1593--1599. & \nolinkurl{doi:10.1126/science.275.5306.1593} & Crossref & A C N n\\
Shi, F. et Evans, J. (2023). Surprising combinations of research contents and contexts are related to impact and emerge with scientific outsiders from distant disciplines. \emph{Nature Communications} 14, 1641. & \nolinkurl{doi:10.1038/s41467-023-36741-4} & Crossref (résumé lu) & N n\\
Simonton, D. K. (1979). Multiple discovery and invention: Zeitgeist, genius, or chance? \emph{Journal of Personality and Social Psychology} 37(9), 1603--1616. & \nolinkurl{doi:10.1037/0022-3514.37.9.1603} & Crossref & N n\\
Simonton, D. K. (2004). \emph{Creativity in Science: Chance, Logic, Genius, and Zeitgeist}. Cambridge University Press. & \nolinkurl{doi:10.1017/CBO9781139165358} & Crossref & A C n\\
Sourati, J. et Evans, J. A. (2023). Accelerating science with human-aware artificial intelligence. \emph{Nature Human Behaviour} 7(10), 1682--1696. & \nolinkurl{doi:10.1038/s41562-023-01648-z} & Crossref ; résumé lu (Europe PMC) & n\\
Spivak, D. I. (2014). \emph{Category Theory for the Sciences}. MIT Press. & OpenLibrary OL20783371W & OpenLibrary & C\\
Spivak, D. I. et Kent, R. E. (2012). Ologs: A categorical framework for knowledge representation. \emph{PLoS ONE} 7(1), e24274. & \nolinkurl{doi:10.1371/journal.pone.0024274} & Crossref & A n\\
Stephan, P. (2012). \emph{How Economics Shapes Science}. Harvard University Press. & \nolinkurl{doi:10.4159/harvard.9780674062757} & Crossref & C\\
Swanson, D. R. (1986a). Fish oil, Raynaud's syndrome, and undiscovered public knowledge. \emph{Perspectives in Biology and Medicine} 30(1), 7--18. & \nolinkurl{doi:10.1353/pbm.1986.0087} & Crossref & C N n\\
Swanson, D. R. (1986b). Undiscovered public knowledge. \emph{The Library Quarterly} 56(2), 103--118. & \nolinkurl{doi:10.1086/601720} & Crossref & A n\\
Swanson, D. R. et Smalheiser, N. R. (1997). An interactive system for finding complementary literatures: A stimulus to scientific discovery. \emph{Artificial Intelligence} 91(2), 183--203. & \nolinkurl{doi:10.1016/S0004-3702(97)00008-8} & Crossref & N n\\
Thagard, P. (1989). Explanatory coherence. \emph{Behavioral and Brain Sciences} 12(3), 435--467. & \nolinkurl{doi:10.1017/S0140525X00057046} & Crossref & N\\
Traag, V. A., Waltman, L. et van Eck, N. J. (2019). From Louvain to Leiden: Guaranteeing well-connected communities. \emph{Scientific Reports} 9, 5233. & \nolinkurl{doi:10.1038/s41598-019-41695-z} & Crossref & A C N\\
Tria, F., Loreto, V., Servedio, V. D. P. et Strogatz, S. H. (2014). The dynamics of correlated novelties. \emph{Scientific Reports} 4, 5890. & \nolinkurl{doi:10.1038/srep05890} & Crossref (résumé lu) & A n\\
Tshitoyan, V., Dagdelen, J., Weston, L., Dunn, A., Rong, Z., Kononova, O., Persson, K. A., Ceder, G. et Jain, A. (2019). Unsupervised word embeddings capture latent knowledge from materials science literature. \emph{Nature} 571, 95--98. & \nolinkurl{doi:10.1038/s41586-019-1335-8} & Crossref ; résumé lu (Europe PMC) & A C N n\\
Unruh, W. G. (1981). Experimental black-hole evaporation? \emph{Physical Review Letters} 46(21), 1351--1353. & \nolinkurl{doi:10.1103/PhysRevLett.46.1351} & Crossref & C n\\
Uzzi, B., Mukherjee, S., Stringer, M. et Jones, B. (2013). Atypical combinations and scientific impact. \emph{Science} 342(6157), 468--472. & \nolinkurl{doi:10.1126/science.1240474} & Crossref (résumé lu) & A C N n\\
Vayer, T., Chapel, L., Flamary, R., Tavenard, R. et Courty, N. (2019). Optimal transport for structured data with application on graphs. \emph{ICML}, PMLR 97. (Neon écrit « Titouan, V. ».) & \nolinkurl{arXiv:1805.09114} & arXiv (résumé lu) & A N n\\
Vayer, T., Chapel, L., Flamary, R., Tavenard, R. et Courty, N. (2020). Fused Gromov-Wasserstein distance for structured objects. \emph{Algorithms} 13(9), 212. & \nolinkurl{doi:10.3390/a13090212} & Crossref & C\\
Wadden, D., Wennberg, U., Luan, Y. et Hajishirzi, H. (2019). Entity, relation, and event extraction with contextualized span representations. \emph{EMNLP-IJCNLP}. & \nolinkurl{doi:10.18653/v1/D19-1585} & Crossref & A\\
Wang, J., Veugelers, R. et Stephan, P. (2017). Bias against novelty in science: A cautionary tale for users of bibliometric indicators. \emph{Research Policy} 46(8), 1416--1436. & \nolinkurl{doi:10.1016/j.respol.2017.06.006} & Crossref & C N\\
Wiser, M. et Carey, S. (1983). When heat and temperature were one. In Gentner et Stevens (dir.), \emph{Mental Models}, 267--297. & livre vérifié (OpenLibrary OL16641020W) ; chapitre non vérifié & \textbf{à vérifier} (chapitre) & N n\\
Wu, L., Wang, D. et Evans, J. A. (2019). Large teams develop and small teams disrupt science and technology. \emph{Nature} 566(7744), 378--382. & \nolinkurl{doi:10.1038/s41586-019-0941-9} & Crossref & A C N\\
\end{longtable}}
\end{document}
